% =====================================================================
%  Chapitre 7 : portes multi-qubits, circuits quantiques et téléportation
%  Fichier importé par main.tex avec % =====================================================================
%  Chapitre 7 : portes multi-qubits, circuits quantiques et téléportation
%  Fichier importé par main.tex avec % =====================================================================
%  Chapitre 7 : portes multi-qubits, circuits quantiques et téléportation
%  Fichier importé par main.tex avec % =====================================================================
%  Chapitre 7 : portes multi-qubits, circuits quantiques et téléportation
%  Fichier importé par main.tex avec \include{chapitre7}.
%  Comme chapitre4.tex, chapitre5.tex et chapitre6.tex, il ouvre une NOUVELLE
%  SECTION (numéro 5). Pour le rattacher à la section précédente, remplacer
%  \section par \subsection et chaque \subsection par \subsubsection.
%  Pas de préambule : environnements (definition, resultat, remarque),
%  macros (\ii, \ee, \dd, \ket, \bra, \braket) et couleurs (insarouge,
%  insagris) sont définis dans main.tex. Les symboles de circuits (contrôle,
%  cible, croix, mesure) sont définis juste après avec \tikzset.
% =====================================================================

% Macros de Dirac (déjà définies dans main.tex récent ; ce bloc évite l'erreur
% « Undefined control sequence \ket » si main.tex est une ancienne version).
\providecommand{\ket}[1]{|#1\rangle}
\providecommand{\bra}[1]{\langle #1|}
\providecommand{\braket}[2]{\langle #1|#2\rangle}

% Symboles de circuits pour les figures de ce chapitre.
\tikzset{
  ctl/.pic   = {\fill[insarouge] (0,0) circle (2.6pt);},
  cible/.pic = {\draw[insarouge, thick, fill=white] (0,0) circle (0.28);
                \draw[insarouge, thick] (-0.28,0) -- (0.28,0);
                \draw[insarouge, thick] (0,-0.28) -- (0,0.28);},
  croix/.pic = {\draw[insarouge, thick] (-0.17,-0.17) -- (0.17,0.17);
                \draw[insarouge, thick] (-0.17,0.17) -- (0.17,-0.17);},
  mesure/.pic = {\draw[thick, fill=white] (-0.3,-0.27) rectangle (0.3,0.27);
                 \draw[thick] (-0.19,-0.13) arc[start angle=155, end angle=25, radius=0.21];
                 \draw[thick, -{Stealth[length=1.3mm]}] (0,-0.17) -- (0.14,0.12);},
  porte/.style = {draw=insarouge, fill=insarouge!8, minimum size=7mm, thick},
}

% =====================================================================
\section{Portes multi-qubits, circuits et téléportation}\label{sec:circuits}
% =====================================================================

La section précédente a décrit les états de deux qubits. On étudie ici les portes qui
agissent sur plusieurs qubits : le CNOT et le SWAP, puis les portes contrôlées à deux et à
trois qubits (contrôlé-$U$, Fredkin, Toffoli). On assemble ensuite des portes dans un
circuit qui fabrique un état de Bell, on mesure l'information contenue dans un tel état,
on distingue les états purs des états mélangés, et on démontre le protocole de
téléportation. Tous les calculs sont détaillés.

\begin{remarque}[Convention pour les circuits]
Dans un circuit, le temps va de gauche à droite et chaque qubit a un fil. Les fils sont
dessinés dans l'ordre \emph{inverse} de l'écriture des kets : le fil du \emph{haut} porte le
\emph{dernier} chiffre du ket, le fil du \emph{bas} le premier. Par exemple, si le qubit $A$
est dessiné au-dessus du qubit $B$, le ket s'écrit $\ket{b,a}=\ket{b}_B\otimes\ket{a}_A$.
La figure~\ref{fig:bell} de la section précédente suit la même convention.
\end{remarque}

% ---------------------------------------------------------------------
\subsection{Porte CNOT}\label{sec:cnot}
% ---------------------------------------------------------------------

\begin{definition}{Porte CNOT (NOT contrôlé)}
La porte CNOT agit sur deux qubits : le qubit $A$ est le \emph{contrôle}, le qubit $B$ est la
\emph{cible}. La cible est inversée (porte $X$) si le contrôle vaut $1$ et reste inchangée
s'il vaut $0$. Avec le système écrit $(B,A)$, soit $\ket{b,a}=\ket{b}_B\otimes\ket{a}_A$ :
\[
  \mathrm{CNOT}\,\ket{b,a}=\ket{b\oplus a,\,a},
  \qquad
  \mathrm{CNOT}=I\otimes\ket{0}\bra{0}+X\otimes\ket{1}\bra{1},
\]
où $\oplus$ est l'addition modulo $2$ ($0\oplus0=1\oplus1=0$ et $0\oplus1=1\oplus0=1$).
\end{definition}

Sur le circuit de la figure~\ref{fig:cnot}, le point noir posé sur le fil de $A$ marque le
contrôle et le symbole $\oplus$ posé sur le fil de $B$ marque la cible.

\begin{figure}[!ht]
\centering
\begin{tikzpicture}[>=Stealth, font=\small]
  \draw (0,1.2) -- (6,1.2);
  \draw (0,0) -- (6,0);
  \node[anchor=east] at (0,1.2) {$\ket{a}_A$};
  \node[anchor=east] at (0,0) {$\ket{b}_B$};
  \node[anchor=west] at (6,1.2) {$\ket{a}_A$};
  \node[anchor=west] at (6,0) {$\ket{b\oplus a}_B$};
  \draw[insarouge, thick] (3,1.2) -- (3,0);
  \pic at (3,1.2) {ctl};
  \pic at (3,0) {cible};
  \node[anchor=south, text=insagris] at (3,1.5) {qubit de contrôle};
  \node[anchor=north, text=insagris] at (3,-0.5) {qubit cible};
\end{tikzpicture}
\caption{Porte CNOT : le fil du haut ($A$) est le contrôle, le fil du bas ($B$) est la
cible. Le ket s'écrit $\ket{b,a}$ (le fil du haut donne le dernier chiffre).}
\label{fig:cnot}
\end{figure}

\subsubsection{Table de vérité}

\begin{center}
\renewcommand{\arraystretch}{1.3}
\begin{tabular}{@{}cccc@{}}
\toprule
Entrée $\ket{b,a}$ & Contrôle $a$ & Sortie $\ket{b\oplus a,a}$ & Effet sur la cible $B$ \\
\midrule
$\ket{00}$ & $0$ & $\ket{00}$ & aucun \\
$\ket{01}$ & $1$ & $\ket{11}$ & $0\to1$ \\
$\ket{10}$ & $0$ & $\ket{10}$ & aucun \\
$\ket{11}$ & $1$ & $\ket{01}$ & $1\to0$ \\
\bottomrule
\end{tabular}
\end{center}

\subsubsection{Matrice du CNOT}

La $j$-ième colonne d'une matrice $4\times4$ est l'image du $j$-ième vecteur de base
($\ket{00},\ket{01},\ket{10},\ket{11}$, dans cet ordre). La table donne :
\begin{align*}
  \mathrm{CNOT}\ket{00}&=\ket{00}=\begin{pmatrix}1\\0\\0\\0\end{pmatrix},&
  \mathrm{CNOT}\ket{01}&=\ket{11}=\begin{pmatrix}0\\0\\0\\1\end{pmatrix},\\[1.5ex]
  \mathrm{CNOT}\ket{10}&=\ket{10}=\begin{pmatrix}0\\0\\1\\0\end{pmatrix},&
  \mathrm{CNOT}\ket{11}&=\ket{01}=\begin{pmatrix}0\\1\\0\\0\end{pmatrix},
\end{align*}
d'où, en juxtaposant ces quatre colonnes,
\[
  \mathrm{CNOT}=\begin{pmatrix}
    1 & 0 & 0 & 0\\ 0 & 0 & 0 & 1\\ 0 & 0 & 1 & 0\\ 0 & 1 & 0 & 0
  \end{pmatrix}.
\]

On retrouve cette matrice avec la formule $I\otimes\ket{0}\bra{0}+X\otimes\ket{1}\bra{1}$ et
la règle $A\otimes B=(A_{ij}B)$ (section~\ref{sec:op2}), où
$\ket{0}\bra{0}=\begin{pmatrix}1&0\\0&0\end{pmatrix}$ et
$\ket{1}\bra{1}=\begin{pmatrix}0&0\\0&1\end{pmatrix}$ :
\[
  I\otimes\ket{0}\bra{0}=\begin{pmatrix}1&0&0&0\\0&0&0&0\\0&0&1&0\\0&0&0&0\end{pmatrix},
  \qquad
  X\otimes\ket{1}\bra{1}=\begin{pmatrix}0&0&0&0\\0&0&0&1\\0&0&0&0\\0&1&0&0\end{pmatrix},
\]
et la somme de ces deux matrices est la matrice du CNOT.

\subsubsection{Action sur un état quelconque}

Pour $\ket{\Psi}=c_{00}\ket{00}+c_{01}\ket{01}+c_{10}\ket{10}+c_{11}\ket{11}$ :
\[
  \mathrm{CNOT}\begin{pmatrix}c_{00}\\c_{01}\\c_{10}\\c_{11}\end{pmatrix}
  =\begin{pmatrix}
    1\cdot c_{00}+0\cdot c_{01}+0\cdot c_{10}+0\cdot c_{11}\\
    0\cdot c_{00}+0\cdot c_{01}+0\cdot c_{10}+1\cdot c_{11}\\
    0\cdot c_{00}+0\cdot c_{01}+1\cdot c_{10}+0\cdot c_{11}\\
    0\cdot c_{00}+1\cdot c_{01}+0\cdot c_{10}+0\cdot c_{11}
  \end{pmatrix}
  =\begin{pmatrix}c_{00}\\c_{11}\\c_{10}\\c_{01}\end{pmatrix}.
\]
Le CNOT échange les amplitudes de $\ket{01}$ et de $\ket{11}$, c'est-à-dire des deux états
dont le contrôle vaut $1$, et laisse les autres.

\paragraph{Exemple : contrôle en superposition.} Soit $\ket{0}_B\otimes\ket{+}_A
=\frac{1}{\sqrt2}\big(\ket{00}+\ket{01}\big)=\frac{1}{\sqrt2}(1,1,0,0)^{\mathsf T}$,
un état séparable. Alors
\[
  \mathrm{CNOT}\,\frac{1}{\sqrt2}\begin{pmatrix}1\\1\\0\\0\end{pmatrix}
  =\frac{1}{\sqrt2}\begin{pmatrix}1\\0\\0\\1\end{pmatrix}
  =\frac{1}{\sqrt2}\big(\ket{00}+\ket{11}\big)=\ket{\Phi^+}.
\]
Le résultat est intriqué ($D=c_{00}c_{11}-c_{01}c_{10}=\frac12\neq0$ ; section~\ref{sec:separable}) :
un CNOT dont le contrôle est en superposition crée de l'intrication.

\paragraph{Réversibilité.} Deux CNOT successifs redonnent l'état initial :
$\ket{b,a}\to\ket{b\oplus a,a}\to\ket{b\oplus a\oplus a,a}=\ket{b,a}$, donc
$\mathrm{CNOT}^2=\mathrm{Id}$. La matrice est réelle et symétrique, donc
$\mathrm{CNOT}^\dagger\mathrm{CNOT}=\mathrm{CNOT}^2=\mathrm{Id}$ : elle est bien unitaire.

\begin{remarque}[Contrôle en dernier ou en premier]
La matrice dépend de l'ordre d'écriture du ket. Notons $\mathrm{CNOT}_{21}$ le CNOT dont le
contrôle est le \emph{second} chiffre du ket (c'est la matrice ci-dessus) et
$\mathrm{CNOT}_{12}$ celui dont le contrôle est le \emph{premier} chiffre, utilisé à la
section précédente :
\[
  \mathrm{CNOT}_{12}=\ket{0}\bra{0}\otimes I+\ket{1}\bra{1}\otimes X
  =\begin{pmatrix}1&0&0&0\\0&1&0&0\\0&0&0&1\\0&0&1&0\end{pmatrix}.
\]
Un même circuit (contrôle $A$, cible $B$) a donc deux matrices, selon que l'on écrit
$A$ en dernier ($\mathrm{CNOT}_{21}$) ou en premier ($\mathrm{CNOT}_{12}$). Les deux sont
reliées par le SWAP (section suivante).
\end{remarque}

% ---------------------------------------------------------------------
\subsection{Porte SWAP}\label{sec:swap}
% ---------------------------------------------------------------------

\begin{definition}{Porte SWAP (échange)}
La porte SWAP échange les états des deux qubits :
\[
  \mathrm{SWAP}\big(\ket{\psi}\otimes\ket{\phi}\big)=\ket{\phi}\otimes\ket{\psi}.
\]
Sur les états de base, $\mathrm{SWAP}\ket{a,b}=\ket{b,a}$ ; on prolonge par linéarité.
\end{definition}

Son symbole est formé de deux croix reliées par un trait vertical
(figure~\ref{fig:swap}).

\begin{figure}[!ht]
\centering
\begin{tikzpicture}[>=Stealth, font=\small]
  \draw (0,1.2) -- (6,1.2);
  \draw (0,0) -- (6,0);
  \node[anchor=east] at (0,1.2) {$\ket{a}$};
  \node[anchor=east] at (0,0) {$\ket{b}$};
  \node[anchor=west] at (6,1.2) {$\ket{b}$};
  \node[anchor=west] at (6,0) {$\ket{a}$};
  \draw[insarouge, thick] (3,1.2) -- (3,0);
  \pic at (3,1.2) {croix};
  \pic at (3,0) {croix};
\end{tikzpicture}
\caption{Porte SWAP : les deux qubits échangent leurs états.}
\label{fig:swap}
\end{figure}

\subsubsection{Table et matrice}

\begin{center}
\renewcommand{\arraystretch}{1.3}
\begin{tabular}{@{}cc@{}}
\toprule
Entrée & Sortie \\
\midrule
$\ket{00}$ & $\ket{00}$ \\
$\ket{01}$ & $\ket{10}$ \\
$\ket{10}$ & $\ket{01}$ \\
$\ket{11}$ & $\ket{11}$ \\
\bottomrule
\end{tabular}
\end{center}

Les images des quatre vecteurs de base sont les colonnes de la matrice :
\[
  \mathrm{SWAP}\ket{00}=\begin{pmatrix}1\\0\\0\\0\end{pmatrix},\quad
  \mathrm{SWAP}\ket{01}=\begin{pmatrix}0\\0\\1\\0\end{pmatrix},\quad
  \mathrm{SWAP}\ket{10}=\begin{pmatrix}0\\1\\0\\0\end{pmatrix},\quad
  \mathrm{SWAP}\ket{11}=\begin{pmatrix}0\\0\\0\\1\end{pmatrix},
\]
\[
  \mathrm{SWAP}=\begin{pmatrix}
    1 & 0 & 0 & 0\\ 0 & 0 & 1 & 0\\ 0 & 1 & 0 & 0\\ 0 & 0 & 0 & 1
  \end{pmatrix}.
\]
Appliquée à $\ket{01}=(0,1,0,0)^{\mathsf T}$, cette matrice sélectionne sa deuxième colonne :
\[
  \begin{pmatrix}
    1 & 0 & 0 & 0\\ 0 & 0 & 1 & 0\\ 0 & 1 & 0 & 0\\ 0 & 0 & 0 & 1
  \end{pmatrix}
  \begin{pmatrix}0\\1\\0\\0\end{pmatrix}
  =\begin{pmatrix}0\\0\\1\\0\end{pmatrix}=\ket{10}.
\]
Sur un état quelconque, $\mathrm{SWAP}\,(c_{00},c_{01},c_{10},c_{11})^{\mathsf T}
=(c_{00},c_{10},c_{01},c_{11})^{\mathsf T}$ : les amplitudes de $\ket{01}$ et $\ket{10}$
sont échangées.

\subsubsection{Accord avec la définition}

Pour deux états $\ket{\psi}=(a_0,a_1)^{\mathsf T}$ et $\ket{\phi}=(b_0,b_1)^{\mathsf T}$ :
\[
  \mathrm{SWAP}\,(\ket{\psi}\otimes\ket{\phi})
  =\mathrm{SWAP}\begin{pmatrix}a_0b_0\\a_0b_1\\a_1b_0\\a_1b_1\end{pmatrix}
  =\begin{pmatrix}a_0b_0\\a_1b_0\\a_0b_1\\a_1b_1\end{pmatrix}
  =\begin{pmatrix}b_0\\b_1\end{pmatrix}\otimes\begin{pmatrix}a_0\\a_1\end{pmatrix}
  =\ket{\phi}\otimes\ket{\psi},
\]
ce qui est bien la définition. La porte s'écrit aussi avec des projecteurs :
\[
  \mathrm{SWAP}=\sum_{a,b\in\{0,1\}}\ket{a}\bra{b}\otimes\ket{b}\bra{a}.
\]
En effet, pour $\ket{c}\otimes\ket{d}$ on trouve
$\sum_{a,b}\ket{a}\braket{b}{c}\otimes\ket{b}\braket{a}{d}$, et seul le terme $b=c$, $a=d$
est non nul : le résultat est $\ket{d}\otimes\ket{c}$.

\paragraph{Exemples.}
\begin{itemize}[nosep]
  \item Deux états différents : $\mathrm{SWAP}\big(\ket{0}\otimes\ket{+}\big)
        =\mathrm{SWAP}\,\frac{1}{\sqrt2}(1,1,0,0)^{\mathsf T}=\frac{1}{\sqrt2}(1,0,1,0)^{\mathsf T}
        =\ket{+}\otimes\ket{0}$.
  \item L'état $\ket{\psi}=\frac{1}{\sqrt2}\ket{00}+\frac{1}{\sqrt2}\ket{11}$ ne change pas :
        $\mathrm{SWAP}\ket{\psi}=\frac{1}{\sqrt2}\ket{00}+\frac{1}{\sqrt2}\ket{11}=\ket{\psi}$,
        car $\ket{00}$ et $\ket{11}$ sont invariants. De même
        $\mathrm{SWAP}\ket{\Phi^-}=\ket{\Phi^-}$ et $\mathrm{SWAP}\ket{\Psi^+}=\ket{\Psi^+}$, alors que
        $\mathrm{SWAP}\ket{\Psi^-}=\frac{1}{\sqrt2}(\ket{10}-\ket{01})=-\ket{\Psi^-}$.
\end{itemize}

\paragraph{Propriétés.} Échanger deux fois redonne l'état de départ :
$\mathrm{SWAP}^2=\mathrm{Id}$ ; la matrice est réelle et symétrique, donc
$\mathrm{SWAP}^\dagger\mathrm{SWAP}=\mathrm{SWAP}^2=\mathrm{Id}$ : elle est unitaire.

\subsubsection{Lien avec le CNOT}

\begin{resultat}{SWAP échange les rôles de contrôle et de cible}
\[
  \mathrm{CNOT}_{21}=\mathrm{SWAP}\;\mathrm{CNOT}_{12}\;\mathrm{SWAP}.
\]
\end{resultat}

\emph{Sur les états de base.}
$\mathrm{SWAP}\,\mathrm{CNOT}_{12}\,\mathrm{SWAP}\ket{b,a}
=\mathrm{SWAP}\,\mathrm{CNOT}_{12}\ket{a,b}
=\mathrm{SWAP}\ket{a,\,a\oplus b}=\ket{a\oplus b,\,a}=\mathrm{CNOT}_{21}\ket{b,a}$.

\emph{Par les matrices.} Multiplier à gauche par $\mathrm{SWAP}$ échange les lignes $2$ et
$3$, multiplier à droite l'échange les colonnes $2$ et $3$ :
\[
  \mathrm{SWAP}\,\mathrm{CNOT}_{12}
  =\begin{pmatrix}1&0&0&0\\0&0&0&1\\0&1&0&0\\0&0&1&0\end{pmatrix},
  \qquad
  \mathrm{SWAP}\,\mathrm{CNOT}_{12}\,\mathrm{SWAP}
  =\begin{pmatrix}1&0&0&0\\0&0&0&1\\0&0&1&0\\0&1&0&0\end{pmatrix}=\mathrm{CNOT}_{21}.
\]

\begin{remarque}[SWAP avec trois CNOT]
$\mathrm{SWAP}=\mathrm{CNOT}_{12}\,\mathrm{CNOT}_{21}\,\mathrm{CNOT}_{12}$. En effet, sur
$\ket{a,b}$ (on applique les portes de droite à gauche) :
\[
  \ket{a,b}\ \xrightarrow{\ \mathrm{CNOT}_{12}\ }\ \ket{a,\,a\oplus b}
  \ \xrightarrow{\ \mathrm{CNOT}_{21}\ }\ \ket{a\oplus(a\oplus b),\,a\oplus b}=\ket{b,\,a\oplus b}
  \ \xrightarrow{\ \mathrm{CNOT}_{12}\ }\ \ket{b,\,b\oplus(a\oplus b)}=\ket{b,a}.
\]
Les deux membres coïncident sur les états de base, donc sur tous les états par linéarité.
\end{remarque}

% ---------------------------------------------------------------------
\subsection{Un circuit quantique : le générateur d'états de Bell}\label{sec:circuit-bell}
% ---------------------------------------------------------------------

On assemble une porte de Hadamard et un CNOT. Le système est $(A,B)$ avec le ket
$\ket{a,b}=\ket{a}_A\otimes\ket{b}_B$ ; d'après la convention de la section, $B$ est dessiné
en haut et $A$ en bas (figure~\ref{fig:circuit-bell}). Le circuit a deux étapes :
\begin{enumerate}[nosep]
  \item la porte $I\otimes H$ : identité sur $A$, Hadamard sur $B$ ;
  \item un CNOT de contrôle $B$ et de cible $A$ (la matrice $\mathrm{CNOT}_{21}$ de la
        section~\ref{sec:cnot}, puisque le contrôle est le second chiffre).
\end{enumerate}
Les traits pointillés délimitent les états $\ket{\Pi_0}$, $\ket{\Pi_1}$ et $\ket{\Pi_2}$ :
avant les portes, après $I\otimes H$, après le CNOT.

\begin{figure}[!ht]
\centering
\begin{tikzpicture}[>=Stealth, font=\small]
  % fils : B en haut, A en bas
  \draw (0,1.4) -- (7.6,1.4);
  \draw (0,0) -- (7.6,0);
  \node[anchor=east] at (0,1.4) {$\ket{0}_B$};
  \node[anchor=east] at (0,0) {$\ket{0}_A$};
  % première étape : H sur B, identité (pointillée) sur A
  \node[porte] at (2.2,1.4) {$H$};
  \node[draw=insagris, dashed, fill=white, minimum size=7mm] at (2.2,0) {$I$};
  % CNOT : contrôle B (haut), cible A (bas)
  \draw[insarouge, thick] (5,1.4) -- (5,0);
  \pic at (5,1.4) {ctl};
  \pic at (5,0) {cible};
  % accolade de sortie
  \draw[insagris, thick] (7.75,1.6) -- (7.9,1.6) -- (7.9,-0.2) -- (7.75,-0.2);
  \node[anchor=west] at (7.9,0.7) {$\ket{\beta^{00}}_{AB}$};
  % états intermédiaires
  \foreach \x in {0.7,3.6,6.4}
    \draw[dashed, insagris!70] (\x,1.9) -- (\x,-0.7);
  \node[anchor=north, font=\footnotesize] at (0.7,-0.7) {$\ket{\Pi_0}=\ket{00}$};
  \node[anchor=north, font=\footnotesize] at (3.6,-0.7) {$\ket{\Pi_1}=\ket{0}\otimes\ket{+}$};
  \node[anchor=north, font=\footnotesize] at (6.4,-0.7) {$\ket{\Pi_2}=\ket{\beta^{00}}$};
\end{tikzpicture}
\caption{Générateur d'états de Bell : $I\otimes H$, puis un CNOT de contrôle $B$ (fil du haut)
et de cible $A$ (fil du bas). L'entrée représentée est $\ket{00}$.}
\label{fig:circuit-bell}
\end{figure}

\subsubsection{Entrée \texorpdfstring{$\ket{00}$}{|00>} : calcul pas à pas}

\begin{align*}
  \ket{\Pi_0}&=\ket{0}_A\otimes\ket{0}_B=\ket{00},\\[1ex]
  \ket{\Pi_1}&=(I\otimes H)\ket{\Pi_0}=I\ket{0}_A\otimes H\ket{0}_B=\ket{0}\otimes\ket{+}
    =\frac{\ket{00}+\ket{01}}{\sqrt2},\\[1ex]
  \ket{\Pi_2}&=\mathrm{CNOT}\ket{\Pi_1}=\frac{\mathrm{CNOT}\ket{00}+\mathrm{CNOT}\ket{01}}{\sqrt2}
    =\frac{\ket{00}+\ket{11}}{\sqrt2}=\ket{\beta^{00}}=\ket{\Phi^+}.
\end{align*}
Dans $\ket{\Pi_1}$, le premier chiffre est celui de $A$ et le second celui de $B$. Le CNOT a
pour contrôle $B$ : $\ket{00}$ ($B=0$) ne change pas, $\ket{01}$ ($B=1$) devient $\ket{11}$
puisque $A$ est inversé. Le résultat est un \emph{état de Bell}.

Avec les matrices, la règle $A\otimes B=(A_{ij}B)$ donne
\[
  I\otimes H=\begin{pmatrix}1\cdot H&0\\0&1\cdot H\end{pmatrix}
  =\frac{1}{\sqrt2}\begin{pmatrix}1&1&0&0\\1&-1&0&0\\0&0&1&1\\0&0&1&-1\end{pmatrix}.
\]
Sa première colonne est $\ket{\Pi_1}=\frac{1}{\sqrt2}(1,1,0,0)^{\mathsf T}$ (image de $\ket{00}$).
Le CNOT échange les amplitudes de $\ket{01}$ et de $\ket{11}$ (section~\ref{sec:cnot}), donc
\[
  \ket{\Pi_2}=\mathrm{CNOT}\,\frac{1}{\sqrt2}\begin{pmatrix}1\\1\\0\\0\end{pmatrix}
  =\frac{1}{\sqrt2}\begin{pmatrix}1\\0\\0\\1\end{pmatrix}.
\]

\subsubsection{Les quatre entrées}

Prenons $B$ dans l'état $\ket{x}$ et $A$ dans l'état $\ket{y}$, avec $x,y\in\{0,1\}$. Alors
$\ket{\Pi_0}=\ket{y}_A\otimes\ket{x}_B=\ket{y,x}$. Comme
$H\ket{x}=\frac{1}{\sqrt2}\big(\ket{0}+(-1)^x\ket{1}\big)$ (c'est $\ket{+}$ pour $x=0$ et
$\ket{-}$ pour $x=1$),
\[
  \ket{\Pi_1}=\ket{y}\otimes H\ket{x}=\frac{1}{\sqrt2}\Big(\ket{y,0}+(-1)^x\ket{y,1}\Big).
\]
Le CNOT (contrôle $B$) ne change pas $\ket{y,0}$ et transforme $\ket{y,1}$ en
$\ket{\bar y,1}$, avec $\bar y=1-y$ :
\[
  \ket{\Pi_2}=\ket{\beta^{xy}}=\frac{1}{\sqrt2}\Big(\ket{y,0}+(-1)^x\ket{\bar y,1}\Big).
\]
Le premier indice de $\ket{\beta^{xy}}$ est l'entrée de $B$ (le qubit qui reçoit $H$), le
second celle de $A$ (la cible). Les quatre cas :

\begin{center}
\renewcommand{\arraystretch}{1.7}
\begin{tabular}{@{}ccccc@{}}
\toprule
$(B,A)=(x,y)$ & $\ket{\Pi_0}=\ket{y,x}$ & $\ket{\Pi_1}$ & $\ket{\Pi_2}=\ket{\beta^{xy}}$ & Nom \\
\midrule
$(0,0)$ & $\ket{00}$ & $\frac{1}{\sqrt2}(\ket{00}+\ket{01})$ & $\frac{1}{\sqrt2}(\ket{00}+\ket{11})$ & $\ket{\Phi^+}$ \\
$(1,0)$ & $\ket{01}$ & $\frac{1}{\sqrt2}(\ket{00}-\ket{01})$ & $\frac{1}{\sqrt2}(\ket{00}-\ket{11})$ & $\ket{\Phi^-}$ \\
$(0,1)$ & $\ket{10}$ & $\frac{1}{\sqrt2}(\ket{10}+\ket{11})$ & $\frac{1}{\sqrt2}(\ket{10}+\ket{01})$ & $\ket{\Psi^+}$ \\
$(1,1)$ & $\ket{11}$ & $\frac{1}{\sqrt2}(\ket{10}-\ket{11})$ & $\frac{1}{\sqrt2}(\ket{10}-\ket{01})$ & $-\ket{\Psi^-}$ \\
\bottomrule
\end{tabular}
\end{center}
Détail de la dernière ligne : $\ket{\Pi_1}=\ket{1}\otimes\ket{-}=\frac{1}{\sqrt2}(\ket{10}-\ket{11})$ ;
le CNOT laisse $\ket{10}$ et envoie $\ket{11}$ sur $\ket{01}$, d'où
$\frac{1}{\sqrt2}(\ket{10}-\ket{01})=-\frac{1}{\sqrt2}(\ket{01}-\ket{10})=-\ket{\Psi^-}$.

\paragraph{Forme matricielle.} $\mathrm{CNOT}_{21}$ échange les lignes $2$ et $4$ de
$I\otimes H$ (il échange les amplitudes de $\ket{01}$ et $\ket{11}$), donc
\[
  \mathrm{CNOT}_{21}\,(I\otimes H)=\frac{1}{\sqrt2}\begin{pmatrix}
    1 & 1 & 0 & 0\\ 0 & 0 & 1 & -1\\ 0 & 0 & 1 & 1\\ 1 & -1 & 0 & 0
  \end{pmatrix}.
\]
Ses colonnes sont les images de $\ket{00},\ket{01},\ket{10},\ket{11}$, c'est-à-dire de
$(B,A)=(0,0),(1,0),(0,1),(1,1)$ : ce sont $\ket{\beta^{00}}$, $\ket{\beta^{10}}$,
$\ket{\beta^{01}}$, $\ket{\beta^{11}}$, soit $\ket{\Phi^+}$, $\ket{\Phi^-}$, $\ket{\Psi^+}$
et $-\ket{\Psi^-}$.

\begin{remarque}[Comparaison avec la section précédente]
Le circuit de la figure~\ref{fig:bell} applique $H$ au \emph{premier} chiffre et prend celui-ci
pour contrôle ; ici $H$ et le contrôle portent sur le \emph{dernier} chiffre. Les deux
circuits se déduisent l'un de l'autre en échangeant les deux fils (conjugaison par le SWAP,
section~\ref{sec:swap}). Les quatre états de Bell obtenus sont les mêmes, avec les mêmes
indices, sauf le dernier : $-\ket{\Psi^-}$ ici, $+\ket{\Psi^-}$ là-bas. Ce signe est une phase
globale, sans effet sur les mesures (elle disparaît de $\rho$, section~\ref{sec:rho}).
\end{remarque}

\subsubsection{Matrice de densité de l'état de sortie}

Pour $\ket{\Psi}=\ket{\beta^{00}}=\frac{1}{\sqrt2}\ket{00}+\frac{1}{\sqrt2}\ket{11}$ :
\[
  \rho_{AB}=\ket{\Psi}\bra{\Psi}
  =\begin{pmatrix}\frac{1}{\sqrt2}\\0\\0\\\frac{1}{\sqrt2}\end{pmatrix}
   \begin{pmatrix}\frac{1}{\sqrt2}&0&0&\frac{1}{\sqrt2}\end{pmatrix}
  =\frac12\begin{pmatrix}1&0&0&1\\0&0&0&0\\0&0&0&0\\1&0&0&1\end{pmatrix}.
\]
Sous forme de ket-bra :
$\rho_{AB}=\frac12\big(\ket{00}\bra{00}+\ket{00}\bra{11}+\ket{11}\bra{00}+\ket{11}\bra{11}\big)$.
On vérifie $\operatorname{Tr}\rho_{AB}=\frac12(1+1)=1$ et
\[
  \rho_{AB}^2=\frac14\begin{pmatrix}2&0&0&2\\0&0&0&0\\0&0&0&0\\2&0&0&2\end{pmatrix}=\rho_{AB},
\]
donc $\operatorname{Tr}\rho_{AB}^2=1$ : le couple $(A,B)$ est dans un état \emph{pur}.

% ---------------------------------------------------------------------
\subsection{Information : variables aléatoires et états quantiques}\label{sec:info}
% ---------------------------------------------------------------------

\begin{definition}{Entropie, information mutuelle, entropie conditionnelle}
Soit $A$ une variable aléatoire de loi $p(a)$. Son \emph{entropie de Shannon}, en bits, est
\[
  H(A)=-\sum_a p(a)\log_2p(a)\qquad(\text{avec }0\log_20=0).
\]
Elle mesure l'incertitude moyenne sur $A$ : une pièce équilibrée a $H=1$ bit. Pour un couple
$(A,B)$ de loi $p(a,b)$, $H(A,B)=-\sum_{a,b}p(a,b)\log_2p(a,b)$, et
\begin{align*}
  I(A;B)&=H(A)+H(B)-H(A,B) &&\text{(information mutuelle)},\\
  H(A|B)&=H(A,B)-H(B)=H(A)-I(A;B) &&\text{(entropie conditionnelle)}.
\end{align*}
\end{definition}

Les deux écritures de $H(A|B)$ sont égales : la définition de $I(A;B)$ donne
$H(A,B)=H(A)+H(B)-I(A;B)$, donc $H(A,B)-H(B)=H(A)-I(A;B)$. L'information mutuelle $I(A;B)$
mesure ce que la connaissance de $B$ apprend sur $A$ ; $H(A|B)$ est l'incertitude qui reste sur
$A$ quand on connaît $B$.

\subsubsection{Deux variables aléatoires indépendantes}

Alice et Bob lancent chacun une pièce équilibrée, sans lien entre les deux lancers :
$p(a,b)=\frac14$ pour les quatre couples. Chaque pièce donne $0$ ou $1$ avec la probabilité $\frac12$, donc
\begin{align*}
  H(A)&=-\Big(\tfrac12\log_2\tfrac12+\tfrac12\log_2\tfrac12\Big)=1\text{ bit},\\
  H(B)&=1\text{ bit (même calcul)},\\
  H(A,B)&=-4\cdot\tfrac14\log_2\tfrac14=2\text{ bits},\\
  I(A;B)&=H(A)+H(B)-H(A,B)=1+1-2=0,\\
  H(A|B)&=H(A)-I(A;B)=1-0=H(A)=1\text{ bit}.
\end{align*}
Connaître $B$ n'apprend rien sur $A$.

\subsubsection{Deux variables solidaires}

Les deux pièces tombent toujours du même côté : $p(0,0)=p(1,1)=\frac12$ et
$p(0,1)=p(1,0)=0$. Les lois de chaque pièce sont inchangées ($p(a{=}0)=p(0,0)+p(0,1)=\frac12$,
etc.), donc $H(A)=H(B)=1$ bit, mais
\begin{align*}
  H(A,B)&=-2\cdot\tfrac12\log_2\tfrac12-2\cdot0\log_20=1\text{ bit},\\
  I(A;B)&=1+1-1=1\text{ bit},\qquad H(A|B)=1-1=0.
\end{align*}
Connaître $B$ détermine $A$ : plus aucune incertitude ne reste.

\subsubsection{Deux qubits dans un état de Bell}

Pour un système quantique, on remplace la loi de probabilité par la matrice de densité.

\begin{definition}{Trace partielle et entropie de von Neumann}
L'état du seul qubit $A$ d'un système $(A,B)$ de matrice de densité $\rho_{AB}$ est obtenu en
\emph{traçant sur $B$} : $\rho_A=\operatorname{Tr}_B(\rho_{AB})$. La trace partielle est
linéaire et remplace la partie $B$ de chaque terme par sa trace :
\[
  \operatorname{Tr}_B\big(\ket{a}\bra{a'}\otimes\ket{b}\bra{b'}\big)=\ket{a}\bra{a'}\;\braket{b'}{b}.
\]
L'\emph{entropie de von Neumann} d'un état $\rho$, de valeurs propres $\lambda_i$, est
$S(\rho)=-\operatorname{Tr}(\rho\log_2\rho)=-\sum_i\lambda_i\log_2\lambda_i$. En quantique,
$H(A)=S(\rho_A)$, $H(B)=S(\rho_B)$ et $H(A,B)=S(\rho_{AB})$ ; les définitions de
$I(A;B)$ et de $H(A|B)$ sont les mêmes.
\end{definition}

On part de l'état de Bell
$\rho_{AB}=\frac12\big(\ket{00}\bra{00}+\ket{00}\bra{11}+\ket{11}\bra{00}+\ket{11}\bra{11}\big)$
produit par le circuit de la section~\ref{sec:circuit-bell}. Chaque terme est un produit
d'opérateurs à un qubit :
$\ket{00}\bra{11}=\ket{0}\bra{1}\otimes\ket{0}\bra{1}$, etc. Alors
\begin{align*}
  \rho_A=\operatorname{Tr}_B(\rho_{AB})
  &=\tfrac12\Big(\ket{0}\bra{0}\,\braket{0}{0}+\ket{0}\bra{1}\,\braket{1}{0}
      +\ket{1}\bra{0}\,\braket{0}{1}+\ket{1}\bra{1}\,\braket{1}{1}\Big)\\
  &=\tfrac12\Big(\ket{0}\bra{0}\cdot1+\ket{0}\bra{1}\cdot0+\ket{1}\bra{0}\cdot0+\ket{1}\bra{1}\cdot1\Big)\\
  &=\tfrac12\big(\ket{0}\bra{0}+\ket{1}\bra{1}\big)=\tfrac12\,\mathrm{Id}
  =\begin{pmatrix}\frac12&0\\0&\frac12\end{pmatrix}.
\end{align*}
La formule des composantes de la section~\ref{sec:separable},
$(\rho_A)_{ij}=\rho_{(i0),(j0)}+\rho_{(i1),(j1)}$, donne le même résultat. En traçant sur $A$
au lieu de $B$, la règle est $\operatorname{Tr}_A(\ket{a}\bra{a'}\otimes\ket{b}\bra{b'})
=\braket{a'}{a}\,\ket{b}\bra{b'}$, et on trouve de même $\rho_B=\frac12\mathrm{Id}$.

Calcul des entropies :
\begin{itemize}[nosep]
  \item $\rho_{AB}$ est pur : ses valeurs propres sont $1,0,0,0$, donc
        $H(A,B)=S(\rho_{AB})=-1\cdot\log_21=0$ ;
  \item $\rho_A$ et $\rho_B$ ont pour valeurs propres $\frac12,\frac12$, donc
        $H(A)=H(B)=-\big(\frac12\log_2\frac12+\frac12\log_2\frac12\big)=1$ bit.
\end{itemize}
D'où
\begin{align*}
  I(A;B)&=H(A)+H(B)-H(A,B)=1+1-0=2\text{ bits},\\
  H(A|B)&=H(A,B)-H(B)=0-1=-1\text{ bit}\qquad(=H(A)-I(A;B)=1-2).
\end{align*}

\begin{remarque}[Une entropie conditionnelle négative]
Pour des variables classiques, $p(a,b)\leq p(b)$, donc $-\log_2p(a,b)\geq-\log_2p(b)$ et
$H(A,B)\geq H(B)$ : on a toujours $H(A|B)\geq0$, et $I(A;B)\leq H(A)$. Ici $H(A|B)=-1$ et
$I(A;B)=2>H(A)$ : ces valeurs n'ont pas d'équivalent classique, c'est une signature de
l'intrication. Le résultat est surprenant mais mathématiquement correct.
\end{remarque}

\begin{center}
\renewcommand{\arraystretch}{1.4}
\begin{tabular}{@{}lccc@{}}
\toprule
 & indépendantes & solidaires & état de Bell \\
\midrule
$H(A)$ & $1$ & $1$ & $1$ \\
$H(B)$ & $1$ & $1$ & $1$ \\
$H(A,B)$ & $2$ & $1$ & $0$ \\
$I(A;B)$ & $0$ & $1$ & $2$ \\
$H(A|B)$ & $1$ & $0$ & $-1$ \\
\bottomrule
\end{tabular}
\end{center}
(en bits ; les deux premières colonnes sont classiques, la troisième est quantique).

% ---------------------------------------------------------------------
\subsection{Système pur et système mélangé}\label{sec:pur-melange}
% ---------------------------------------------------------------------

Rappel (section~\ref{sec:pur-mixte}) : si $\lambda_i$ sont les valeurs propres de la matrice de
densité $\rho$, alors $\operatorname{Tr}\rho^2=\sum_i\lambda_i^2\leq1$. Le système est
\emph{pur} si $\operatorname{Tr}\rho^2=1$ (alors $\rho=\ket{\psi}\bra{\psi}$ et $\rho^2=\rho$),
et \emph{mélangé} si $\operatorname{Tr}\rho^2<1$. Pour un qubit,
$\operatorname{Tr}\rho^2=\frac{1+|\vec r|^2}{2}$ (section~\ref{sec:bloch}) est compris entre
$\frac12$ et $1$.

\subsubsection{Un système pur : exercice}

\emph{Énoncé.} Soit $\ket{\psi}=\frac{1}{\sqrt2}\ket{0}+\frac{1}{\sqrt2}\ket{1}$.
\begin{enumerate}[label=\alph*), nosep]
  \item Calculer la matrice de densité $\rho$.
  \item Calculer la pureté du système quantique.
  \item Calculer $\langle\sigma_x\rangle$ et $\langle\sigma_z\rangle$.
  \item Écrire $\ket{\psi}$ dans la base $X$.
\end{enumerate}

\paragraph{a) Matrice de densité.} $\rho=\ket{\psi}\bra{\psi}$ est le produit d'un vecteur
colonne par un vecteur ligne (les amplitudes sont réelles) :
\[
  \rho=\begin{pmatrix}\frac{1}{\sqrt2}\\[0.4ex]\frac{1}{\sqrt2}\end{pmatrix}
       \begin{pmatrix}\frac{1}{\sqrt2}&\frac{1}{\sqrt2}\end{pmatrix}
  =\begin{pmatrix}\frac12&\frac12\\[0.4ex]\frac12&\frac12\end{pmatrix}.
\]

\paragraph{b) Pureté.}
\[
  \rho^2=\begin{pmatrix}\frac12&\frac12\\[0.4ex]\frac12&\frac12\end{pmatrix}
         \begin{pmatrix}\frac12&\frac12\\[0.4ex]\frac12&\frac12\end{pmatrix}
  =\begin{pmatrix}\frac14+\frac14&\frac14+\frac14\\[0.4ex]\frac14+\frac14&\frac14+\frac14\end{pmatrix}
  =\begin{pmatrix}\frac12&\frac12\\[0.4ex]\frac12&\frac12\end{pmatrix}=\rho,
\]
donc $\operatorname{Tr}\rho^2=\operatorname{Tr}\rho=\frac12+\frac12=1$ : le système est \emph{pur}.

\paragraph{c) Valeurs moyennes.} Avec $\sigma_x=X=\begin{pmatrix}0&1\\1&0\end{pmatrix}$ et
$\sigma_z=Z=\begin{pmatrix}1&0\\0&-1\end{pmatrix}$ (section~\ref{sec:op-dirac}), la valeur
moyenne est $\langle\sigma\rangle=\operatorname{Tr}(\rho\,\sigma)$ (section~\ref{sec:rho}).
\begin{align*}
  \rho\,\sigma_x&=\begin{pmatrix}\frac12&\frac12\\[0.4ex]\frac12&\frac12\end{pmatrix}
     \begin{pmatrix}0&1\\1&0\end{pmatrix}
     =\begin{pmatrix}\frac12&\frac12\\[0.4ex]\frac12&\frac12\end{pmatrix},
     &\langle\sigma_x\rangle&=\operatorname{Tr}(\rho\,\sigma_x)=\tfrac12+\tfrac12=1,\\[1.5ex]
  \rho\,\sigma_z&=\begin{pmatrix}\frac12&\frac12\\[0.4ex]\frac12&\frac12\end{pmatrix}
     \begin{pmatrix}1&0\\0&-1\end{pmatrix}
     =\begin{pmatrix}\frac12&-\frac12\\[0.4ex]\frac12&-\frac12\end{pmatrix},
     &\langle\sigma_z\rangle&=\operatorname{Tr}(\rho\,\sigma_z)=\tfrac12-\tfrac12=0.
\end{align*}
On retrouve $\langle\sigma_z\rangle$ avec les probabilités de mesure : dans la base $Z$,
$\sigma_z$ vaut $+1$ pour $\ket{0}$ et $-1$ pour $\ket{1}$, et
$|\braket{0}{\psi}|^2=|\braket{1}{\psi}|^2=\frac12$, donc
$\langle\sigma_z\rangle=\frac12\times(+1)+\frac12\times(-1)=0$.

\paragraph{d) Écriture dans la base $X$.} La base $X$ est $\{\ket{+},\ket{-}\}$, orthonormée, donc
$\ket{+}\bra{+}+\ket{-}\bra{-}=\mathrm{Id}$ et
$\ket{\psi}=\ket{+}\braket{+}{\psi}+\ket{-}\braket{-}{\psi}$, avec
\[
  \braket{+}{\psi}=\tfrac{1}{\sqrt2}\Big(\tfrac{1}{\sqrt2}+\tfrac{1}{\sqrt2}\Big)=1,
  \qquad
  \braket{-}{\psi}=\tfrac{1}{\sqrt2}\Big(\tfrac{1}{\sqrt2}-\tfrac{1}{\sqrt2}\Big)=0 .
\]
Donc $\ket{\psi}=1\cdot\ket{+}+0\cdot\ket{-}=\ket{+}$ : dans la base $X$, l'état n'a qu'une
seule composante. Comme $\sigma_x\ket{+}=+\ket{+}$ et $\sigma_x\ket{-}=-\ket{-}$, une
mesure de $\sigma_x$ donne $+1$ avec certitude, et
$\langle\sigma_x\rangle=1\cdot(+1)+0\cdot(-1)=1$ : c'est le résultat de la question c).
Le vecteur de Bloch de $\ket{+}$ est $(\langle\sigma_x\rangle,\langle\sigma_y\rangle,\langle\sigma_z\rangle)=(1,0,0)$,
de norme $1$ : l'état est sur la sphère, en accord avec $\operatorname{Tr}\rho^2=\frac{1+1}{2}=1$.

\subsubsection{Un système mélangé : un qubit de la paire de Bell}

À la section~\ref{sec:info}, la trace sur $B$ de l'état de Bell a donné
$\rho_A=\frac12\mathrm{Id}=\begin{pmatrix}\frac12&0\\0&\frac12\end{pmatrix}$. Alors
\[
  \rho_A^2=\begin{pmatrix}\frac14&0\\0&\frac14\end{pmatrix},
  \qquad
  \operatorname{Tr}\big(\rho_A^2\big)=\frac14+\frac14=\frac12<1 :
\]
le qubit $A$ est dans un état \emph{mélangé}, et c'est le plus mélangé possible pour un qubit
(centre de la boule de Bloch, $\vec r=\vec0$). Pourtant le couple $(A,B)$ est pur
($\operatorname{Tr}\rho_{AB}^2=1$) : on connaît parfaitement l'état du couple, mais pas celui d'un
qubit pris seul. L'information est entièrement dans les corrélations entre $A$ et $B$, et
$S(\rho_A)=1$ bit.

\subsubsection{Un système mélangé : le mélange statistique de Bob}

Bob prépare un qubit dans l'état $\ket{0}$ avec la probabilité $p_1=\frac34$ (75\,\%) et dans
l'état $\ket{1}$ avec la probabilité $p_2=\frac14$ (25\,\%), sans que l'on sache lequel des deux a
été préparé. Ce n'est pas un état pur $\ket{\psi}$ : on décrit le système par la matrice de
densité $\rho=\sum_ip_i\ket{\psi_i}\bra{\psi_i}$ (section~\ref{sec:pur-mixte}) :
\[
  \rho=\underbrace{\tfrac34}_{p_1}\ket{0}\bra{0}+\underbrace{\tfrac14}_{p_2}\ket{1}\bra{1}
  =\begin{pmatrix}\frac34&0\\0&\frac14\end{pmatrix}.
\]
Alors
\[
  \rho^2=\begin{pmatrix}\frac{9}{16}&0\\0&\frac{1}{16}\end{pmatrix},
  \qquad
  \operatorname{Tr}\big(\rho^2\big)=\frac{9}{16}+\frac{1}{16}=\frac{10}{16}=\frac58<1 :
\]
le système est \emph{mélangé}.

Les valeurs moyennes sont les moyennes pondérées sur la préparation :
$\langle\sigma_z\rangle=\operatorname{Tr}(\rho\,\sigma_z)=\frac34-\frac14=\frac12
=p_1\cdot(+1)+p_2\cdot(-1)$, et $\langle\sigma_x\rangle=\operatorname{Tr}(\rho\,\sigma_x)=0$
(les termes hors diagonale de $\rho$ sont nuls). Le vecteur de Bloch est $(0,0,\frac12)$,
à l'intérieur de la boule, et $\operatorname{Tr}\rho^2=\frac{1+1/4}{2}=\frac58$. Son entropie
est $S(\rho)=-\frac34\log_2\frac34-\frac14\log_2\frac14=2-\frac34\log_23\approx0{,}811$ bit.

\begin{remarque}[Mélange et superposition]
Considérons l'état pur $\ket{\psi'}=\frac{\sqrt3}{2}\ket{0}+\frac12\ket{1}$, qui a les mêmes
probabilités en base $Z$ ($P(0)=\frac34$, $P(1)=\frac14$). Sa matrice de densité est
\[
  \rho'=\begin{pmatrix}\frac34&\frac{\sqrt3}{4}\\[0.4ex]\frac{\sqrt3}{4}&\frac14\end{pmatrix},
  \qquad
  \operatorname{Tr}\rho'^2=\frac{9}{16}+\frac1{16}+2\cdot\frac{3}{16}=1,
  \qquad
  \langle\sigma_x\rangle=2\cdot\frac{\sqrt3}{4}=\frac{\sqrt3}{2}.
\]
Le mélange de Bob a la même diagonale mais aucun terme hors diagonale ($\operatorname{Tr}\rho^2=\frac58$,
$\langle\sigma_x\rangle=0$) : une superposition n'est pas un tirage au hasard entre $\ket{0}$ et
$\ket{1}$. De plus, pour un état pur il existe une mesure au résultat certain (ici $\sigma_x$ pour
$\ket{+}$) ; pour le mélange de Bob, la meilleure probabilité d'obtenir un même résultat est
$75\,\%$ (mesure en base $Z$, résultat $\ket{0}$).
\end{remarque}

\subsubsection{Bilan}

\begin{center}
\renewcommand{\arraystretch}{1.5}
\begin{tabular}{@{}lcccl@{}}
\toprule
\textbf{Système} & $\rho$ & $\operatorname{Tr}\rho^2$ & $S(\rho)$ & \textbf{Nature} \\
\midrule
$\ket{\psi}=\ket{+}$ & ${\renewcommand{\arraystretch}{1}\begin{pmatrix}\frac12&\frac12\\\frac12&\frac12\end{pmatrix}}$ & $1$ & $0$ & pur \\[2ex]
Qubit $A$ d'une paire de Bell & ${\renewcommand{\arraystretch}{1}\begin{pmatrix}\frac12&0\\0&\frac12\end{pmatrix}}$ & $\frac12$ & $1$ bit & mélangé \\[2ex]
Source de Bob (75\,\%/25\,\%) & ${\renewcommand{\arraystretch}{1}\begin{pmatrix}\frac34&0\\0&\frac14\end{pmatrix}}$ & $\frac58$ & $\approx0{,}811$ bit & mélangé \\
\bottomrule
\end{tabular}
\end{center}

% ---------------------------------------------------------------------
\subsection{Algorithme quantique : la téléportation}\label{sec:teleportation}
% ---------------------------------------------------------------------

\emph{Problème.} Alice détient un qubit $Q$ dans un état qu'elle ne connaît pas,
$\ket{\varphi}=\alpha\ket{0}+\beta\ket{1}$ avec $|\alpha|^2+|\beta|^2=1$, et elle veut que Bob
obtienne cet état sur l'un de ses qubits. Elle ne peut pas mesurer $Q$ pour connaître
$\alpha$ et $\beta$ (une mesure ne donne qu'un bit et détruit l'état), elle ne peut pas le
copier (théorème de non-clonage) et elle n'envoie pas le qubit lui-même.

\begin{definition}{Protocole de téléportation quantique}
\textbf{Ressources.}
\begin{itemize}[nosep]
  \item Un \emph{e-bit} : Alice et Bob partagent une paire de qubits $(A,B)$ dans l'état de Bell
        $\ket{\Phi^+}_{BA}=\frac{1}{\sqrt2}\big(\ket{00}+\ket{11}\big)$ ; Alice détient $A$,
        Bob détient $B$.
  \item Le qubit $Q$ à téléporter, chez Alice.
  \item Un canal classique (deux bits).
\end{itemize}
Le système est écrit $(B,A,Q)$ : $\ket{b,a,q}=\ket{b}_B\otimes\ket{a}_A\otimes\ket{q}_Q$.

\textbf{Étapes.}
\begin{enumerate}[nosep]
  \item Alice applique un CNOT de contrôle $Q$ et de cible $A$.
  \item Alice applique une porte $H$ sur $Q$.
  \item Alice mesure $A$ et $Q$ dans la base de calcul, obtient deux bits $(a,q)$ et les envoie
        à Bob.
  \item Bob applique $X$ si $a=1$, puis $Z$ si $q=1$, à son qubit $B$.
\end{enumerate}
\end{definition}

\begin{figure}[!ht]
\centering
\begin{adjustbox}{max width=\linewidth}
\begin{tikzpicture}[>=Stealth, font=\small]
  % fils quantiques : Q en haut, A au milieu, B en bas
  \draw (0,2.4) -- (5.3,2.4);
  \draw (0,1.2) -- (5.3,1.2);
  \draw (0,0) -- (10.6,0);
  \node[anchor=east] at (0,2.4) {$\ket{\varphi}_Q$};
  \node[anchor=south west, text=insagris, font=\footnotesize] at (0.05,1.2) {$A$};
  \node[anchor=south west, text=insagris, font=\footnotesize] at (0.05,0) {$B$};
  % paire de Bell partagée
  \draw[insagris, thick] (-0.1,1.5) -- (-0.25,1.5) -- (-0.25,-0.3) -- (-0.1,-0.3);
  \node[anchor=east] at (-0.35,0.6) {$\ket{\Phi^+}_{BA}$};
  % CNOT : contrôle Q, cible A
  \draw[insarouge, thick] (2.6,2.4) -- (2.6,1.2);
  \pic at (2.6,2.4) {ctl};
  \pic at (2.6,1.2) {cible};
  % Hadamard sur Q
  \node[porte] at (4.2,2.4) {$H$};
  % mesures d'Alice
  \pic at (5.6,2.4) {mesure};
  \pic at (5.6,1.2) {mesure};
  % fils classiques (doubles)
  \draw[double, double distance=1.2pt, insagris] (5.9,1.2) -- (7.6,1.2) -- (7.6,0.35);
  \draw[double, double distance=1.2pt, insagris] (5.9,2.4) -- (9.6,2.4) -- (9.6,0.35);
  \node[above, font=\footnotesize] at (6.5,1.2) {$a$};
  \node[above, font=\footnotesize] at (6.5,2.4) {$q$};
  % corrections de Bob
  \node[porte] at (7.6,0) {$X$};
  \node[porte] at (9.6,0) {$Z$};
  \node[anchor=west] at (10.6,0) {$\ket{\varphi}_B$};
  % états intermédiaires
  \foreach \x in {1.4,3.4,4.9}
    \draw[dashed, insagris!70] (\x,2.8) -- (\x,-0.5);
  \node[anchor=north, font=\footnotesize] at (1.4,-0.5) {$\ket{\Pi_0}$};
  \node[anchor=north, font=\footnotesize] at (3.4,-0.5) {$\ket{\Pi_1}$};
  \node[anchor=north, font=\footnotesize] at (4.9,-0.5) {$\ket{\Pi_2}$};
\end{tikzpicture}
\end{adjustbox}
\caption{Circuit de téléportation. Alice détient $Q$ (dans $\ket{\varphi}=\alpha\ket{0}+\beta\ket{1}$)
et $A$, Bob détient $B$ ; $A$ et $B$ forment la paire $\ket{\Phi^+}_{BA}$. Les traits doubles
sont des fils \emph{classiques} qui portent les résultats $a$ (mesure de $A$) et $q$ (mesure de
$Q$) ; $X$ est appliqué si $a=1$, $Z$ si $q=1$. Le ket s'écrit $\ket{b,a,q}$ : le fil du haut
($Q$) est le dernier chiffre.}
\label{fig:teleportation}
\end{figure}

\subsubsection{Calcul pas à pas}

\paragraph{État initial $\ket{\Pi_0}$.} Alice et Bob partagent la paire, et $Q$ est indépendant
d'elle :
\begin{align*}
  \ket{\Pi_0}&=\ket{\Phi^+}_{BA}\otimes\ket{\varphi}_Q
    =\frac{1}{\sqrt2}\big(\ket{00}+\ket{11}\big)\otimes\big(\alpha\ket{0}+\beta\ket{1}\big)\\
  &=\frac{1}{\sqrt2}\Big(\alpha\ket{00}\ket{0}+\beta\ket{00}\ket{1}+\alpha\ket{11}\ket{0}+\beta\ket{11}\ket{1}\Big)\\
  &=\frac{\alpha\ket{000}+\beta\ket{001}+\alpha\ket{110}+\beta\ket{111}}{\sqrt2}.
\end{align*}

\paragraph{Après le CNOT $Q\to A$ : $\ket{\Pi_1}$.} Le CNOT (contrôle $Q$, dernier chiffre ;
cible $A$, chiffre du milieu) agit ainsi : $\ket{b,a,q}\to\ket{b,\,a\oplus q,\,q}$. Terme par terme :
\begin{center}
\renewcommand{\arraystretch}{1.3}
\begin{tabular}{@{}cccc@{}}
\toprule
Terme & $q$ & Effet sur $A$ & Résultat \\
\midrule
$\alpha\ket{000}$ & $0$ & aucun & $\alpha\ket{000}$ \\
$\beta\ket{001}$ & $1$ & $0\to1$ & $\beta\ket{011}$ \\
$\alpha\ket{110}$ & $0$ & aucun & $\alpha\ket{110}$ \\
$\beta\ket{111}$ & $1$ & $1\to0$ & $\beta\ket{101}$ \\
\bottomrule
\end{tabular}
\end{center}
donc
\[
  \ket{\Pi_1}=\frac{\alpha\ket{000}+\beta\ket{011}+\alpha\ket{110}+\beta\ket{101}}{\sqrt2}.
\]
On isole ensuite le dernier chiffre, celui de $Q$, qui va changer à l'étape suivante
(le premier ket est celui du couple $\ket{b,a}$) :
\[
  \ket{\Pi_1}=\frac{\alpha\ket{00}\otimes\ket{0}+\beta\ket{01}\otimes\ket{1}
    +\alpha\ket{11}\otimes\ket{0}+\beta\ket{10}\otimes\ket{1}}{\sqrt2}.
\]

\paragraph{Après la porte $H$ sur $Q$ : $\ket{\Pi_2}$.} On remplace $\ket{0}$ par
$H\ket{0}=\ket{+}$ et $\ket{1}$ par $H\ket{1}=\ket{-}$ dans le facteur $Q$ :
\[
  \ket{\Pi_2}=\frac{1}{\sqrt2}\Big(\alpha\ket{00}\otimes\ket{+}+\beta\ket{01}\otimes\ket{-}
    +\alpha\ket{11}\otimes\ket{+}+\beta\ket{10}\otimes\ket{-}\Big).
\]
Avec $\ket{\pm}=\frac{1}{\sqrt2}(\ket{0}\pm\ket{1})$, le préfacteur devient
$\frac{1}{\sqrt2}\cdot\frac{1}{\sqrt2}=\frac12$ :
\begin{align*}
  \ket{\Pi_2}&=\frac12\Big(\alpha\ket{00}\big(\ket{0}+\ket{1}\big)+\beta\ket{01}\big(\ket{0}-\ket{1}\big)
     +\alpha\ket{11}\big(\ket{0}+\ket{1}\big)+\beta\ket{10}\big(\ket{0}-\ket{1}\big)\Big)\\
  &=\frac12\Big(\alpha\ket{000}+\alpha\ket{001}+\beta\ket{010}-\beta\ket{011}
     +\alpha\ket{110}+\alpha\ket{111}+\beta\ket{100}-\beta\ket{101}\Big).
\end{align*}

\paragraph{Regroupement selon le résultat d'Alice.} Alice mesurera $A$ et $Q$, donc les deux
derniers chiffres. On regroupe les huit termes qui ont les mêmes chiffres $(a,q)$, en mettant
le ket de Bob en premier :
\begin{center}
\renewcommand{\arraystretch}{1.3}
\begin{tabular}{@{}cll@{}}
\toprule
$(a,q)$ & Termes de $\ket{\Pi_2}$ (facteur $\frac12$ omis) & Ket de Bob \\
\midrule
$(0,0)$ & $\alpha\ket{000}+\beta\ket{100}$ & $(\alpha\ket{0}+\beta\ket{1})\ket{00}$ \\
$(0,1)$ & $\alpha\ket{001}-\beta\ket{101}$ & $(\alpha\ket{0}-\beta\ket{1})\ket{01}$ \\
$(1,0)$ & $\beta\ket{010}+\alpha\ket{110}$ & $(\beta\ket{0}+\alpha\ket{1})\ket{10}$ \\
$(1,1)$ & $-\beta\ket{011}+\alpha\ket{111}$ & $(-\beta\ket{0}+\alpha\ket{1})\ket{11}$ \\
\bottomrule
\end{tabular}
\end{center}
d'où
\[
  \ket{\Pi_2}=\frac12\Big[\big(\alpha\ket{0}+\beta\ket{1}\big)\ket{00}
    +\big(\alpha\ket{0}-\beta\ket{1}\big)\ket{01}
    +\big(\beta\ket{0}+\alpha\ket{1}\big)\ket{10}
    +\big(\alpha\ket{1}-\beta\ket{0}\big)\ket{11}\Big].
\]
Dans chaque terme, le ket de gauche est l'état de $B$ et le ket de droite $\ket{a,q}$ celui du
couple $(A,Q)$ d'Alice.

\subsubsection{Mesure d'Alice et corrections de Bob}

Alice mesure $A$ et $Q$. Comme à la section~\ref{sec:mesure2}, la probabilité d'un résultat
$(a,q)$ est le carré de la norme du terme correspondant, et l'état de Bob devient le ket
associé (déjà normé). Ici
\[
  P(a,q)=\Big\|\tfrac12\ket{\phi_{aq}}\Big\|^2=\tfrac14\big(|\alpha|^2+|\beta|^2\big)=\tfrac14
  \quad\text{pour les quatre résultats,}
\]
qui sont donc équiprobables et \emph{indépendants} de $\alpha$ et $\beta$ : les deux bits
envoyés ne révèlent rien sur $\ket{\varphi}$. Les kets de Bob sont
$\ket{\phi_{00}}=\alpha\ket{0}+\beta\ket{1}$, $\ket{\phi_{01}}=\alpha\ket{0}-\beta\ket{1}$,
$\ket{\phi_{10}}=\beta\ket{0}+\alpha\ket{1}$ et $\ket{\phi_{11}}=-\beta\ket{0}+\alpha\ket{1}$.

\begin{center}
\renewcommand{\arraystretch}{1.5}
\begin{tabular}{@{}cccll@{}}
\toprule
Résultat $(a,q)$ & Probabilité & État de Bob & Lien avec $\ket{\varphi}$ & Correction de Bob \\
\midrule
$(0,0)$ & $\frac14$ & $\alpha\ket{0}+\beta\ket{1}$ & $\ket{\varphi}$ & $I$ (rien) \\
$(0,1)$ & $\frac14$ & $\alpha\ket{0}-\beta\ket{1}$ & $Z\ket{\varphi}$ & $Z$ \\
$(1,0)$ & $\frac14$ & $\beta\ket{0}+\alpha\ket{1}$ & $X\ket{\varphi}$ & $X$ \\
$(1,1)$ & $\frac14$ & $-\beta\ket{0}+\alpha\ket{1}$ & $XZ\ket{\varphi}$ & $X$ puis $Z$ (opérateur $ZX$) \\
\bottomrule
\end{tabular}
\end{center}

\paragraph{Vérification des quatre corrections.} Avec $X\ket{0}=\ket{1}$, $X\ket{1}=\ket{0}$,
$Z\ket{0}=\ket{0}$ et $Z\ket{1}=-\ket{1}$ :
\begin{align*}
  (0,0):\quad & I\big(\alpha\ket{0}+\beta\ket{1}\big)=\alpha\ket{0}+\beta\ket{1},\\
  (0,1):\quad & Z\big(\alpha\ket{0}-\beta\ket{1}\big)=\alpha\ket{0}-\beta Z\ket{1}=\alpha\ket{0}+\beta\ket{1},\\
  (1,0):\quad & X\big(\beta\ket{0}+\alpha\ket{1}\big)=\beta\ket{1}+\alpha\ket{0}=\alpha\ket{0}+\beta\ket{1},\\
  (1,1):\quad & X\big(-\beta\ket{0}+\alpha\ket{1}\big)=-\beta\ket{1}+\alpha\ket{0}=\alpha\ket{0}-\beta\ket{1},
     \ \text{puis}\ Z\big(\alpha\ket{0}-\beta\ket{1}\big)=\alpha\ket{0}+\beta\ket{1}.
\end{align*}
Dans le dernier cas, avec les matrices,
$ZX=\begin{pmatrix}1&0\\0&-1\end{pmatrix}\begin{pmatrix}0&1\\1&0\end{pmatrix}
=\begin{pmatrix}0&1\\-1&0\end{pmatrix}$ et
$ZX\begin{pmatrix}-\beta\\ \alpha\end{pmatrix}=\begin{pmatrix}\alpha\\ \beta\end{pmatrix}$.
L'ordre compte : l'état de Bob est $XZ\ket{\varphi}$ ($Z\ket{\varphi}=\alpha\ket{0}-\beta\ket{1}$,
puis $X$ donne $\alpha\ket{1}-\beta\ket{0}$), et on l'annule avec
$(XZ)^{-1}=Z^{-1}X^{-1}=ZX$ : on applique d'abord $X$, puis $Z$, comme dans le circuit.
Dans les quatre cas, Bob retrouve exactement $\ket{\varphi}=\alpha\ket{0}+\beta\ket{1}$.

\subsubsection{Ce que Bob sait avant de recevoir les bits}

Tant que Bob ignore $(a,q)$, son qubit est dans le mélange des quatre kets $\ket{\phi_{aq}}$, chacun
avec la probabilité $\frac14$ :
\begin{align*}
  \rho_B&=\frac14\sum_{a,q}\ket{\phi_{aq}}\bra{\phi_{aq}}\\
  &=\frac14\Bigg[\begin{pmatrix}|\alpha|^2&\alpha\beta^*\\ \alpha^*\beta&|\beta|^2\end{pmatrix}
    +\begin{pmatrix}|\alpha|^2&-\alpha\beta^*\\ -\alpha^*\beta&|\beta|^2\end{pmatrix}
    +\begin{pmatrix}|\beta|^2&\alpha^*\beta\\ \alpha\beta^*&|\alpha|^2\end{pmatrix}
    +\begin{pmatrix}|\beta|^2&-\alpha^*\beta\\ -\alpha\beta^*&|\alpha|^2\end{pmatrix}\Bigg]\\
  &=\frac14\begin{pmatrix}2(|\alpha|^2+|\beta|^2)&0\\0&2(|\alpha|^2+|\beta|^2)\end{pmatrix}
   =\frac12\begin{pmatrix}1&0\\0&1\end{pmatrix}=\frac12\,\mathrm{Id}.
\end{align*}
Sans les deux bits, Bob a donc un qubit complètement mélangé, qui ne dépend pas de $\alpha$ ni de
$\beta$ : il n'a reçu aucune information sur $\ket{\varphi}$. C'est l'envoi classique de $(a,q)$
qui permet de la retrouver, donc rien ne voyage plus vite que la lumière.

\begin{remarque}[Bilan de la téléportation]
\begin{itemize}[nosep]
  \item Ressources : $1$ e-bit (la paire de Bell) et $2$ bits classiques pour transmettre
        \emph{un qubit} d'état quelconque, sans qu'Alice connaisse $\alpha$ ni $\beta$.
  \item Pas de clonage : après la mesure, $A$ et $Q$ sont dans l'état de base $\ket{a,q}$.
        Alice n'a plus $\ket{\varphi}$ : l'état a été \emph{déplacé}, pas copié. L'intrication
        entre $A$ et $B$ a été consommée.
  \item Les quatre corrections $I,Z,X,ZX$ sont des portes contrôlées par des bits classiques
        (section~\ref{sec:controlees}).
\end{itemize}
\end{remarque}

% ---------------------------------------------------------------------
\subsection{Portes contrôlées : contrôlé-\texorpdfstring{$U$}{U}, Fredkin et Toffoli}\label{sec:controlees}
% ---------------------------------------------------------------------

Le CNOT est le cas le plus simple d'une idée générale : appliquer une opération à un ou
plusieurs qubits \emph{cibles} seulement si un ou plusieurs qubits de \emph{contrôle} valent $1$.
Dans cette section, comme à la section~\ref{sec:op2}, le contrôle est le \emph{premier} chiffre
du ket. Les qubits sont notés $x$, $y$, $z$ (le qubit $x$ est le contrôle) et
$\ket{x,y,z}=\ket{x}\otimes\ket{y}\otimes\ket{z}$ ; d'après la convention de la section, le
contrôle est dessiné en bas (figure~\ref{fig:controlees}). Pour éviter la confusion avec les
matrices de Pauli, on écrit ici $\sigma_x=X$ pour la porte NOT.

\begin{figure}[!ht]
\centering
\begin{adjustbox}{max width=\linewidth}
\begin{tikzpicture}[>=Stealth, font=\small]
  % (a) contrôlé-U : contrôle x en bas, cible y en haut
  \draw (0,1.2) -- (3.4,1.2);
  \draw (0,0) -- (3.4,0);
  \node[anchor=east] at (0,1.2) {$y$};
  \node[anchor=east] at (0,0) {$x$};
  \draw[insarouge, thick] (1.7,0) -- (1.7,0.85);
  \pic at (1.7,0) {ctl};
  \node[porte] at (1.7,1.2) {$U$};
  \node[anchor=north, font=\footnotesize] at (1.7,-0.5) {(a) contrôlé-$U$};
  % (b) contrôlé-SWAP (Fredkin) : x contrôle, y et z échangés
  \draw (5,2.4) -- (8.4,2.4);
  \draw (5,1.2) -- (8.4,1.2);
  \draw (5,0) -- (8.4,0);
  \node[anchor=east] at (5,2.4) {$z$};
  \node[anchor=east] at (5,1.2) {$y$};
  \node[anchor=east] at (5,0) {$x$};
  \draw[insarouge, thick] (6.7,0) -- (6.7,2.4);
  \pic at (6.7,0) {ctl};
  \pic at (6.7,1.2) {croix};
  \pic at (6.7,2.4) {croix};
  \node[anchor=north, font=\footnotesize] at (6.7,-0.5) {(b) Fredkin (C-SWAP)};
  % (c) Toffoli (CCNOT) : x et y contrôlent, z cible
  \draw (10,2.4) -- (13.4,2.4);
  \draw (10,1.2) -- (13.4,1.2);
  \draw (10,0) -- (13.4,0);
  \node[anchor=east] at (10,2.4) {$z$};
  \node[anchor=east] at (10,1.2) {$y$};
  \node[anchor=east] at (10,0) {$x$};
  \draw[insarouge, thick] (11.7,0) -- (11.7,2.4);
  \pic at (11.7,0) {ctl};
  \pic at (11.7,1.2) {ctl};
  \pic at (11.7,2.4) {cible};
  \node[anchor=north, font=\footnotesize] at (11.7,-0.5) {(c) Toffoli (CCNOT)};
\end{tikzpicture}
\end{adjustbox}
\caption{Portes contrôlées. Le contrôle $x$ est le premier chiffre du ket, donc le fil du bas.
(a) Le qubit $y$ subit $U$ si $x=1$ ; pour $U=X$ on retrouve le CNOT. (b) Si $x=1$, les qubits
$y$ et $z$ sont échangés. (c) Si $x=y=1$, le qubit $z$ est inversé.}
\label{fig:controlees}
\end{figure}

\subsubsection{Contrôlé-\texorpdfstring{$U$}{U}}

\begin{definition}{Porte contrôlée-$U$}
Soit $U$ une matrice unitaire $2\times2$. Sur le système $(x,y)$, de contrôle $x$ et de cible $y$,
\[
  \mathrm{C}U=\ket{0}\bra{0}_x\otimes I_y+\ket{1}\bra{1}_x\otimes U_y .
\]
Le premier terme agit quand le contrôle est $\ket{0}$ : la cible ne change pas ($I$). Le second
agit quand le contrôle est $\ket{1}$ : la cible subit $U$.
\end{definition}

Comme $\ket{0}\bra{0}\ket{x}=\delta_{x0}\ket{0}$ et $\ket{1}\bra{1}\ket{x}=\delta_{x1}\ket{1}$,
\[
  \mathrm{C}U\big(\ket{x}\otimes\ket{y}\big)=\ket{x}\otimes U^x\ket{y},\qquad U^0=I,\ U^1=U .
\]
Par la règle $A\otimes B=(A_{ij}B)$, $\ket{0}\bra{0}\otimes I=\begin{pmatrix}I&0\\0&0\end{pmatrix}$
et $\ket{1}\bra{1}\otimes U=\begin{pmatrix}0&0\\0&U\end{pmatrix}$ (blocs $2\times2$), donc
\[
  \mathrm{C}U=\begin{pmatrix}I&0\\0&U\end{pmatrix}
  =\begin{pmatrix}1&0&0&0\\0&1&0&0\\0&0&u_{00}&u_{01}\\0&0&u_{10}&u_{11}\end{pmatrix}
  \qquad\text{pour }U=\begin{pmatrix}u_{00}&u_{01}\\u_{10}&u_{11}\end{pmatrix}.
\]
Cette matrice est unitaire car $I^\dagger I=I$ et $U^\dagger U=I$.

\paragraph{Le CNOT est un cas particulier.} Pour $U=\sigma_x$,
$\mathrm{C}\sigma_x=\ket{0}\bra{0}_x\otimes I_y+\ket{1}\bra{1}_x\otimes\sigma_x$ est le CNOT de
contrôle $x$ (premier chiffre), c'est-à-dire $\mathrm{CNOT}_{12}$ de la section~\ref{sec:cnot}.
On calcule l'action sur les quatre états de base, avec $\braket{0}{0}=\braket{1}{1}=1$,
$\braket{0}{1}=\braket{1}{0}=0$, $\sigma_x\ket{0}=\ket{1}$ et $\sigma_x\ket{1}=\ket{0}$ :
\begin{align*}
  \mathrm{C}\sigma_x\ket{00}&=\big[\ket{0}\bra{0}_x\otimes I_y\big]\ket{0}_x\ket{0}_y
      +\big[\ket{1}\bra{1}_x\otimes\sigma_x\big]\ket{0}_x\ket{0}_y\\
    &=\ket{0}\braket{0}{0}\otimes I\ket{0}+\ket{1}\braket{1}{0}\otimes\sigma_x\ket{0}
     =\ket{0}\otimes\ket{0}+0=\ket{00},\\[1.2ex]
  \mathrm{C}\sigma_x\ket{01}&=\ket{0}\braket{0}{0}\otimes I\ket{1}+\ket{1}\braket{1}{0}\otimes\sigma_x\ket{1}
     =\ket{0}\otimes\ket{1}+0=\ket{01},\\[1.2ex]
  \mathrm{C}\sigma_x\ket{10}&=\ket{0}\braket{0}{1}\otimes I\ket{0}+\ket{1}\braket{1}{1}\otimes\sigma_x\ket{0}
     =0+\ket{1}\otimes\ket{1}=\ket{11},\\[1.2ex]
  \mathrm{C}\sigma_x\ket{11}&=\ket{0}\braket{0}{1}\otimes I\ket{1}+\ket{1}\braket{1}{1}\otimes\sigma_x\ket{1}
     =0+\ket{1}\otimes\ket{0}=\ket{10}.
\end{align*}
On retrouve la table du CNOT : la cible n'est inversée que si le contrôle vaut $1$.

\paragraph{Un autre exemple : $U=Z$.} Ici
\[
  \mathrm{C}Z=\begin{pmatrix}I&0\\0&Z\end{pmatrix}=\mathrm{diag}(1,1,1,-1) :
\]
la porte multiplie $\ket{11}$ par $-1$ et ne touche pas les autres états de base. Elle est
symétrique entre les deux qubits. Elle se déduit du CNOT par
$\mathrm{C}Z=(I\otimes H)\,\mathrm{C}X\,(I\otimes H)$, où $\mathrm{C}X$ est le contrôlé-$X$.
En effet, $I\otimes H=\mathrm{diag}(H,H)$ (deux blocs $2\times2$ égaux à $H$), et le produit de
matrices diagonales par blocs se fait bloc par bloc :
\[
  (I\otimes H)\begin{pmatrix}I&0\\0&X\end{pmatrix}(I\otimes H)
  =\begin{pmatrix}H\,I\,H&0\\0&H\,X\,H\end{pmatrix}
  =\begin{pmatrix}I&0\\0&HXH\end{pmatrix},
\]
car $H^2=I$, et
\[
  HXH=\frac12\begin{pmatrix}1&1\\1&-1\end{pmatrix}\begin{pmatrix}0&1\\1&0\end{pmatrix}\begin{pmatrix}1&1\\1&-1\end{pmatrix}
  =\frac12\begin{pmatrix}1&1\\1&-1\end{pmatrix}\begin{pmatrix}1&-1\\1&1\end{pmatrix}
  =\frac12\begin{pmatrix}2&0\\0&-2\end{pmatrix}=Z
\]
(section~\ref{sec:hadamard} : $H$ échange les axes $x$ et $z$).

\subsubsection{Contrôlé-SWAP : la porte de Fredkin}

\begin{definition}{Porte de Fredkin (C-SWAP)}
Sur trois qubits $(x,y,z)$ (espace de dimension $2^3=8$), avec $x$ pour contrôle,
\[
  \mathrm{C\text{-}SWAP}=\ket{0}\bra{0}_x\otimes I_y\otimes I_z+\ket{1}\bra{1}_x\otimes\mathrm{SWAP}_{yz}.
\]
Si $x=0$, rien ne change ; si $x=1$, les qubits $y$ et $z$ sont échangés.
\end{definition}

Calcul sur quatre états de base (avec $I\ket{ab}=\ket{ab}$ et $\mathrm{SWAP}\ket{ab}=\ket{ba}$) :
\begin{align*}
  \mathrm{C\text{-}SWAP}\ket{000}&=\ket{0}\braket{0}{0}\otimes I\ket{00}+\ket{1}\braket{1}{0}\otimes\mathrm{SWAP}\ket{00}
     =\ket{0}\otimes\ket{00}+0=\ket{000},\\[1.2ex]
  \mathrm{C\text{-}SWAP}\ket{101}&=\ket{0}\braket{0}{1}\otimes I\ket{01}+\ket{1}\braket{1}{1}\otimes\mathrm{SWAP}\ket{01}
     =0+\ket{1}\otimes\ket{10}=\ket{110},\\[1.2ex]
  \mathrm{C\text{-}SWAP}\ket{110}&=\ket{0}\braket{0}{1}\otimes I\ket{10}+\ket{1}\braket{1}{1}\otimes\mathrm{SWAP}\ket{10}
     =0+\ket{1}\otimes\ket{01}=\ket{101},\\[1.2ex]
  \mathrm{C\text{-}SWAP}\ket{010}&=\ket{0}\braket{0}{0}\otimes I\ket{10}+\ket{1}\braket{1}{0}\otimes\mathrm{SWAP}\ket{10}
     =\ket{0}\otimes\ket{10}+0=\ket{010}.
\end{align*}
Les huit états de base donnent :
\begin{center}
\renewcommand{\arraystretch}{1.3}
\begin{tabular}{@{}ccccccccc@{}}
\toprule
Entrée & $\ket{000}$ & $\ket{001}$ & $\ket{010}$ & $\ket{011}$ & $\ket{100}$ & $\ket{101}$ & $\ket{110}$ & $\ket{111}$ \\
\midrule
Sortie & $\ket{000}$ & $\ket{001}$ & $\ket{010}$ & $\ket{011}$ & $\ket{100}$ & $\ket{110}$ & $\ket{101}$ & $\ket{111}$ \\
\bottomrule
\end{tabular}
\end{center}
Seuls $\ket{101}$ et $\ket{110}$ changent (contrôle à $1$ et $y\neq z$) ; pour $\ket{100}$ et
$\ket{111}$ l'échange de deux qubits égaux ne change rien. Dans la base
$\ket{000},\ket{001},\dots,\ket{111}$ (ordre binaire), la matrice est
$\mathrm{diag}(I_4,\mathrm{SWAP})$ (deux blocs $4\times4$) :
\[
  \mathrm{C\text{-}SWAP}=\begingroup\setlength{\arraycolsep}{3pt}\begin{pmatrix}
    1&0&0&0&0&0&0&0\\
    0&1&0&0&0&0&0&0\\
    0&0&1&0&0&0&0&0\\
    0&0&0&1&0&0&0&0\\
    0&0&0&0&1&0&0&0\\
    0&0&0&0&0&0&1&0\\
    0&0&0&0&0&1&0&0\\
    0&0&0&0&0&0&0&1
  \end{pmatrix}\endgroup .
\]
C'est l'identité dont les lignes $6$ et $7$ sont échangées.

\paragraph{Contrôle en superposition.} Avec $\ket{+}_x\otimes\ket{0}_y\otimes\ket{1}_z
=\frac{1}{\sqrt2}\big(\ket{001}+\ket{101}\big)$ :
\[
  \mathrm{C\text{-}SWAP}\,\frac{1}{\sqrt2}\big(\ket{001}+\ket{101}\big)
  =\frac{1}{\sqrt2}\big(\ket{001}+\ket{110}\big).
\]
Le résultat n'est pas un produit : les coefficients de $\ket{x}\otimes\ket{yz}$ forment la
matrice (lignes $x=0,1$, colonnes $yz=00,01,10,11$)
\[
  \frac{1}{\sqrt2}\begin{pmatrix}0&1&0&0\\0&0&1&0\end{pmatrix},
\]
dont les deux lignes ne sont pas proportionnelles. Le qubit $x$ est intriqué avec la paire
$(y,z)$.

\subsubsection{Contrôlé-contrôlé-NOT : la porte de Toffoli}

\begin{definition}{Porte de Toffoli (CCNOT)}
Sur trois qubits $(x,y,z)$, avec $x$ et $y$ pour contrôles et $z$ pour cible,
\[
  \mathrm{CCNOT}=\ket{0}\bra{0}_x\otimes I_y\otimes I_z
   +\ket{1}\bra{1}_x\otimes\Big[\ket{0}\bra{0}_y\otimes I_z+\ket{1}\bra{1}_y\otimes\sigma_x\Big].
\]
Le crochet est un CNOT (contrôle $y$, cible $z$), appliqué seulement si $x=1$ : si $x=0$, rien ;
si $x=1$ et $y=0$, rien ; si $x=1$ et $y=1$, le qubit $z$ est inversé.
\end{definition}

Calcul de $\mathrm{CCNOT}\ket{110}$ :
\begin{align*}
  \mathrm{CCNOT}\ket{110}&=\ket{0}\braket{0}{1}\otimes I\ket{1}\otimes I\ket{0}
     +\ket{1}\braket{1}{1}\otimes\Big[\ket{0}\braket{0}{1}\otimes I\ket{0}
       +\ket{1}\braket{1}{1}\otimes\sigma_x\ket{0}\Big]\\
  &=0+\ket{1}\otimes\Big[0+\ket{1}\otimes\ket{1}\Big]=\ket{111}.
\end{align*}
Pour les huit états de base, $\mathrm{CCNOT}\ket{x,y,z}=\ket{x,\,y,\,z\oplus xy}$ :
\begin{center}
\renewcommand{\arraystretch}{1.3}
\begin{tabular}{@{}ccccccccc@{}}
\toprule
Entrée & $\ket{000}$ & $\ket{001}$ & $\ket{010}$ & $\ket{011}$ & $\ket{100}$ & $\ket{101}$ & $\ket{110}$ & $\ket{111}$ \\
\midrule
Sortie & $\ket{000}$ & $\ket{001}$ & $\ket{010}$ & $\ket{011}$ & $\ket{100}$ & $\ket{101}$ & $\ket{111}$ & $\ket{110}$ \\
\bottomrule
\end{tabular}
\end{center}
Seuls $\ket{110}$ et $\ket{111}$ sont échangés : la matrice est $\mathrm{diag}(I_6,X)$
(un bloc $6\times6$ égal à l'identité, puis $X$),
\[
  \mathrm{CCNOT}=\begingroup\setlength{\arraycolsep}{3pt}\begin{pmatrix}
    1&0&0&0&0&0&0&0\\
    0&1&0&0&0&0&0&0\\
    0&0&1&0&0&0&0&0\\
    0&0&0&1&0&0&0&0\\
    0&0&0&0&1&0&0&0\\
    0&0&0&0&0&1&0&0\\
    0&0&0&0&0&0&0&1\\
    0&0&0&0&0&0&1&0
  \end{pmatrix}\endgroup .
\]

\begin{remarque}[Portes réversibles]
Les portes de Fredkin et de Toffoli sont des permutations des états de base qui échangent
seulement deux d'entre eux : elles sont leurs propres inverses ($\mathrm{C\text{-}SWAP}^2=\mathrm{CCNOT}^2=\mathrm{Id}$)
et unitaires. Avec $z=0$, la Toffoli calcule le ET logique $x\wedge y$ sans rien effacer :
$\ket{x,y,0}\to\ket{x,y,xy}$ ; avec $z=1$, elle calcule le NON-ET. Elle sert à écrire tout calcul
classique de façon réversible.
\end{remarque}

% ---------------------------------------------------------------------
\subsection{Récapitulatif du chapitre}
% ---------------------------------------------------------------------

\begin{center}
\renewcommand{\arraystretch}{1.5}
\begin{tabular}{@{}lp{12cm}@{}}
\toprule
\textbf{Notion} & \textbf{Formule ou résultat} \\
\midrule
CNOT & $\ket{b,a}\to\ket{b\oplus a,a}$ (contrôle $a$ écrit en dernier), $\mathrm{CNOT}^2=\mathrm{Id}$ \\
SWAP & $\mathrm{SWAP}\ket{a,b}=\ket{b,a}$, \quad $\mathrm{CNOT}_{21}=\mathrm{SWAP}\,\mathrm{CNOT}_{12}\,\mathrm{SWAP}$ \\
Circuit de Bell & $(I\otimes H)$ puis CNOT : $\ket{\beta^{xy}}=\frac{1}{\sqrt2}\big(\ket{y,0}+(-1)^x\ket{\bar y,1}\big)$, $\ket{\beta^{00}}=\ket{\Phi^+}$ \\
Information & $I(A;B)=H(A)+H(B)-H(A,B)$, \quad $H(A|B)=H(A,B)-H(B)$ \\
Trace partielle & $\operatorname{Tr}_B\big(\ket{a}\bra{a'}\otimes\ket{b}\bra{b'}\big)=\ket{a}\bra{a'}\braket{b'}{b}$ ; état de Bell : $\rho_A=\frac12\mathrm{Id}$ \\
État de Bell (bits) & $H(A,B)=0$, \quad $H(A)=H(B)=1$, \quad $I(A;B)=2$, \quad $H(A|B)=-1$ \\
Pur / mélangé & $\operatorname{Tr}\rho^2=1$ / $<1$ ; $\ket{+}$ : $1$ ; qubit d'une paire de Bell : $\frac12$ ; source de Bob : $\frac58$ \\
Téléportation & $1$ e-bit $+\ 2$ bits classiques ; correction : $X$ si $a=1$, puis $Z$ si $q=1$ ; $P(a,q)=\frac14$ \\
Contrôlé-$U$ & $\ket{0}\bra{0}\otimes I+\ket{1}\bra{1}\otimes U=\mathrm{diag}(I,U)$ ; $\ket{x,y}\to\ket{x}\otimes U^x\ket{y}$ \\
Fredkin & $\ket{0}\bra{0}\otimes I\otimes I+\ket{1}\bra{1}\otimes\mathrm{SWAP}$ ; $\ket{101}\leftrightarrow\ket{110}$ \\
Toffoli & $\ket{x,y,z}\to\ket{x,y,z\oplus xy}$ ; $\ket{110}\leftrightarrow\ket{111}$ \\
\bottomrule
\end{tabular}
\end{center}
.
%  Comme chapitre4.tex, chapitre5.tex et chapitre6.tex, il ouvre une NOUVELLE
%  SECTION (numéro 5). Pour le rattacher à la section précédente, remplacer
%  \section par \subsection et chaque \subsection par \subsubsection.
%  Pas de préambule : environnements (definition, resultat, remarque),
%  macros (\ii, \ee, \dd, \ket, \bra, \braket) et couleurs (insarouge,
%  insagris) sont définis dans main.tex. Les symboles de circuits (contrôle,
%  cible, croix, mesure) sont définis juste après avec \tikzset.
% =====================================================================

% Macros de Dirac (déjà définies dans main.tex récent ; ce bloc évite l'erreur
% « Undefined control sequence \ket » si main.tex est une ancienne version).
\providecommand{\ket}[1]{|#1\rangle}
\providecommand{\bra}[1]{\langle #1|}
\providecommand{\braket}[2]{\langle #1|#2\rangle}

% Symboles de circuits pour les figures de ce chapitre.
\tikzset{
  ctl/.pic   = {\fill[insarouge] (0,0) circle (2.6pt);},
  cible/.pic = {\draw[insarouge, thick, fill=white] (0,0) circle (0.28);
                \draw[insarouge, thick] (-0.28,0) -- (0.28,0);
                \draw[insarouge, thick] (0,-0.28) -- (0,0.28);},
  croix/.pic = {\draw[insarouge, thick] (-0.17,-0.17) -- (0.17,0.17);
                \draw[insarouge, thick] (-0.17,0.17) -- (0.17,-0.17);},
  mesure/.pic = {\draw[thick, fill=white] (-0.3,-0.27) rectangle (0.3,0.27);
                 \draw[thick] (-0.19,-0.13) arc[start angle=155, end angle=25, radius=0.21];
                 \draw[thick, -{Stealth[length=1.3mm]}] (0,-0.17) -- (0.14,0.12);},
  porte/.style = {draw=insarouge, fill=insarouge!8, minimum size=7mm, thick},
}

% =====================================================================
\section{Portes multi-qubits, circuits et téléportation}\label{sec:circuits}
% =====================================================================

La section précédente a décrit les états de deux qubits. On étudie ici les portes qui
agissent sur plusieurs qubits : le CNOT et le SWAP, puis les portes contrôlées à deux et à
trois qubits (contrôlé-$U$, Fredkin, Toffoli). On assemble ensuite des portes dans un
circuit qui fabrique un état de Bell, on mesure l'information contenue dans un tel état,
on distingue les états purs des états mélangés, et on démontre le protocole de
téléportation. Tous les calculs sont détaillés.

\begin{remarque}[Convention pour les circuits]
Dans un circuit, le temps va de gauche à droite et chaque qubit a un fil. Les fils sont
dessinés dans l'ordre \emph{inverse} de l'écriture des kets : le fil du \emph{haut} porte le
\emph{dernier} chiffre du ket, le fil du \emph{bas} le premier. Par exemple, si le qubit $A$
est dessiné au-dessus du qubit $B$, le ket s'écrit $\ket{b,a}=\ket{b}_B\otimes\ket{a}_A$.
La figure~\ref{fig:bell} de la section précédente suit la même convention.
\end{remarque}

% ---------------------------------------------------------------------
\subsection{Porte CNOT}\label{sec:cnot}
% ---------------------------------------------------------------------

\begin{definition}{Porte CNOT (NOT contrôlé)}
La porte CNOT agit sur deux qubits : le qubit $A$ est le \emph{contrôle}, le qubit $B$ est la
\emph{cible}. La cible est inversée (porte $X$) si le contrôle vaut $1$ et reste inchangée
s'il vaut $0$. Avec le système écrit $(B,A)$, soit $\ket{b,a}=\ket{b}_B\otimes\ket{a}_A$ :
\[
  \mathrm{CNOT}\,\ket{b,a}=\ket{b\oplus a,\,a},
  \qquad
  \mathrm{CNOT}=I\otimes\ket{0}\bra{0}+X\otimes\ket{1}\bra{1},
\]
où $\oplus$ est l'addition modulo $2$ ($0\oplus0=1\oplus1=0$ et $0\oplus1=1\oplus0=1$).
\end{definition}

Sur le circuit de la figure~\ref{fig:cnot}, le point noir posé sur le fil de $A$ marque le
contrôle et le symbole $\oplus$ posé sur le fil de $B$ marque la cible.

\begin{figure}[!ht]
\centering
\begin{tikzpicture}[>=Stealth, font=\small]
  \draw (0,1.2) -- (6,1.2);
  \draw (0,0) -- (6,0);
  \node[anchor=east] at (0,1.2) {$\ket{a}_A$};
  \node[anchor=east] at (0,0) {$\ket{b}_B$};
  \node[anchor=west] at (6,1.2) {$\ket{a}_A$};
  \node[anchor=west] at (6,0) {$\ket{b\oplus a}_B$};
  \draw[insarouge, thick] (3,1.2) -- (3,0);
  \pic at (3,1.2) {ctl};
  \pic at (3,0) {cible};
  \node[anchor=south, text=insagris] at (3,1.5) {qubit de contrôle};
  \node[anchor=north, text=insagris] at (3,-0.5) {qubit cible};
\end{tikzpicture}
\caption{Porte CNOT : le fil du haut ($A$) est le contrôle, le fil du bas ($B$) est la
cible. Le ket s'écrit $\ket{b,a}$ (le fil du haut donne le dernier chiffre).}
\label{fig:cnot}
\end{figure}

\subsubsection{Table de vérité}

\begin{center}
\renewcommand{\arraystretch}{1.3}
\begin{tabular}{@{}cccc@{}}
\toprule
Entrée $\ket{b,a}$ & Contrôle $a$ & Sortie $\ket{b\oplus a,a}$ & Effet sur la cible $B$ \\
\midrule
$\ket{00}$ & $0$ & $\ket{00}$ & aucun \\
$\ket{01}$ & $1$ & $\ket{11}$ & $0\to1$ \\
$\ket{10}$ & $0$ & $\ket{10}$ & aucun \\
$\ket{11}$ & $1$ & $\ket{01}$ & $1\to0$ \\
\bottomrule
\end{tabular}
\end{center}

\subsubsection{Matrice du CNOT}

La $j$-ième colonne d'une matrice $4\times4$ est l'image du $j$-ième vecteur de base
($\ket{00},\ket{01},\ket{10},\ket{11}$, dans cet ordre). La table donne :
\begin{align*}
  \mathrm{CNOT}\ket{00}&=\ket{00}=\begin{pmatrix}1\\0\\0\\0\end{pmatrix},&
  \mathrm{CNOT}\ket{01}&=\ket{11}=\begin{pmatrix}0\\0\\0\\1\end{pmatrix},\\[1.5ex]
  \mathrm{CNOT}\ket{10}&=\ket{10}=\begin{pmatrix}0\\0\\1\\0\end{pmatrix},&
  \mathrm{CNOT}\ket{11}&=\ket{01}=\begin{pmatrix}0\\1\\0\\0\end{pmatrix},
\end{align*}
d'où, en juxtaposant ces quatre colonnes,
\[
  \mathrm{CNOT}=\begin{pmatrix}
    1 & 0 & 0 & 0\\ 0 & 0 & 0 & 1\\ 0 & 0 & 1 & 0\\ 0 & 1 & 0 & 0
  \end{pmatrix}.
\]

On retrouve cette matrice avec la formule $I\otimes\ket{0}\bra{0}+X\otimes\ket{1}\bra{1}$ et
la règle $A\otimes B=(A_{ij}B)$ (section~\ref{sec:op2}), où
$\ket{0}\bra{0}=\begin{pmatrix}1&0\\0&0\end{pmatrix}$ et
$\ket{1}\bra{1}=\begin{pmatrix}0&0\\0&1\end{pmatrix}$ :
\[
  I\otimes\ket{0}\bra{0}=\begin{pmatrix}1&0&0&0\\0&0&0&0\\0&0&1&0\\0&0&0&0\end{pmatrix},
  \qquad
  X\otimes\ket{1}\bra{1}=\begin{pmatrix}0&0&0&0\\0&0&0&1\\0&0&0&0\\0&1&0&0\end{pmatrix},
\]
et la somme de ces deux matrices est la matrice du CNOT.

\subsubsection{Action sur un état quelconque}

Pour $\ket{\Psi}=c_{00}\ket{00}+c_{01}\ket{01}+c_{10}\ket{10}+c_{11}\ket{11}$ :
\[
  \mathrm{CNOT}\begin{pmatrix}c_{00}\\c_{01}\\c_{10}\\c_{11}\end{pmatrix}
  =\begin{pmatrix}
    1\cdot c_{00}+0\cdot c_{01}+0\cdot c_{10}+0\cdot c_{11}\\
    0\cdot c_{00}+0\cdot c_{01}+0\cdot c_{10}+1\cdot c_{11}\\
    0\cdot c_{00}+0\cdot c_{01}+1\cdot c_{10}+0\cdot c_{11}\\
    0\cdot c_{00}+1\cdot c_{01}+0\cdot c_{10}+0\cdot c_{11}
  \end{pmatrix}
  =\begin{pmatrix}c_{00}\\c_{11}\\c_{10}\\c_{01}\end{pmatrix}.
\]
Le CNOT échange les amplitudes de $\ket{01}$ et de $\ket{11}$, c'est-à-dire des deux états
dont le contrôle vaut $1$, et laisse les autres.

\paragraph{Exemple : contrôle en superposition.} Soit $\ket{0}_B\otimes\ket{+}_A
=\frac{1}{\sqrt2}\big(\ket{00}+\ket{01}\big)=\frac{1}{\sqrt2}(1,1,0,0)^{\mathsf T}$,
un état séparable. Alors
\[
  \mathrm{CNOT}\,\frac{1}{\sqrt2}\begin{pmatrix}1\\1\\0\\0\end{pmatrix}
  =\frac{1}{\sqrt2}\begin{pmatrix}1\\0\\0\\1\end{pmatrix}
  =\frac{1}{\sqrt2}\big(\ket{00}+\ket{11}\big)=\ket{\Phi^+}.
\]
Le résultat est intriqué ($D=c_{00}c_{11}-c_{01}c_{10}=\frac12\neq0$ ; section~\ref{sec:separable}) :
un CNOT dont le contrôle est en superposition crée de l'intrication.

\paragraph{Réversibilité.} Deux CNOT successifs redonnent l'état initial :
$\ket{b,a}\to\ket{b\oplus a,a}\to\ket{b\oplus a\oplus a,a}=\ket{b,a}$, donc
$\mathrm{CNOT}^2=\mathrm{Id}$. La matrice est réelle et symétrique, donc
$\mathrm{CNOT}^\dagger\mathrm{CNOT}=\mathrm{CNOT}^2=\mathrm{Id}$ : elle est bien unitaire.

\begin{remarque}[Contrôle en dernier ou en premier]
La matrice dépend de l'ordre d'écriture du ket. Notons $\mathrm{CNOT}_{21}$ le CNOT dont le
contrôle est le \emph{second} chiffre du ket (c'est la matrice ci-dessus) et
$\mathrm{CNOT}_{12}$ celui dont le contrôle est le \emph{premier} chiffre, utilisé à la
section précédente :
\[
  \mathrm{CNOT}_{12}=\ket{0}\bra{0}\otimes I+\ket{1}\bra{1}\otimes X
  =\begin{pmatrix}1&0&0&0\\0&1&0&0\\0&0&0&1\\0&0&1&0\end{pmatrix}.
\]
Un même circuit (contrôle $A$, cible $B$) a donc deux matrices, selon que l'on écrit
$A$ en dernier ($\mathrm{CNOT}_{21}$) ou en premier ($\mathrm{CNOT}_{12}$). Les deux sont
reliées par le SWAP (section suivante).
\end{remarque}

% ---------------------------------------------------------------------
\subsection{Porte SWAP}\label{sec:swap}
% ---------------------------------------------------------------------

\begin{definition}{Porte SWAP (échange)}
La porte SWAP échange les états des deux qubits :
\[
  \mathrm{SWAP}\big(\ket{\psi}\otimes\ket{\phi}\big)=\ket{\phi}\otimes\ket{\psi}.
\]
Sur les états de base, $\mathrm{SWAP}\ket{a,b}=\ket{b,a}$ ; on prolonge par linéarité.
\end{definition}

Son symbole est formé de deux croix reliées par un trait vertical
(figure~\ref{fig:swap}).

\begin{figure}[!ht]
\centering
\begin{tikzpicture}[>=Stealth, font=\small]
  \draw (0,1.2) -- (6,1.2);
  \draw (0,0) -- (6,0);
  \node[anchor=east] at (0,1.2) {$\ket{a}$};
  \node[anchor=east] at (0,0) {$\ket{b}$};
  \node[anchor=west] at (6,1.2) {$\ket{b}$};
  \node[anchor=west] at (6,0) {$\ket{a}$};
  \draw[insarouge, thick] (3,1.2) -- (3,0);
  \pic at (3,1.2) {croix};
  \pic at (3,0) {croix};
\end{tikzpicture}
\caption{Porte SWAP : les deux qubits échangent leurs états.}
\label{fig:swap}
\end{figure}

\subsubsection{Table et matrice}

\begin{center}
\renewcommand{\arraystretch}{1.3}
\begin{tabular}{@{}cc@{}}
\toprule
Entrée & Sortie \\
\midrule
$\ket{00}$ & $\ket{00}$ \\
$\ket{01}$ & $\ket{10}$ \\
$\ket{10}$ & $\ket{01}$ \\
$\ket{11}$ & $\ket{11}$ \\
\bottomrule
\end{tabular}
\end{center}

Les images des quatre vecteurs de base sont les colonnes de la matrice :
\[
  \mathrm{SWAP}\ket{00}=\begin{pmatrix}1\\0\\0\\0\end{pmatrix},\quad
  \mathrm{SWAP}\ket{01}=\begin{pmatrix}0\\0\\1\\0\end{pmatrix},\quad
  \mathrm{SWAP}\ket{10}=\begin{pmatrix}0\\1\\0\\0\end{pmatrix},\quad
  \mathrm{SWAP}\ket{11}=\begin{pmatrix}0\\0\\0\\1\end{pmatrix},
\]
\[
  \mathrm{SWAP}=\begin{pmatrix}
    1 & 0 & 0 & 0\\ 0 & 0 & 1 & 0\\ 0 & 1 & 0 & 0\\ 0 & 0 & 0 & 1
  \end{pmatrix}.
\]
Appliquée à $\ket{01}=(0,1,0,0)^{\mathsf T}$, cette matrice sélectionne sa deuxième colonne :
\[
  \begin{pmatrix}
    1 & 0 & 0 & 0\\ 0 & 0 & 1 & 0\\ 0 & 1 & 0 & 0\\ 0 & 0 & 0 & 1
  \end{pmatrix}
  \begin{pmatrix}0\\1\\0\\0\end{pmatrix}
  =\begin{pmatrix}0\\0\\1\\0\end{pmatrix}=\ket{10}.
\]
Sur un état quelconque, $\mathrm{SWAP}\,(c_{00},c_{01},c_{10},c_{11})^{\mathsf T}
=(c_{00},c_{10},c_{01},c_{11})^{\mathsf T}$ : les amplitudes de $\ket{01}$ et $\ket{10}$
sont échangées.

\subsubsection{Accord avec la définition}

Pour deux états $\ket{\psi}=(a_0,a_1)^{\mathsf T}$ et $\ket{\phi}=(b_0,b_1)^{\mathsf T}$ :
\[
  \mathrm{SWAP}\,(\ket{\psi}\otimes\ket{\phi})
  =\mathrm{SWAP}\begin{pmatrix}a_0b_0\\a_0b_1\\a_1b_0\\a_1b_1\end{pmatrix}
  =\begin{pmatrix}a_0b_0\\a_1b_0\\a_0b_1\\a_1b_1\end{pmatrix}
  =\begin{pmatrix}b_0\\b_1\end{pmatrix}\otimes\begin{pmatrix}a_0\\a_1\end{pmatrix}
  =\ket{\phi}\otimes\ket{\psi},
\]
ce qui est bien la définition. La porte s'écrit aussi avec des projecteurs :
\[
  \mathrm{SWAP}=\sum_{a,b\in\{0,1\}}\ket{a}\bra{b}\otimes\ket{b}\bra{a}.
\]
En effet, pour $\ket{c}\otimes\ket{d}$ on trouve
$\sum_{a,b}\ket{a}\braket{b}{c}\otimes\ket{b}\braket{a}{d}$, et seul le terme $b=c$, $a=d$
est non nul : le résultat est $\ket{d}\otimes\ket{c}$.

\paragraph{Exemples.}
\begin{itemize}[nosep]
  \item Deux états différents : $\mathrm{SWAP}\big(\ket{0}\otimes\ket{+}\big)
        =\mathrm{SWAP}\,\frac{1}{\sqrt2}(1,1,0,0)^{\mathsf T}=\frac{1}{\sqrt2}(1,0,1,0)^{\mathsf T}
        =\ket{+}\otimes\ket{0}$.
  \item L'état $\ket{\psi}=\frac{1}{\sqrt2}\ket{00}+\frac{1}{\sqrt2}\ket{11}$ ne change pas :
        $\mathrm{SWAP}\ket{\psi}=\frac{1}{\sqrt2}\ket{00}+\frac{1}{\sqrt2}\ket{11}=\ket{\psi}$,
        car $\ket{00}$ et $\ket{11}$ sont invariants. De même
        $\mathrm{SWAP}\ket{\Phi^-}=\ket{\Phi^-}$ et $\mathrm{SWAP}\ket{\Psi^+}=\ket{\Psi^+}$, alors que
        $\mathrm{SWAP}\ket{\Psi^-}=\frac{1}{\sqrt2}(\ket{10}-\ket{01})=-\ket{\Psi^-}$.
\end{itemize}

\paragraph{Propriétés.} Échanger deux fois redonne l'état de départ :
$\mathrm{SWAP}^2=\mathrm{Id}$ ; la matrice est réelle et symétrique, donc
$\mathrm{SWAP}^\dagger\mathrm{SWAP}=\mathrm{SWAP}^2=\mathrm{Id}$ : elle est unitaire.

\subsubsection{Lien avec le CNOT}

\begin{resultat}{SWAP échange les rôles de contrôle et de cible}
\[
  \mathrm{CNOT}_{21}=\mathrm{SWAP}\;\mathrm{CNOT}_{12}\;\mathrm{SWAP}.
\]
\end{resultat}

\emph{Sur les états de base.}
$\mathrm{SWAP}\,\mathrm{CNOT}_{12}\,\mathrm{SWAP}\ket{b,a}
=\mathrm{SWAP}\,\mathrm{CNOT}_{12}\ket{a,b}
=\mathrm{SWAP}\ket{a,\,a\oplus b}=\ket{a\oplus b,\,a}=\mathrm{CNOT}_{21}\ket{b,a}$.

\emph{Par les matrices.} Multiplier à gauche par $\mathrm{SWAP}$ échange les lignes $2$ et
$3$, multiplier à droite l'échange les colonnes $2$ et $3$ :
\[
  \mathrm{SWAP}\,\mathrm{CNOT}_{12}
  =\begin{pmatrix}1&0&0&0\\0&0&0&1\\0&1&0&0\\0&0&1&0\end{pmatrix},
  \qquad
  \mathrm{SWAP}\,\mathrm{CNOT}_{12}\,\mathrm{SWAP}
  =\begin{pmatrix}1&0&0&0\\0&0&0&1\\0&0&1&0\\0&1&0&0\end{pmatrix}=\mathrm{CNOT}_{21}.
\]

\begin{remarque}[SWAP avec trois CNOT]
$\mathrm{SWAP}=\mathrm{CNOT}_{12}\,\mathrm{CNOT}_{21}\,\mathrm{CNOT}_{12}$. En effet, sur
$\ket{a,b}$ (on applique les portes de droite à gauche) :
\[
  \ket{a,b}\ \xrightarrow{\ \mathrm{CNOT}_{12}\ }\ \ket{a,\,a\oplus b}
  \ \xrightarrow{\ \mathrm{CNOT}_{21}\ }\ \ket{a\oplus(a\oplus b),\,a\oplus b}=\ket{b,\,a\oplus b}
  \ \xrightarrow{\ \mathrm{CNOT}_{12}\ }\ \ket{b,\,b\oplus(a\oplus b)}=\ket{b,a}.
\]
Les deux membres coïncident sur les états de base, donc sur tous les états par linéarité.
\end{remarque}

% ---------------------------------------------------------------------
\subsection{Un circuit quantique : le générateur d'états de Bell}\label{sec:circuit-bell}
% ---------------------------------------------------------------------

On assemble une porte de Hadamard et un CNOT. Le système est $(A,B)$ avec le ket
$\ket{a,b}=\ket{a}_A\otimes\ket{b}_B$ ; d'après la convention de la section, $B$ est dessiné
en haut et $A$ en bas (figure~\ref{fig:circuit-bell}). Le circuit a deux étapes :
\begin{enumerate}[nosep]
  \item la porte $I\otimes H$ : identité sur $A$, Hadamard sur $B$ ;
  \item un CNOT de contrôle $B$ et de cible $A$ (la matrice $\mathrm{CNOT}_{21}$ de la
        section~\ref{sec:cnot}, puisque le contrôle est le second chiffre).
\end{enumerate}
Les traits pointillés délimitent les états $\ket{\Pi_0}$, $\ket{\Pi_1}$ et $\ket{\Pi_2}$ :
avant les portes, après $I\otimes H$, après le CNOT.

\begin{figure}[!ht]
\centering
\begin{tikzpicture}[>=Stealth, font=\small]
  % fils : B en haut, A en bas
  \draw (0,1.4) -- (7.6,1.4);
  \draw (0,0) -- (7.6,0);
  \node[anchor=east] at (0,1.4) {$\ket{0}_B$};
  \node[anchor=east] at (0,0) {$\ket{0}_A$};
  % première étape : H sur B, identité (pointillée) sur A
  \node[porte] at (2.2,1.4) {$H$};
  \node[draw=insagris, dashed, fill=white, minimum size=7mm] at (2.2,0) {$I$};
  % CNOT : contrôle B (haut), cible A (bas)
  \draw[insarouge, thick] (5,1.4) -- (5,0);
  \pic at (5,1.4) {ctl};
  \pic at (5,0) {cible};
  % accolade de sortie
  \draw[insagris, thick] (7.75,1.6) -- (7.9,1.6) -- (7.9,-0.2) -- (7.75,-0.2);
  \node[anchor=west] at (7.9,0.7) {$\ket{\beta^{00}}_{AB}$};
  % états intermédiaires
  \foreach \x in {0.7,3.6,6.4}
    \draw[dashed, insagris!70] (\x,1.9) -- (\x,-0.7);
  \node[anchor=north, font=\footnotesize] at (0.7,-0.7) {$\ket{\Pi_0}=\ket{00}$};
  \node[anchor=north, font=\footnotesize] at (3.6,-0.7) {$\ket{\Pi_1}=\ket{0}\otimes\ket{+}$};
  \node[anchor=north, font=\footnotesize] at (6.4,-0.7) {$\ket{\Pi_2}=\ket{\beta^{00}}$};
\end{tikzpicture}
\caption{Générateur d'états de Bell : $I\otimes H$, puis un CNOT de contrôle $B$ (fil du haut)
et de cible $A$ (fil du bas). L'entrée représentée est $\ket{00}$.}
\label{fig:circuit-bell}
\end{figure}

\subsubsection{Entrée \texorpdfstring{$\ket{00}$}{|00>} : calcul pas à pas}

\begin{align*}
  \ket{\Pi_0}&=\ket{0}_A\otimes\ket{0}_B=\ket{00},\\[1ex]
  \ket{\Pi_1}&=(I\otimes H)\ket{\Pi_0}=I\ket{0}_A\otimes H\ket{0}_B=\ket{0}\otimes\ket{+}
    =\frac{\ket{00}+\ket{01}}{\sqrt2},\\[1ex]
  \ket{\Pi_2}&=\mathrm{CNOT}\ket{\Pi_1}=\frac{\mathrm{CNOT}\ket{00}+\mathrm{CNOT}\ket{01}}{\sqrt2}
    =\frac{\ket{00}+\ket{11}}{\sqrt2}=\ket{\beta^{00}}=\ket{\Phi^+}.
\end{align*}
Dans $\ket{\Pi_1}$, le premier chiffre est celui de $A$ et le second celui de $B$. Le CNOT a
pour contrôle $B$ : $\ket{00}$ ($B=0$) ne change pas, $\ket{01}$ ($B=1$) devient $\ket{11}$
puisque $A$ est inversé. Le résultat est un \emph{état de Bell}.

Avec les matrices, la règle $A\otimes B=(A_{ij}B)$ donne
\[
  I\otimes H=\begin{pmatrix}1\cdot H&0\\0&1\cdot H\end{pmatrix}
  =\frac{1}{\sqrt2}\begin{pmatrix}1&1&0&0\\1&-1&0&0\\0&0&1&1\\0&0&1&-1\end{pmatrix}.
\]
Sa première colonne est $\ket{\Pi_1}=\frac{1}{\sqrt2}(1,1,0,0)^{\mathsf T}$ (image de $\ket{00}$).
Le CNOT échange les amplitudes de $\ket{01}$ et de $\ket{11}$ (section~\ref{sec:cnot}), donc
\[
  \ket{\Pi_2}=\mathrm{CNOT}\,\frac{1}{\sqrt2}\begin{pmatrix}1\\1\\0\\0\end{pmatrix}
  =\frac{1}{\sqrt2}\begin{pmatrix}1\\0\\0\\1\end{pmatrix}.
\]

\subsubsection{Les quatre entrées}

Prenons $B$ dans l'état $\ket{x}$ et $A$ dans l'état $\ket{y}$, avec $x,y\in\{0,1\}$. Alors
$\ket{\Pi_0}=\ket{y}_A\otimes\ket{x}_B=\ket{y,x}$. Comme
$H\ket{x}=\frac{1}{\sqrt2}\big(\ket{0}+(-1)^x\ket{1}\big)$ (c'est $\ket{+}$ pour $x=0$ et
$\ket{-}$ pour $x=1$),
\[
  \ket{\Pi_1}=\ket{y}\otimes H\ket{x}=\frac{1}{\sqrt2}\Big(\ket{y,0}+(-1)^x\ket{y,1}\Big).
\]
Le CNOT (contrôle $B$) ne change pas $\ket{y,0}$ et transforme $\ket{y,1}$ en
$\ket{\bar y,1}$, avec $\bar y=1-y$ :
\[
  \ket{\Pi_2}=\ket{\beta^{xy}}=\frac{1}{\sqrt2}\Big(\ket{y,0}+(-1)^x\ket{\bar y,1}\Big).
\]
Le premier indice de $\ket{\beta^{xy}}$ est l'entrée de $B$ (le qubit qui reçoit $H$), le
second celle de $A$ (la cible). Les quatre cas :

\begin{center}
\renewcommand{\arraystretch}{1.7}
\begin{tabular}{@{}ccccc@{}}
\toprule
$(B,A)=(x,y)$ & $\ket{\Pi_0}=\ket{y,x}$ & $\ket{\Pi_1}$ & $\ket{\Pi_2}=\ket{\beta^{xy}}$ & Nom \\
\midrule
$(0,0)$ & $\ket{00}$ & $\frac{1}{\sqrt2}(\ket{00}+\ket{01})$ & $\frac{1}{\sqrt2}(\ket{00}+\ket{11})$ & $\ket{\Phi^+}$ \\
$(1,0)$ & $\ket{01}$ & $\frac{1}{\sqrt2}(\ket{00}-\ket{01})$ & $\frac{1}{\sqrt2}(\ket{00}-\ket{11})$ & $\ket{\Phi^-}$ \\
$(0,1)$ & $\ket{10}$ & $\frac{1}{\sqrt2}(\ket{10}+\ket{11})$ & $\frac{1}{\sqrt2}(\ket{10}+\ket{01})$ & $\ket{\Psi^+}$ \\
$(1,1)$ & $\ket{11}$ & $\frac{1}{\sqrt2}(\ket{10}-\ket{11})$ & $\frac{1}{\sqrt2}(\ket{10}-\ket{01})$ & $-\ket{\Psi^-}$ \\
\bottomrule
\end{tabular}
\end{center}
Détail de la dernière ligne : $\ket{\Pi_1}=\ket{1}\otimes\ket{-}=\frac{1}{\sqrt2}(\ket{10}-\ket{11})$ ;
le CNOT laisse $\ket{10}$ et envoie $\ket{11}$ sur $\ket{01}$, d'où
$\frac{1}{\sqrt2}(\ket{10}-\ket{01})=-\frac{1}{\sqrt2}(\ket{01}-\ket{10})=-\ket{\Psi^-}$.

\paragraph{Forme matricielle.} $\mathrm{CNOT}_{21}$ échange les lignes $2$ et $4$ de
$I\otimes H$ (il échange les amplitudes de $\ket{01}$ et $\ket{11}$), donc
\[
  \mathrm{CNOT}_{21}\,(I\otimes H)=\frac{1}{\sqrt2}\begin{pmatrix}
    1 & 1 & 0 & 0\\ 0 & 0 & 1 & -1\\ 0 & 0 & 1 & 1\\ 1 & -1 & 0 & 0
  \end{pmatrix}.
\]
Ses colonnes sont les images de $\ket{00},\ket{01},\ket{10},\ket{11}$, c'est-à-dire de
$(B,A)=(0,0),(1,0),(0,1),(1,1)$ : ce sont $\ket{\beta^{00}}$, $\ket{\beta^{10}}$,
$\ket{\beta^{01}}$, $\ket{\beta^{11}}$, soit $\ket{\Phi^+}$, $\ket{\Phi^-}$, $\ket{\Psi^+}$
et $-\ket{\Psi^-}$.

\begin{remarque}[Comparaison avec la section précédente]
Le circuit de la figure~\ref{fig:bell} applique $H$ au \emph{premier} chiffre et prend celui-ci
pour contrôle ; ici $H$ et le contrôle portent sur le \emph{dernier} chiffre. Les deux
circuits se déduisent l'un de l'autre en échangeant les deux fils (conjugaison par le SWAP,
section~\ref{sec:swap}). Les quatre états de Bell obtenus sont les mêmes, avec les mêmes
indices, sauf le dernier : $-\ket{\Psi^-}$ ici, $+\ket{\Psi^-}$ là-bas. Ce signe est une phase
globale, sans effet sur les mesures (elle disparaît de $\rho$, section~\ref{sec:rho}).
\end{remarque}

\subsubsection{Matrice de densité de l'état de sortie}

Pour $\ket{\Psi}=\ket{\beta^{00}}=\frac{1}{\sqrt2}\ket{00}+\frac{1}{\sqrt2}\ket{11}$ :
\[
  \rho_{AB}=\ket{\Psi}\bra{\Psi}
  =\begin{pmatrix}\frac{1}{\sqrt2}\\0\\0\\\frac{1}{\sqrt2}\end{pmatrix}
   \begin{pmatrix}\frac{1}{\sqrt2}&0&0&\frac{1}{\sqrt2}\end{pmatrix}
  =\frac12\begin{pmatrix}1&0&0&1\\0&0&0&0\\0&0&0&0\\1&0&0&1\end{pmatrix}.
\]
Sous forme de ket-bra :
$\rho_{AB}=\frac12\big(\ket{00}\bra{00}+\ket{00}\bra{11}+\ket{11}\bra{00}+\ket{11}\bra{11}\big)$.
On vérifie $\operatorname{Tr}\rho_{AB}=\frac12(1+1)=1$ et
\[
  \rho_{AB}^2=\frac14\begin{pmatrix}2&0&0&2\\0&0&0&0\\0&0&0&0\\2&0&0&2\end{pmatrix}=\rho_{AB},
\]
donc $\operatorname{Tr}\rho_{AB}^2=1$ : le couple $(A,B)$ est dans un état \emph{pur}.

% ---------------------------------------------------------------------
\subsection{Information : variables aléatoires et états quantiques}\label{sec:info}
% ---------------------------------------------------------------------

\begin{definition}{Entropie, information mutuelle, entropie conditionnelle}
Soit $A$ une variable aléatoire de loi $p(a)$. Son \emph{entropie de Shannon}, en bits, est
\[
  H(A)=-\sum_a p(a)\log_2p(a)\qquad(\text{avec }0\log_20=0).
\]
Elle mesure l'incertitude moyenne sur $A$ : une pièce équilibrée a $H=1$ bit. Pour un couple
$(A,B)$ de loi $p(a,b)$, $H(A,B)=-\sum_{a,b}p(a,b)\log_2p(a,b)$, et
\begin{align*}
  I(A;B)&=H(A)+H(B)-H(A,B) &&\text{(information mutuelle)},\\
  H(A|B)&=H(A,B)-H(B)=H(A)-I(A;B) &&\text{(entropie conditionnelle)}.
\end{align*}
\end{definition}

Les deux écritures de $H(A|B)$ sont égales : la définition de $I(A;B)$ donne
$H(A,B)=H(A)+H(B)-I(A;B)$, donc $H(A,B)-H(B)=H(A)-I(A;B)$. L'information mutuelle $I(A;B)$
mesure ce que la connaissance de $B$ apprend sur $A$ ; $H(A|B)$ est l'incertitude qui reste sur
$A$ quand on connaît $B$.

\subsubsection{Deux variables aléatoires indépendantes}

Alice et Bob lancent chacun une pièce équilibrée, sans lien entre les deux lancers :
$p(a,b)=\frac14$ pour les quatre couples. Chaque pièce donne $0$ ou $1$ avec la probabilité $\frac12$, donc
\begin{align*}
  H(A)&=-\Big(\tfrac12\log_2\tfrac12+\tfrac12\log_2\tfrac12\Big)=1\text{ bit},\\
  H(B)&=1\text{ bit (même calcul)},\\
  H(A,B)&=-4\cdot\tfrac14\log_2\tfrac14=2\text{ bits},\\
  I(A;B)&=H(A)+H(B)-H(A,B)=1+1-2=0,\\
  H(A|B)&=H(A)-I(A;B)=1-0=H(A)=1\text{ bit}.
\end{align*}
Connaître $B$ n'apprend rien sur $A$.

\subsubsection{Deux variables solidaires}

Les deux pièces tombent toujours du même côté : $p(0,0)=p(1,1)=\frac12$ et
$p(0,1)=p(1,0)=0$. Les lois de chaque pièce sont inchangées ($p(a{=}0)=p(0,0)+p(0,1)=\frac12$,
etc.), donc $H(A)=H(B)=1$ bit, mais
\begin{align*}
  H(A,B)&=-2\cdot\tfrac12\log_2\tfrac12-2\cdot0\log_20=1\text{ bit},\\
  I(A;B)&=1+1-1=1\text{ bit},\qquad H(A|B)=1-1=0.
\end{align*}
Connaître $B$ détermine $A$ : plus aucune incertitude ne reste.

\subsubsection{Deux qubits dans un état de Bell}

Pour un système quantique, on remplace la loi de probabilité par la matrice de densité.

\begin{definition}{Trace partielle et entropie de von Neumann}
L'état du seul qubit $A$ d'un système $(A,B)$ de matrice de densité $\rho_{AB}$ est obtenu en
\emph{traçant sur $B$} : $\rho_A=\operatorname{Tr}_B(\rho_{AB})$. La trace partielle est
linéaire et remplace la partie $B$ de chaque terme par sa trace :
\[
  \operatorname{Tr}_B\big(\ket{a}\bra{a'}\otimes\ket{b}\bra{b'}\big)=\ket{a}\bra{a'}\;\braket{b'}{b}.
\]
L'\emph{entropie de von Neumann} d'un état $\rho$, de valeurs propres $\lambda_i$, est
$S(\rho)=-\operatorname{Tr}(\rho\log_2\rho)=-\sum_i\lambda_i\log_2\lambda_i$. En quantique,
$H(A)=S(\rho_A)$, $H(B)=S(\rho_B)$ et $H(A,B)=S(\rho_{AB})$ ; les définitions de
$I(A;B)$ et de $H(A|B)$ sont les mêmes.
\end{definition}

On part de l'état de Bell
$\rho_{AB}=\frac12\big(\ket{00}\bra{00}+\ket{00}\bra{11}+\ket{11}\bra{00}+\ket{11}\bra{11}\big)$
produit par le circuit de la section~\ref{sec:circuit-bell}. Chaque terme est un produit
d'opérateurs à un qubit :
$\ket{00}\bra{11}=\ket{0}\bra{1}\otimes\ket{0}\bra{1}$, etc. Alors
\begin{align*}
  \rho_A=\operatorname{Tr}_B(\rho_{AB})
  &=\tfrac12\Big(\ket{0}\bra{0}\,\braket{0}{0}+\ket{0}\bra{1}\,\braket{1}{0}
      +\ket{1}\bra{0}\,\braket{0}{1}+\ket{1}\bra{1}\,\braket{1}{1}\Big)\\
  &=\tfrac12\Big(\ket{0}\bra{0}\cdot1+\ket{0}\bra{1}\cdot0+\ket{1}\bra{0}\cdot0+\ket{1}\bra{1}\cdot1\Big)\\
  &=\tfrac12\big(\ket{0}\bra{0}+\ket{1}\bra{1}\big)=\tfrac12\,\mathrm{Id}
  =\begin{pmatrix}\frac12&0\\0&\frac12\end{pmatrix}.
\end{align*}
La formule des composantes de la section~\ref{sec:separable},
$(\rho_A)_{ij}=\rho_{(i0),(j0)}+\rho_{(i1),(j1)}$, donne le même résultat. En traçant sur $A$
au lieu de $B$, la règle est $\operatorname{Tr}_A(\ket{a}\bra{a'}\otimes\ket{b}\bra{b'})
=\braket{a'}{a}\,\ket{b}\bra{b'}$, et on trouve de même $\rho_B=\frac12\mathrm{Id}$.

Calcul des entropies :
\begin{itemize}[nosep]
  \item $\rho_{AB}$ est pur : ses valeurs propres sont $1,0,0,0$, donc
        $H(A,B)=S(\rho_{AB})=-1\cdot\log_21=0$ ;
  \item $\rho_A$ et $\rho_B$ ont pour valeurs propres $\frac12,\frac12$, donc
        $H(A)=H(B)=-\big(\frac12\log_2\frac12+\frac12\log_2\frac12\big)=1$ bit.
\end{itemize}
D'où
\begin{align*}
  I(A;B)&=H(A)+H(B)-H(A,B)=1+1-0=2\text{ bits},\\
  H(A|B)&=H(A,B)-H(B)=0-1=-1\text{ bit}\qquad(=H(A)-I(A;B)=1-2).
\end{align*}

\begin{remarque}[Une entropie conditionnelle négative]
Pour des variables classiques, $p(a,b)\leq p(b)$, donc $-\log_2p(a,b)\geq-\log_2p(b)$ et
$H(A,B)\geq H(B)$ : on a toujours $H(A|B)\geq0$, et $I(A;B)\leq H(A)$. Ici $H(A|B)=-1$ et
$I(A;B)=2>H(A)$ : ces valeurs n'ont pas d'équivalent classique, c'est une signature de
l'intrication. Le résultat est surprenant mais mathématiquement correct.
\end{remarque}

\begin{center}
\renewcommand{\arraystretch}{1.4}
\begin{tabular}{@{}lccc@{}}
\toprule
 & indépendantes & solidaires & état de Bell \\
\midrule
$H(A)$ & $1$ & $1$ & $1$ \\
$H(B)$ & $1$ & $1$ & $1$ \\
$H(A,B)$ & $2$ & $1$ & $0$ \\
$I(A;B)$ & $0$ & $1$ & $2$ \\
$H(A|B)$ & $1$ & $0$ & $-1$ \\
\bottomrule
\end{tabular}
\end{center}
(en bits ; les deux premières colonnes sont classiques, la troisième est quantique).

% ---------------------------------------------------------------------
\subsection{Système pur et système mélangé}\label{sec:pur-melange}
% ---------------------------------------------------------------------

Rappel (section~\ref{sec:pur-mixte}) : si $\lambda_i$ sont les valeurs propres de la matrice de
densité $\rho$, alors $\operatorname{Tr}\rho^2=\sum_i\lambda_i^2\leq1$. Le système est
\emph{pur} si $\operatorname{Tr}\rho^2=1$ (alors $\rho=\ket{\psi}\bra{\psi}$ et $\rho^2=\rho$),
et \emph{mélangé} si $\operatorname{Tr}\rho^2<1$. Pour un qubit,
$\operatorname{Tr}\rho^2=\frac{1+|\vec r|^2}{2}$ (section~\ref{sec:bloch}) est compris entre
$\frac12$ et $1$.

\subsubsection{Un système pur : exercice}

\emph{Énoncé.} Soit $\ket{\psi}=\frac{1}{\sqrt2}\ket{0}+\frac{1}{\sqrt2}\ket{1}$.
\begin{enumerate}[label=\alph*), nosep]
  \item Calculer la matrice de densité $\rho$.
  \item Calculer la pureté du système quantique.
  \item Calculer $\langle\sigma_x\rangle$ et $\langle\sigma_z\rangle$.
  \item Écrire $\ket{\psi}$ dans la base $X$.
\end{enumerate}

\paragraph{a) Matrice de densité.} $\rho=\ket{\psi}\bra{\psi}$ est le produit d'un vecteur
colonne par un vecteur ligne (les amplitudes sont réelles) :
\[
  \rho=\begin{pmatrix}\frac{1}{\sqrt2}\\[0.4ex]\frac{1}{\sqrt2}\end{pmatrix}
       \begin{pmatrix}\frac{1}{\sqrt2}&\frac{1}{\sqrt2}\end{pmatrix}
  =\begin{pmatrix}\frac12&\frac12\\[0.4ex]\frac12&\frac12\end{pmatrix}.
\]

\paragraph{b) Pureté.}
\[
  \rho^2=\begin{pmatrix}\frac12&\frac12\\[0.4ex]\frac12&\frac12\end{pmatrix}
         \begin{pmatrix}\frac12&\frac12\\[0.4ex]\frac12&\frac12\end{pmatrix}
  =\begin{pmatrix}\frac14+\frac14&\frac14+\frac14\\[0.4ex]\frac14+\frac14&\frac14+\frac14\end{pmatrix}
  =\begin{pmatrix}\frac12&\frac12\\[0.4ex]\frac12&\frac12\end{pmatrix}=\rho,
\]
donc $\operatorname{Tr}\rho^2=\operatorname{Tr}\rho=\frac12+\frac12=1$ : le système est \emph{pur}.

\paragraph{c) Valeurs moyennes.} Avec $\sigma_x=X=\begin{pmatrix}0&1\\1&0\end{pmatrix}$ et
$\sigma_z=Z=\begin{pmatrix}1&0\\0&-1\end{pmatrix}$ (section~\ref{sec:op-dirac}), la valeur
moyenne est $\langle\sigma\rangle=\operatorname{Tr}(\rho\,\sigma)$ (section~\ref{sec:rho}).
\begin{align*}
  \rho\,\sigma_x&=\begin{pmatrix}\frac12&\frac12\\[0.4ex]\frac12&\frac12\end{pmatrix}
     \begin{pmatrix}0&1\\1&0\end{pmatrix}
     =\begin{pmatrix}\frac12&\frac12\\[0.4ex]\frac12&\frac12\end{pmatrix},
     &\langle\sigma_x\rangle&=\operatorname{Tr}(\rho\,\sigma_x)=\tfrac12+\tfrac12=1,\\[1.5ex]
  \rho\,\sigma_z&=\begin{pmatrix}\frac12&\frac12\\[0.4ex]\frac12&\frac12\end{pmatrix}
     \begin{pmatrix}1&0\\0&-1\end{pmatrix}
     =\begin{pmatrix}\frac12&-\frac12\\[0.4ex]\frac12&-\frac12\end{pmatrix},
     &\langle\sigma_z\rangle&=\operatorname{Tr}(\rho\,\sigma_z)=\tfrac12-\tfrac12=0.
\end{align*}
On retrouve $\langle\sigma_z\rangle$ avec les probabilités de mesure : dans la base $Z$,
$\sigma_z$ vaut $+1$ pour $\ket{0}$ et $-1$ pour $\ket{1}$, et
$|\braket{0}{\psi}|^2=|\braket{1}{\psi}|^2=\frac12$, donc
$\langle\sigma_z\rangle=\frac12\times(+1)+\frac12\times(-1)=0$.

\paragraph{d) Écriture dans la base $X$.} La base $X$ est $\{\ket{+},\ket{-}\}$, orthonormée, donc
$\ket{+}\bra{+}+\ket{-}\bra{-}=\mathrm{Id}$ et
$\ket{\psi}=\ket{+}\braket{+}{\psi}+\ket{-}\braket{-}{\psi}$, avec
\[
  \braket{+}{\psi}=\tfrac{1}{\sqrt2}\Big(\tfrac{1}{\sqrt2}+\tfrac{1}{\sqrt2}\Big)=1,
  \qquad
  \braket{-}{\psi}=\tfrac{1}{\sqrt2}\Big(\tfrac{1}{\sqrt2}-\tfrac{1}{\sqrt2}\Big)=0 .
\]
Donc $\ket{\psi}=1\cdot\ket{+}+0\cdot\ket{-}=\ket{+}$ : dans la base $X$, l'état n'a qu'une
seule composante. Comme $\sigma_x\ket{+}=+\ket{+}$ et $\sigma_x\ket{-}=-\ket{-}$, une
mesure de $\sigma_x$ donne $+1$ avec certitude, et
$\langle\sigma_x\rangle=1\cdot(+1)+0\cdot(-1)=1$ : c'est le résultat de la question c).
Le vecteur de Bloch de $\ket{+}$ est $(\langle\sigma_x\rangle,\langle\sigma_y\rangle,\langle\sigma_z\rangle)=(1,0,0)$,
de norme $1$ : l'état est sur la sphère, en accord avec $\operatorname{Tr}\rho^2=\frac{1+1}{2}=1$.

\subsubsection{Un système mélangé : un qubit de la paire de Bell}

À la section~\ref{sec:info}, la trace sur $B$ de l'état de Bell a donné
$\rho_A=\frac12\mathrm{Id}=\begin{pmatrix}\frac12&0\\0&\frac12\end{pmatrix}$. Alors
\[
  \rho_A^2=\begin{pmatrix}\frac14&0\\0&\frac14\end{pmatrix},
  \qquad
  \operatorname{Tr}\big(\rho_A^2\big)=\frac14+\frac14=\frac12<1 :
\]
le qubit $A$ est dans un état \emph{mélangé}, et c'est le plus mélangé possible pour un qubit
(centre de la boule de Bloch, $\vec r=\vec0$). Pourtant le couple $(A,B)$ est pur
($\operatorname{Tr}\rho_{AB}^2=1$) : on connaît parfaitement l'état du couple, mais pas celui d'un
qubit pris seul. L'information est entièrement dans les corrélations entre $A$ et $B$, et
$S(\rho_A)=1$ bit.

\subsubsection{Un système mélangé : le mélange statistique de Bob}

Bob prépare un qubit dans l'état $\ket{0}$ avec la probabilité $p_1=\frac34$ (75\,\%) et dans
l'état $\ket{1}$ avec la probabilité $p_2=\frac14$ (25\,\%), sans que l'on sache lequel des deux a
été préparé. Ce n'est pas un état pur $\ket{\psi}$ : on décrit le système par la matrice de
densité $\rho=\sum_ip_i\ket{\psi_i}\bra{\psi_i}$ (section~\ref{sec:pur-mixte}) :
\[
  \rho=\underbrace{\tfrac34}_{p_1}\ket{0}\bra{0}+\underbrace{\tfrac14}_{p_2}\ket{1}\bra{1}
  =\begin{pmatrix}\frac34&0\\0&\frac14\end{pmatrix}.
\]
Alors
\[
  \rho^2=\begin{pmatrix}\frac{9}{16}&0\\0&\frac{1}{16}\end{pmatrix},
  \qquad
  \operatorname{Tr}\big(\rho^2\big)=\frac{9}{16}+\frac{1}{16}=\frac{10}{16}=\frac58<1 :
\]
le système est \emph{mélangé}.

Les valeurs moyennes sont les moyennes pondérées sur la préparation :
$\langle\sigma_z\rangle=\operatorname{Tr}(\rho\,\sigma_z)=\frac34-\frac14=\frac12
=p_1\cdot(+1)+p_2\cdot(-1)$, et $\langle\sigma_x\rangle=\operatorname{Tr}(\rho\,\sigma_x)=0$
(les termes hors diagonale de $\rho$ sont nuls). Le vecteur de Bloch est $(0,0,\frac12)$,
à l'intérieur de la boule, et $\operatorname{Tr}\rho^2=\frac{1+1/4}{2}=\frac58$. Son entropie
est $S(\rho)=-\frac34\log_2\frac34-\frac14\log_2\frac14=2-\frac34\log_23\approx0{,}811$ bit.

\begin{remarque}[Mélange et superposition]
Considérons l'état pur $\ket{\psi'}=\frac{\sqrt3}{2}\ket{0}+\frac12\ket{1}$, qui a les mêmes
probabilités en base $Z$ ($P(0)=\frac34$, $P(1)=\frac14$). Sa matrice de densité est
\[
  \rho'=\begin{pmatrix}\frac34&\frac{\sqrt3}{4}\\[0.4ex]\frac{\sqrt3}{4}&\frac14\end{pmatrix},
  \qquad
  \operatorname{Tr}\rho'^2=\frac{9}{16}+\frac1{16}+2\cdot\frac{3}{16}=1,
  \qquad
  \langle\sigma_x\rangle=2\cdot\frac{\sqrt3}{4}=\frac{\sqrt3}{2}.
\]
Le mélange de Bob a la même diagonale mais aucun terme hors diagonale ($\operatorname{Tr}\rho^2=\frac58$,
$\langle\sigma_x\rangle=0$) : une superposition n'est pas un tirage au hasard entre $\ket{0}$ et
$\ket{1}$. De plus, pour un état pur il existe une mesure au résultat certain (ici $\sigma_x$ pour
$\ket{+}$) ; pour le mélange de Bob, la meilleure probabilité d'obtenir un même résultat est
$75\,\%$ (mesure en base $Z$, résultat $\ket{0}$).
\end{remarque}

\subsubsection{Bilan}

\begin{center}
\renewcommand{\arraystretch}{1.5}
\begin{tabular}{@{}lcccl@{}}
\toprule
\textbf{Système} & $\rho$ & $\operatorname{Tr}\rho^2$ & $S(\rho)$ & \textbf{Nature} \\
\midrule
$\ket{\psi}=\ket{+}$ & ${\renewcommand{\arraystretch}{1}\begin{pmatrix}\frac12&\frac12\\\frac12&\frac12\end{pmatrix}}$ & $1$ & $0$ & pur \\[2ex]
Qubit $A$ d'une paire de Bell & ${\renewcommand{\arraystretch}{1}\begin{pmatrix}\frac12&0\\0&\frac12\end{pmatrix}}$ & $\frac12$ & $1$ bit & mélangé \\[2ex]
Source de Bob (75\,\%/25\,\%) & ${\renewcommand{\arraystretch}{1}\begin{pmatrix}\frac34&0\\0&\frac14\end{pmatrix}}$ & $\frac58$ & $\approx0{,}811$ bit & mélangé \\
\bottomrule
\end{tabular}
\end{center}

% ---------------------------------------------------------------------
\subsection{Algorithme quantique : la téléportation}\label{sec:teleportation}
% ---------------------------------------------------------------------

\emph{Problème.} Alice détient un qubit $Q$ dans un état qu'elle ne connaît pas,
$\ket{\varphi}=\alpha\ket{0}+\beta\ket{1}$ avec $|\alpha|^2+|\beta|^2=1$, et elle veut que Bob
obtienne cet état sur l'un de ses qubits. Elle ne peut pas mesurer $Q$ pour connaître
$\alpha$ et $\beta$ (une mesure ne donne qu'un bit et détruit l'état), elle ne peut pas le
copier (théorème de non-clonage) et elle n'envoie pas le qubit lui-même.

\begin{definition}{Protocole de téléportation quantique}
\textbf{Ressources.}
\begin{itemize}[nosep]
  \item Un \emph{e-bit} : Alice et Bob partagent une paire de qubits $(A,B)$ dans l'état de Bell
        $\ket{\Phi^+}_{BA}=\frac{1}{\sqrt2}\big(\ket{00}+\ket{11}\big)$ ; Alice détient $A$,
        Bob détient $B$.
  \item Le qubit $Q$ à téléporter, chez Alice.
  \item Un canal classique (deux bits).
\end{itemize}
Le système est écrit $(B,A,Q)$ : $\ket{b,a,q}=\ket{b}_B\otimes\ket{a}_A\otimes\ket{q}_Q$.

\textbf{Étapes.}
\begin{enumerate}[nosep]
  \item Alice applique un CNOT de contrôle $Q$ et de cible $A$.
  \item Alice applique une porte $H$ sur $Q$.
  \item Alice mesure $A$ et $Q$ dans la base de calcul, obtient deux bits $(a,q)$ et les envoie
        à Bob.
  \item Bob applique $X$ si $a=1$, puis $Z$ si $q=1$, à son qubit $B$.
\end{enumerate}
\end{definition}

\begin{figure}[!ht]
\centering
\begin{adjustbox}{max width=\linewidth}
\begin{tikzpicture}[>=Stealth, font=\small]
  % fils quantiques : Q en haut, A au milieu, B en bas
  \draw (0,2.4) -- (5.3,2.4);
  \draw (0,1.2) -- (5.3,1.2);
  \draw (0,0) -- (10.6,0);
  \node[anchor=east] at (0,2.4) {$\ket{\varphi}_Q$};
  \node[anchor=south west, text=insagris, font=\footnotesize] at (0.05,1.2) {$A$};
  \node[anchor=south west, text=insagris, font=\footnotesize] at (0.05,0) {$B$};
  % paire de Bell partagée
  \draw[insagris, thick] (-0.1,1.5) -- (-0.25,1.5) -- (-0.25,-0.3) -- (-0.1,-0.3);
  \node[anchor=east] at (-0.35,0.6) {$\ket{\Phi^+}_{BA}$};
  % CNOT : contrôle Q, cible A
  \draw[insarouge, thick] (2.6,2.4) -- (2.6,1.2);
  \pic at (2.6,2.4) {ctl};
  \pic at (2.6,1.2) {cible};
  % Hadamard sur Q
  \node[porte] at (4.2,2.4) {$H$};
  % mesures d'Alice
  \pic at (5.6,2.4) {mesure};
  \pic at (5.6,1.2) {mesure};
  % fils classiques (doubles)
  \draw[double, double distance=1.2pt, insagris] (5.9,1.2) -- (7.6,1.2) -- (7.6,0.35);
  \draw[double, double distance=1.2pt, insagris] (5.9,2.4) -- (9.6,2.4) -- (9.6,0.35);
  \node[above, font=\footnotesize] at (6.5,1.2) {$a$};
  \node[above, font=\footnotesize] at (6.5,2.4) {$q$};
  % corrections de Bob
  \node[porte] at (7.6,0) {$X$};
  \node[porte] at (9.6,0) {$Z$};
  \node[anchor=west] at (10.6,0) {$\ket{\varphi}_B$};
  % états intermédiaires
  \foreach \x in {1.4,3.4,4.9}
    \draw[dashed, insagris!70] (\x,2.8) -- (\x,-0.5);
  \node[anchor=north, font=\footnotesize] at (1.4,-0.5) {$\ket{\Pi_0}$};
  \node[anchor=north, font=\footnotesize] at (3.4,-0.5) {$\ket{\Pi_1}$};
  \node[anchor=north, font=\footnotesize] at (4.9,-0.5) {$\ket{\Pi_2}$};
\end{tikzpicture}
\end{adjustbox}
\caption{Circuit de téléportation. Alice détient $Q$ (dans $\ket{\varphi}=\alpha\ket{0}+\beta\ket{1}$)
et $A$, Bob détient $B$ ; $A$ et $B$ forment la paire $\ket{\Phi^+}_{BA}$. Les traits doubles
sont des fils \emph{classiques} qui portent les résultats $a$ (mesure de $A$) et $q$ (mesure de
$Q$) ; $X$ est appliqué si $a=1$, $Z$ si $q=1$. Le ket s'écrit $\ket{b,a,q}$ : le fil du haut
($Q$) est le dernier chiffre.}
\label{fig:teleportation}
\end{figure}

\subsubsection{Calcul pas à pas}

\paragraph{État initial $\ket{\Pi_0}$.} Alice et Bob partagent la paire, et $Q$ est indépendant
d'elle :
\begin{align*}
  \ket{\Pi_0}&=\ket{\Phi^+}_{BA}\otimes\ket{\varphi}_Q
    =\frac{1}{\sqrt2}\big(\ket{00}+\ket{11}\big)\otimes\big(\alpha\ket{0}+\beta\ket{1}\big)\\
  &=\frac{1}{\sqrt2}\Big(\alpha\ket{00}\ket{0}+\beta\ket{00}\ket{1}+\alpha\ket{11}\ket{0}+\beta\ket{11}\ket{1}\Big)\\
  &=\frac{\alpha\ket{000}+\beta\ket{001}+\alpha\ket{110}+\beta\ket{111}}{\sqrt2}.
\end{align*}

\paragraph{Après le CNOT $Q\to A$ : $\ket{\Pi_1}$.} Le CNOT (contrôle $Q$, dernier chiffre ;
cible $A$, chiffre du milieu) agit ainsi : $\ket{b,a,q}\to\ket{b,\,a\oplus q,\,q}$. Terme par terme :
\begin{center}
\renewcommand{\arraystretch}{1.3}
\begin{tabular}{@{}cccc@{}}
\toprule
Terme & $q$ & Effet sur $A$ & Résultat \\
\midrule
$\alpha\ket{000}$ & $0$ & aucun & $\alpha\ket{000}$ \\
$\beta\ket{001}$ & $1$ & $0\to1$ & $\beta\ket{011}$ \\
$\alpha\ket{110}$ & $0$ & aucun & $\alpha\ket{110}$ \\
$\beta\ket{111}$ & $1$ & $1\to0$ & $\beta\ket{101}$ \\
\bottomrule
\end{tabular}
\end{center}
donc
\[
  \ket{\Pi_1}=\frac{\alpha\ket{000}+\beta\ket{011}+\alpha\ket{110}+\beta\ket{101}}{\sqrt2}.
\]
On isole ensuite le dernier chiffre, celui de $Q$, qui va changer à l'étape suivante
(le premier ket est celui du couple $\ket{b,a}$) :
\[
  \ket{\Pi_1}=\frac{\alpha\ket{00}\otimes\ket{0}+\beta\ket{01}\otimes\ket{1}
    +\alpha\ket{11}\otimes\ket{0}+\beta\ket{10}\otimes\ket{1}}{\sqrt2}.
\]

\paragraph{Après la porte $H$ sur $Q$ : $\ket{\Pi_2}$.} On remplace $\ket{0}$ par
$H\ket{0}=\ket{+}$ et $\ket{1}$ par $H\ket{1}=\ket{-}$ dans le facteur $Q$ :
\[
  \ket{\Pi_2}=\frac{1}{\sqrt2}\Big(\alpha\ket{00}\otimes\ket{+}+\beta\ket{01}\otimes\ket{-}
    +\alpha\ket{11}\otimes\ket{+}+\beta\ket{10}\otimes\ket{-}\Big).
\]
Avec $\ket{\pm}=\frac{1}{\sqrt2}(\ket{0}\pm\ket{1})$, le préfacteur devient
$\frac{1}{\sqrt2}\cdot\frac{1}{\sqrt2}=\frac12$ :
\begin{align*}
  \ket{\Pi_2}&=\frac12\Big(\alpha\ket{00}\big(\ket{0}+\ket{1}\big)+\beta\ket{01}\big(\ket{0}-\ket{1}\big)
     +\alpha\ket{11}\big(\ket{0}+\ket{1}\big)+\beta\ket{10}\big(\ket{0}-\ket{1}\big)\Big)\\
  &=\frac12\Big(\alpha\ket{000}+\alpha\ket{001}+\beta\ket{010}-\beta\ket{011}
     +\alpha\ket{110}+\alpha\ket{111}+\beta\ket{100}-\beta\ket{101}\Big).
\end{align*}

\paragraph{Regroupement selon le résultat d'Alice.} Alice mesurera $A$ et $Q$, donc les deux
derniers chiffres. On regroupe les huit termes qui ont les mêmes chiffres $(a,q)$, en mettant
le ket de Bob en premier :
\begin{center}
\renewcommand{\arraystretch}{1.3}
\begin{tabular}{@{}cll@{}}
\toprule
$(a,q)$ & Termes de $\ket{\Pi_2}$ (facteur $\frac12$ omis) & Ket de Bob \\
\midrule
$(0,0)$ & $\alpha\ket{000}+\beta\ket{100}$ & $(\alpha\ket{0}+\beta\ket{1})\ket{00}$ \\
$(0,1)$ & $\alpha\ket{001}-\beta\ket{101}$ & $(\alpha\ket{0}-\beta\ket{1})\ket{01}$ \\
$(1,0)$ & $\beta\ket{010}+\alpha\ket{110}$ & $(\beta\ket{0}+\alpha\ket{1})\ket{10}$ \\
$(1,1)$ & $-\beta\ket{011}+\alpha\ket{111}$ & $(-\beta\ket{0}+\alpha\ket{1})\ket{11}$ \\
\bottomrule
\end{tabular}
\end{center}
d'où
\[
  \ket{\Pi_2}=\frac12\Big[\big(\alpha\ket{0}+\beta\ket{1}\big)\ket{00}
    +\big(\alpha\ket{0}-\beta\ket{1}\big)\ket{01}
    +\big(\beta\ket{0}+\alpha\ket{1}\big)\ket{10}
    +\big(\alpha\ket{1}-\beta\ket{0}\big)\ket{11}\Big].
\]
Dans chaque terme, le ket de gauche est l'état de $B$ et le ket de droite $\ket{a,q}$ celui du
couple $(A,Q)$ d'Alice.

\subsubsection{Mesure d'Alice et corrections de Bob}

Alice mesure $A$ et $Q$. Comme à la section~\ref{sec:mesure2}, la probabilité d'un résultat
$(a,q)$ est le carré de la norme du terme correspondant, et l'état de Bob devient le ket
associé (déjà normé). Ici
\[
  P(a,q)=\Big\|\tfrac12\ket{\phi_{aq}}\Big\|^2=\tfrac14\big(|\alpha|^2+|\beta|^2\big)=\tfrac14
  \quad\text{pour les quatre résultats,}
\]
qui sont donc équiprobables et \emph{indépendants} de $\alpha$ et $\beta$ : les deux bits
envoyés ne révèlent rien sur $\ket{\varphi}$. Les kets de Bob sont
$\ket{\phi_{00}}=\alpha\ket{0}+\beta\ket{1}$, $\ket{\phi_{01}}=\alpha\ket{0}-\beta\ket{1}$,
$\ket{\phi_{10}}=\beta\ket{0}+\alpha\ket{1}$ et $\ket{\phi_{11}}=-\beta\ket{0}+\alpha\ket{1}$.

\begin{center}
\renewcommand{\arraystretch}{1.5}
\begin{tabular}{@{}cccll@{}}
\toprule
Résultat $(a,q)$ & Probabilité & État de Bob & Lien avec $\ket{\varphi}$ & Correction de Bob \\
\midrule
$(0,0)$ & $\frac14$ & $\alpha\ket{0}+\beta\ket{1}$ & $\ket{\varphi}$ & $I$ (rien) \\
$(0,1)$ & $\frac14$ & $\alpha\ket{0}-\beta\ket{1}$ & $Z\ket{\varphi}$ & $Z$ \\
$(1,0)$ & $\frac14$ & $\beta\ket{0}+\alpha\ket{1}$ & $X\ket{\varphi}$ & $X$ \\
$(1,1)$ & $\frac14$ & $-\beta\ket{0}+\alpha\ket{1}$ & $XZ\ket{\varphi}$ & $X$ puis $Z$ (opérateur $ZX$) \\
\bottomrule
\end{tabular}
\end{center}

\paragraph{Vérification des quatre corrections.} Avec $X\ket{0}=\ket{1}$, $X\ket{1}=\ket{0}$,
$Z\ket{0}=\ket{0}$ et $Z\ket{1}=-\ket{1}$ :
\begin{align*}
  (0,0):\quad & I\big(\alpha\ket{0}+\beta\ket{1}\big)=\alpha\ket{0}+\beta\ket{1},\\
  (0,1):\quad & Z\big(\alpha\ket{0}-\beta\ket{1}\big)=\alpha\ket{0}-\beta Z\ket{1}=\alpha\ket{0}+\beta\ket{1},\\
  (1,0):\quad & X\big(\beta\ket{0}+\alpha\ket{1}\big)=\beta\ket{1}+\alpha\ket{0}=\alpha\ket{0}+\beta\ket{1},\\
  (1,1):\quad & X\big(-\beta\ket{0}+\alpha\ket{1}\big)=-\beta\ket{1}+\alpha\ket{0}=\alpha\ket{0}-\beta\ket{1},
     \ \text{puis}\ Z\big(\alpha\ket{0}-\beta\ket{1}\big)=\alpha\ket{0}+\beta\ket{1}.
\end{align*}
Dans le dernier cas, avec les matrices,
$ZX=\begin{pmatrix}1&0\\0&-1\end{pmatrix}\begin{pmatrix}0&1\\1&0\end{pmatrix}
=\begin{pmatrix}0&1\\-1&0\end{pmatrix}$ et
$ZX\begin{pmatrix}-\beta\\ \alpha\end{pmatrix}=\begin{pmatrix}\alpha\\ \beta\end{pmatrix}$.
L'ordre compte : l'état de Bob est $XZ\ket{\varphi}$ ($Z\ket{\varphi}=\alpha\ket{0}-\beta\ket{1}$,
puis $X$ donne $\alpha\ket{1}-\beta\ket{0}$), et on l'annule avec
$(XZ)^{-1}=Z^{-1}X^{-1}=ZX$ : on applique d'abord $X$, puis $Z$, comme dans le circuit.
Dans les quatre cas, Bob retrouve exactement $\ket{\varphi}=\alpha\ket{0}+\beta\ket{1}$.

\subsubsection{Ce que Bob sait avant de recevoir les bits}

Tant que Bob ignore $(a,q)$, son qubit est dans le mélange des quatre kets $\ket{\phi_{aq}}$, chacun
avec la probabilité $\frac14$ :
\begin{align*}
  \rho_B&=\frac14\sum_{a,q}\ket{\phi_{aq}}\bra{\phi_{aq}}\\
  &=\frac14\Bigg[\begin{pmatrix}|\alpha|^2&\alpha\beta^*\\ \alpha^*\beta&|\beta|^2\end{pmatrix}
    +\begin{pmatrix}|\alpha|^2&-\alpha\beta^*\\ -\alpha^*\beta&|\beta|^2\end{pmatrix}
    +\begin{pmatrix}|\beta|^2&\alpha^*\beta\\ \alpha\beta^*&|\alpha|^2\end{pmatrix}
    +\begin{pmatrix}|\beta|^2&-\alpha^*\beta\\ -\alpha\beta^*&|\alpha|^2\end{pmatrix}\Bigg]\\
  &=\frac14\begin{pmatrix}2(|\alpha|^2+|\beta|^2)&0\\0&2(|\alpha|^2+|\beta|^2)\end{pmatrix}
   =\frac12\begin{pmatrix}1&0\\0&1\end{pmatrix}=\frac12\,\mathrm{Id}.
\end{align*}
Sans les deux bits, Bob a donc un qubit complètement mélangé, qui ne dépend pas de $\alpha$ ni de
$\beta$ : il n'a reçu aucune information sur $\ket{\varphi}$. C'est l'envoi classique de $(a,q)$
qui permet de la retrouver, donc rien ne voyage plus vite que la lumière.

\begin{remarque}[Bilan de la téléportation]
\begin{itemize}[nosep]
  \item Ressources : $1$ e-bit (la paire de Bell) et $2$ bits classiques pour transmettre
        \emph{un qubit} d'état quelconque, sans qu'Alice connaisse $\alpha$ ni $\beta$.
  \item Pas de clonage : après la mesure, $A$ et $Q$ sont dans l'état de base $\ket{a,q}$.
        Alice n'a plus $\ket{\varphi}$ : l'état a été \emph{déplacé}, pas copié. L'intrication
        entre $A$ et $B$ a été consommée.
  \item Les quatre corrections $I,Z,X,ZX$ sont des portes contrôlées par des bits classiques
        (section~\ref{sec:controlees}).
\end{itemize}
\end{remarque}

% ---------------------------------------------------------------------
\subsection{Portes contrôlées : contrôlé-\texorpdfstring{$U$}{U}, Fredkin et Toffoli}\label{sec:controlees}
% ---------------------------------------------------------------------

Le CNOT est le cas le plus simple d'une idée générale : appliquer une opération à un ou
plusieurs qubits \emph{cibles} seulement si un ou plusieurs qubits de \emph{contrôle} valent $1$.
Dans cette section, comme à la section~\ref{sec:op2}, le contrôle est le \emph{premier} chiffre
du ket. Les qubits sont notés $x$, $y$, $z$ (le qubit $x$ est le contrôle) et
$\ket{x,y,z}=\ket{x}\otimes\ket{y}\otimes\ket{z}$ ; d'après la convention de la section, le
contrôle est dessiné en bas (figure~\ref{fig:controlees}). Pour éviter la confusion avec les
matrices de Pauli, on écrit ici $\sigma_x=X$ pour la porte NOT.

\begin{figure}[!ht]
\centering
\begin{adjustbox}{max width=\linewidth}
\begin{tikzpicture}[>=Stealth, font=\small]
  % (a) contrôlé-U : contrôle x en bas, cible y en haut
  \draw (0,1.2) -- (3.4,1.2);
  \draw (0,0) -- (3.4,0);
  \node[anchor=east] at (0,1.2) {$y$};
  \node[anchor=east] at (0,0) {$x$};
  \draw[insarouge, thick] (1.7,0) -- (1.7,0.85);
  \pic at (1.7,0) {ctl};
  \node[porte] at (1.7,1.2) {$U$};
  \node[anchor=north, font=\footnotesize] at (1.7,-0.5) {(a) contrôlé-$U$};
  % (b) contrôlé-SWAP (Fredkin) : x contrôle, y et z échangés
  \draw (5,2.4) -- (8.4,2.4);
  \draw (5,1.2) -- (8.4,1.2);
  \draw (5,0) -- (8.4,0);
  \node[anchor=east] at (5,2.4) {$z$};
  \node[anchor=east] at (5,1.2) {$y$};
  \node[anchor=east] at (5,0) {$x$};
  \draw[insarouge, thick] (6.7,0) -- (6.7,2.4);
  \pic at (6.7,0) {ctl};
  \pic at (6.7,1.2) {croix};
  \pic at (6.7,2.4) {croix};
  \node[anchor=north, font=\footnotesize] at (6.7,-0.5) {(b) Fredkin (C-SWAP)};
  % (c) Toffoli (CCNOT) : x et y contrôlent, z cible
  \draw (10,2.4) -- (13.4,2.4);
  \draw (10,1.2) -- (13.4,1.2);
  \draw (10,0) -- (13.4,0);
  \node[anchor=east] at (10,2.4) {$z$};
  \node[anchor=east] at (10,1.2) {$y$};
  \node[anchor=east] at (10,0) {$x$};
  \draw[insarouge, thick] (11.7,0) -- (11.7,2.4);
  \pic at (11.7,0) {ctl};
  \pic at (11.7,1.2) {ctl};
  \pic at (11.7,2.4) {cible};
  \node[anchor=north, font=\footnotesize] at (11.7,-0.5) {(c) Toffoli (CCNOT)};
\end{tikzpicture}
\end{adjustbox}
\caption{Portes contrôlées. Le contrôle $x$ est le premier chiffre du ket, donc le fil du bas.
(a) Le qubit $y$ subit $U$ si $x=1$ ; pour $U=X$ on retrouve le CNOT. (b) Si $x=1$, les qubits
$y$ et $z$ sont échangés. (c) Si $x=y=1$, le qubit $z$ est inversé.}
\label{fig:controlees}
\end{figure}

\subsubsection{Contrôlé-\texorpdfstring{$U$}{U}}

\begin{definition}{Porte contrôlée-$U$}
Soit $U$ une matrice unitaire $2\times2$. Sur le système $(x,y)$, de contrôle $x$ et de cible $y$,
\[
  \mathrm{C}U=\ket{0}\bra{0}_x\otimes I_y+\ket{1}\bra{1}_x\otimes U_y .
\]
Le premier terme agit quand le contrôle est $\ket{0}$ : la cible ne change pas ($I$). Le second
agit quand le contrôle est $\ket{1}$ : la cible subit $U$.
\end{definition}

Comme $\ket{0}\bra{0}\ket{x}=\delta_{x0}\ket{0}$ et $\ket{1}\bra{1}\ket{x}=\delta_{x1}\ket{1}$,
\[
  \mathrm{C}U\big(\ket{x}\otimes\ket{y}\big)=\ket{x}\otimes U^x\ket{y},\qquad U^0=I,\ U^1=U .
\]
Par la règle $A\otimes B=(A_{ij}B)$, $\ket{0}\bra{0}\otimes I=\begin{pmatrix}I&0\\0&0\end{pmatrix}$
et $\ket{1}\bra{1}\otimes U=\begin{pmatrix}0&0\\0&U\end{pmatrix}$ (blocs $2\times2$), donc
\[
  \mathrm{C}U=\begin{pmatrix}I&0\\0&U\end{pmatrix}
  =\begin{pmatrix}1&0&0&0\\0&1&0&0\\0&0&u_{00}&u_{01}\\0&0&u_{10}&u_{11}\end{pmatrix}
  \qquad\text{pour }U=\begin{pmatrix}u_{00}&u_{01}\\u_{10}&u_{11}\end{pmatrix}.
\]
Cette matrice est unitaire car $I^\dagger I=I$ et $U^\dagger U=I$.

\paragraph{Le CNOT est un cas particulier.} Pour $U=\sigma_x$,
$\mathrm{C}\sigma_x=\ket{0}\bra{0}_x\otimes I_y+\ket{1}\bra{1}_x\otimes\sigma_x$ est le CNOT de
contrôle $x$ (premier chiffre), c'est-à-dire $\mathrm{CNOT}_{12}$ de la section~\ref{sec:cnot}.
On calcule l'action sur les quatre états de base, avec $\braket{0}{0}=\braket{1}{1}=1$,
$\braket{0}{1}=\braket{1}{0}=0$, $\sigma_x\ket{0}=\ket{1}$ et $\sigma_x\ket{1}=\ket{0}$ :
\begin{align*}
  \mathrm{C}\sigma_x\ket{00}&=\big[\ket{0}\bra{0}_x\otimes I_y\big]\ket{0}_x\ket{0}_y
      +\big[\ket{1}\bra{1}_x\otimes\sigma_x\big]\ket{0}_x\ket{0}_y\\
    &=\ket{0}\braket{0}{0}\otimes I\ket{0}+\ket{1}\braket{1}{0}\otimes\sigma_x\ket{0}
     =\ket{0}\otimes\ket{0}+0=\ket{00},\\[1.2ex]
  \mathrm{C}\sigma_x\ket{01}&=\ket{0}\braket{0}{0}\otimes I\ket{1}+\ket{1}\braket{1}{0}\otimes\sigma_x\ket{1}
     =\ket{0}\otimes\ket{1}+0=\ket{01},\\[1.2ex]
  \mathrm{C}\sigma_x\ket{10}&=\ket{0}\braket{0}{1}\otimes I\ket{0}+\ket{1}\braket{1}{1}\otimes\sigma_x\ket{0}
     =0+\ket{1}\otimes\ket{1}=\ket{11},\\[1.2ex]
  \mathrm{C}\sigma_x\ket{11}&=\ket{0}\braket{0}{1}\otimes I\ket{1}+\ket{1}\braket{1}{1}\otimes\sigma_x\ket{1}
     =0+\ket{1}\otimes\ket{0}=\ket{10}.
\end{align*}
On retrouve la table du CNOT : la cible n'est inversée que si le contrôle vaut $1$.

\paragraph{Un autre exemple : $U=Z$.} Ici
\[
  \mathrm{C}Z=\begin{pmatrix}I&0\\0&Z\end{pmatrix}=\mathrm{diag}(1,1,1,-1) :
\]
la porte multiplie $\ket{11}$ par $-1$ et ne touche pas les autres états de base. Elle est
symétrique entre les deux qubits. Elle se déduit du CNOT par
$\mathrm{C}Z=(I\otimes H)\,\mathrm{C}X\,(I\otimes H)$, où $\mathrm{C}X$ est le contrôlé-$X$.
En effet, $I\otimes H=\mathrm{diag}(H,H)$ (deux blocs $2\times2$ égaux à $H$), et le produit de
matrices diagonales par blocs se fait bloc par bloc :
\[
  (I\otimes H)\begin{pmatrix}I&0\\0&X\end{pmatrix}(I\otimes H)
  =\begin{pmatrix}H\,I\,H&0\\0&H\,X\,H\end{pmatrix}
  =\begin{pmatrix}I&0\\0&HXH\end{pmatrix},
\]
car $H^2=I$, et
\[
  HXH=\frac12\begin{pmatrix}1&1\\1&-1\end{pmatrix}\begin{pmatrix}0&1\\1&0\end{pmatrix}\begin{pmatrix}1&1\\1&-1\end{pmatrix}
  =\frac12\begin{pmatrix}1&1\\1&-1\end{pmatrix}\begin{pmatrix}1&-1\\1&1\end{pmatrix}
  =\frac12\begin{pmatrix}2&0\\0&-2\end{pmatrix}=Z
\]
(section~\ref{sec:hadamard} : $H$ échange les axes $x$ et $z$).

\subsubsection{Contrôlé-SWAP : la porte de Fredkin}

\begin{definition}{Porte de Fredkin (C-SWAP)}
Sur trois qubits $(x,y,z)$ (espace de dimension $2^3=8$), avec $x$ pour contrôle,
\[
  \mathrm{C\text{-}SWAP}=\ket{0}\bra{0}_x\otimes I_y\otimes I_z+\ket{1}\bra{1}_x\otimes\mathrm{SWAP}_{yz}.
\]
Si $x=0$, rien ne change ; si $x=1$, les qubits $y$ et $z$ sont échangés.
\end{definition}

Calcul sur quatre états de base (avec $I\ket{ab}=\ket{ab}$ et $\mathrm{SWAP}\ket{ab}=\ket{ba}$) :
\begin{align*}
  \mathrm{C\text{-}SWAP}\ket{000}&=\ket{0}\braket{0}{0}\otimes I\ket{00}+\ket{1}\braket{1}{0}\otimes\mathrm{SWAP}\ket{00}
     =\ket{0}\otimes\ket{00}+0=\ket{000},\\[1.2ex]
  \mathrm{C\text{-}SWAP}\ket{101}&=\ket{0}\braket{0}{1}\otimes I\ket{01}+\ket{1}\braket{1}{1}\otimes\mathrm{SWAP}\ket{01}
     =0+\ket{1}\otimes\ket{10}=\ket{110},\\[1.2ex]
  \mathrm{C\text{-}SWAP}\ket{110}&=\ket{0}\braket{0}{1}\otimes I\ket{10}+\ket{1}\braket{1}{1}\otimes\mathrm{SWAP}\ket{10}
     =0+\ket{1}\otimes\ket{01}=\ket{101},\\[1.2ex]
  \mathrm{C\text{-}SWAP}\ket{010}&=\ket{0}\braket{0}{0}\otimes I\ket{10}+\ket{1}\braket{1}{0}\otimes\mathrm{SWAP}\ket{10}
     =\ket{0}\otimes\ket{10}+0=\ket{010}.
\end{align*}
Les huit états de base donnent :
\begin{center}
\renewcommand{\arraystretch}{1.3}
\begin{tabular}{@{}ccccccccc@{}}
\toprule
Entrée & $\ket{000}$ & $\ket{001}$ & $\ket{010}$ & $\ket{011}$ & $\ket{100}$ & $\ket{101}$ & $\ket{110}$ & $\ket{111}$ \\
\midrule
Sortie & $\ket{000}$ & $\ket{001}$ & $\ket{010}$ & $\ket{011}$ & $\ket{100}$ & $\ket{110}$ & $\ket{101}$ & $\ket{111}$ \\
\bottomrule
\end{tabular}
\end{center}
Seuls $\ket{101}$ et $\ket{110}$ changent (contrôle à $1$ et $y\neq z$) ; pour $\ket{100}$ et
$\ket{111}$ l'échange de deux qubits égaux ne change rien. Dans la base
$\ket{000},\ket{001},\dots,\ket{111}$ (ordre binaire), la matrice est
$\mathrm{diag}(I_4,\mathrm{SWAP})$ (deux blocs $4\times4$) :
\[
  \mathrm{C\text{-}SWAP}=\begingroup\setlength{\arraycolsep}{3pt}\begin{pmatrix}
    1&0&0&0&0&0&0&0\\
    0&1&0&0&0&0&0&0\\
    0&0&1&0&0&0&0&0\\
    0&0&0&1&0&0&0&0\\
    0&0&0&0&1&0&0&0\\
    0&0&0&0&0&0&1&0\\
    0&0&0&0&0&1&0&0\\
    0&0&0&0&0&0&0&1
  \end{pmatrix}\endgroup .
\]
C'est l'identité dont les lignes $6$ et $7$ sont échangées.

\paragraph{Contrôle en superposition.} Avec $\ket{+}_x\otimes\ket{0}_y\otimes\ket{1}_z
=\frac{1}{\sqrt2}\big(\ket{001}+\ket{101}\big)$ :
\[
  \mathrm{C\text{-}SWAP}\,\frac{1}{\sqrt2}\big(\ket{001}+\ket{101}\big)
  =\frac{1}{\sqrt2}\big(\ket{001}+\ket{110}\big).
\]
Le résultat n'est pas un produit : les coefficients de $\ket{x}\otimes\ket{yz}$ forment la
matrice (lignes $x=0,1$, colonnes $yz=00,01,10,11$)
\[
  \frac{1}{\sqrt2}\begin{pmatrix}0&1&0&0\\0&0&1&0\end{pmatrix},
\]
dont les deux lignes ne sont pas proportionnelles. Le qubit $x$ est intriqué avec la paire
$(y,z)$.

\subsubsection{Contrôlé-contrôlé-NOT : la porte de Toffoli}

\begin{definition}{Porte de Toffoli (CCNOT)}
Sur trois qubits $(x,y,z)$, avec $x$ et $y$ pour contrôles et $z$ pour cible,
\[
  \mathrm{CCNOT}=\ket{0}\bra{0}_x\otimes I_y\otimes I_z
   +\ket{1}\bra{1}_x\otimes\Big[\ket{0}\bra{0}_y\otimes I_z+\ket{1}\bra{1}_y\otimes\sigma_x\Big].
\]
Le crochet est un CNOT (contrôle $y$, cible $z$), appliqué seulement si $x=1$ : si $x=0$, rien ;
si $x=1$ et $y=0$, rien ; si $x=1$ et $y=1$, le qubit $z$ est inversé.
\end{definition}

Calcul de $\mathrm{CCNOT}\ket{110}$ :
\begin{align*}
  \mathrm{CCNOT}\ket{110}&=\ket{0}\braket{0}{1}\otimes I\ket{1}\otimes I\ket{0}
     +\ket{1}\braket{1}{1}\otimes\Big[\ket{0}\braket{0}{1}\otimes I\ket{0}
       +\ket{1}\braket{1}{1}\otimes\sigma_x\ket{0}\Big]\\
  &=0+\ket{1}\otimes\Big[0+\ket{1}\otimes\ket{1}\Big]=\ket{111}.
\end{align*}
Pour les huit états de base, $\mathrm{CCNOT}\ket{x,y,z}=\ket{x,\,y,\,z\oplus xy}$ :
\begin{center}
\renewcommand{\arraystretch}{1.3}
\begin{tabular}{@{}ccccccccc@{}}
\toprule
Entrée & $\ket{000}$ & $\ket{001}$ & $\ket{010}$ & $\ket{011}$ & $\ket{100}$ & $\ket{101}$ & $\ket{110}$ & $\ket{111}$ \\
\midrule
Sortie & $\ket{000}$ & $\ket{001}$ & $\ket{010}$ & $\ket{011}$ & $\ket{100}$ & $\ket{101}$ & $\ket{111}$ & $\ket{110}$ \\
\bottomrule
\end{tabular}
\end{center}
Seuls $\ket{110}$ et $\ket{111}$ sont échangés : la matrice est $\mathrm{diag}(I_6,X)$
(un bloc $6\times6$ égal à l'identité, puis $X$),
\[
  \mathrm{CCNOT}=\begingroup\setlength{\arraycolsep}{3pt}\begin{pmatrix}
    1&0&0&0&0&0&0&0\\
    0&1&0&0&0&0&0&0\\
    0&0&1&0&0&0&0&0\\
    0&0&0&1&0&0&0&0\\
    0&0&0&0&1&0&0&0\\
    0&0&0&0&0&1&0&0\\
    0&0&0&0&0&0&0&1\\
    0&0&0&0&0&0&1&0
  \end{pmatrix}\endgroup .
\]

\begin{remarque}[Portes réversibles]
Les portes de Fredkin et de Toffoli sont des permutations des états de base qui échangent
seulement deux d'entre eux : elles sont leurs propres inverses ($\mathrm{C\text{-}SWAP}^2=\mathrm{CCNOT}^2=\mathrm{Id}$)
et unitaires. Avec $z=0$, la Toffoli calcule le ET logique $x\wedge y$ sans rien effacer :
$\ket{x,y,0}\to\ket{x,y,xy}$ ; avec $z=1$, elle calcule le NON-ET. Elle sert à écrire tout calcul
classique de façon réversible.
\end{remarque}

% ---------------------------------------------------------------------
\subsection{Récapitulatif du chapitre}
% ---------------------------------------------------------------------

\begin{center}
\renewcommand{\arraystretch}{1.5}
\begin{tabular}{@{}lp{12cm}@{}}
\toprule
\textbf{Notion} & \textbf{Formule ou résultat} \\
\midrule
CNOT & $\ket{b,a}\to\ket{b\oplus a,a}$ (contrôle $a$ écrit en dernier), $\mathrm{CNOT}^2=\mathrm{Id}$ \\
SWAP & $\mathrm{SWAP}\ket{a,b}=\ket{b,a}$, \quad $\mathrm{CNOT}_{21}=\mathrm{SWAP}\,\mathrm{CNOT}_{12}\,\mathrm{SWAP}$ \\
Circuit de Bell & $(I\otimes H)$ puis CNOT : $\ket{\beta^{xy}}=\frac{1}{\sqrt2}\big(\ket{y,0}+(-1)^x\ket{\bar y,1}\big)$, $\ket{\beta^{00}}=\ket{\Phi^+}$ \\
Information & $I(A;B)=H(A)+H(B)-H(A,B)$, \quad $H(A|B)=H(A,B)-H(B)$ \\
Trace partielle & $\operatorname{Tr}_B\big(\ket{a}\bra{a'}\otimes\ket{b}\bra{b'}\big)=\ket{a}\bra{a'}\braket{b'}{b}$ ; état de Bell : $\rho_A=\frac12\mathrm{Id}$ \\
État de Bell (bits) & $H(A,B)=0$, \quad $H(A)=H(B)=1$, \quad $I(A;B)=2$, \quad $H(A|B)=-1$ \\
Pur / mélangé & $\operatorname{Tr}\rho^2=1$ / $<1$ ; $\ket{+}$ : $1$ ; qubit d'une paire de Bell : $\frac12$ ; source de Bob : $\frac58$ \\
Téléportation & $1$ e-bit $+\ 2$ bits classiques ; correction : $X$ si $a=1$, puis $Z$ si $q=1$ ; $P(a,q)=\frac14$ \\
Contrôlé-$U$ & $\ket{0}\bra{0}\otimes I+\ket{1}\bra{1}\otimes U=\mathrm{diag}(I,U)$ ; $\ket{x,y}\to\ket{x}\otimes U^x\ket{y}$ \\
Fredkin & $\ket{0}\bra{0}\otimes I\otimes I+\ket{1}\bra{1}\otimes\mathrm{SWAP}$ ; $\ket{101}\leftrightarrow\ket{110}$ \\
Toffoli & $\ket{x,y,z}\to\ket{x,y,z\oplus xy}$ ; $\ket{110}\leftrightarrow\ket{111}$ \\
\bottomrule
\end{tabular}
\end{center}
.
%  Comme chapitre4.tex, chapitre5.tex et chapitre6.tex, il ouvre une NOUVELLE
%  SECTION (numéro 5). Pour le rattacher à la section précédente, remplacer
%  \section par \subsection et chaque \subsection par \subsubsection.
%  Pas de préambule : environnements (definition, resultat, remarque),
%  macros (\ii, \ee, \dd, \ket, \bra, \braket) et couleurs (insarouge,
%  insagris) sont définis dans main.tex. Les symboles de circuits (contrôle,
%  cible, croix, mesure) sont définis juste après avec \tikzset.
% =====================================================================

% Macros de Dirac (déjà définies dans main.tex récent ; ce bloc évite l'erreur
% « Undefined control sequence \ket » si main.tex est une ancienne version).
\providecommand{\ket}[1]{|#1\rangle}
\providecommand{\bra}[1]{\langle #1|}
\providecommand{\braket}[2]{\langle #1|#2\rangle}

% Symboles de circuits pour les figures de ce chapitre.
\tikzset{
  ctl/.pic   = {\fill[insarouge] (0,0) circle (2.6pt);},
  cible/.pic = {\draw[insarouge, thick, fill=white] (0,0) circle (0.28);
                \draw[insarouge, thick] (-0.28,0) -- (0.28,0);
                \draw[insarouge, thick] (0,-0.28) -- (0,0.28);},
  croix/.pic = {\draw[insarouge, thick] (-0.17,-0.17) -- (0.17,0.17);
                \draw[insarouge, thick] (-0.17,0.17) -- (0.17,-0.17);},
  mesure/.pic = {\draw[thick, fill=white] (-0.3,-0.27) rectangle (0.3,0.27);
                 \draw[thick] (-0.19,-0.13) arc[start angle=155, end angle=25, radius=0.21];
                 \draw[thick, -{Stealth[length=1.3mm]}] (0,-0.17) -- (0.14,0.12);},
  porte/.style = {draw=insarouge, fill=insarouge!8, minimum size=7mm, thick},
}

% =====================================================================
\section{Portes multi-qubits, circuits et téléportation}\label{sec:circuits}
% =====================================================================

La section précédente a décrit les états de deux qubits. On étudie ici les portes qui
agissent sur plusieurs qubits : le CNOT et le SWAP, puis les portes contrôlées à deux et à
trois qubits (contrôlé-$U$, Fredkin, Toffoli). On assemble ensuite des portes dans un
circuit qui fabrique un état de Bell, on mesure l'information contenue dans un tel état,
on distingue les états purs des états mélangés, et on démontre le protocole de
téléportation. Tous les calculs sont détaillés.

\begin{remarque}[Convention pour les circuits]
Dans un circuit, le temps va de gauche à droite et chaque qubit a un fil. Les fils sont
dessinés dans l'ordre \emph{inverse} de l'écriture des kets : le fil du \emph{haut} porte le
\emph{dernier} chiffre du ket, le fil du \emph{bas} le premier. Par exemple, si le qubit $A$
est dessiné au-dessus du qubit $B$, le ket s'écrit $\ket{b,a}=\ket{b}_B\otimes\ket{a}_A$.
La figure~\ref{fig:bell} de la section précédente suit la même convention.
\end{remarque}

% ---------------------------------------------------------------------
\subsection{Porte CNOT}\label{sec:cnot}
% ---------------------------------------------------------------------

\begin{definition}{Porte CNOT (NOT contrôlé)}
La porte CNOT agit sur deux qubits : le qubit $A$ est le \emph{contrôle}, le qubit $B$ est la
\emph{cible}. La cible est inversée (porte $X$) si le contrôle vaut $1$ et reste inchangée
s'il vaut $0$. Avec le système écrit $(B,A)$, soit $\ket{b,a}=\ket{b}_B\otimes\ket{a}_A$ :
\[
  \mathrm{CNOT}\,\ket{b,a}=\ket{b\oplus a,\,a},
  \qquad
  \mathrm{CNOT}=I\otimes\ket{0}\bra{0}+X\otimes\ket{1}\bra{1},
\]
où $\oplus$ est l'addition modulo $2$ ($0\oplus0=1\oplus1=0$ et $0\oplus1=1\oplus0=1$).
\end{definition}

Sur le circuit de la figure~\ref{fig:cnot}, le point noir posé sur le fil de $A$ marque le
contrôle et le symbole $\oplus$ posé sur le fil de $B$ marque la cible.

\begin{figure}[!ht]
\centering
\begin{tikzpicture}[>=Stealth, font=\small]
  \draw (0,1.2) -- (6,1.2);
  \draw (0,0) -- (6,0);
  \node[anchor=east] at (0,1.2) {$\ket{a}_A$};
  \node[anchor=east] at (0,0) {$\ket{b}_B$};
  \node[anchor=west] at (6,1.2) {$\ket{a}_A$};
  \node[anchor=west] at (6,0) {$\ket{b\oplus a}_B$};
  \draw[insarouge, thick] (3,1.2) -- (3,0);
  \pic at (3,1.2) {ctl};
  \pic at (3,0) {cible};
  \node[anchor=south, text=insagris] at (3,1.5) {qubit de contrôle};
  \node[anchor=north, text=insagris] at (3,-0.5) {qubit cible};
\end{tikzpicture}
\caption{Porte CNOT : le fil du haut ($A$) est le contrôle, le fil du bas ($B$) est la
cible. Le ket s'écrit $\ket{b,a}$ (le fil du haut donne le dernier chiffre).}
\label{fig:cnot}
\end{figure}

\subsubsection{Table de vérité}

\begin{center}
\renewcommand{\arraystretch}{1.3}
\begin{tabular}{@{}cccc@{}}
\toprule
Entrée $\ket{b,a}$ & Contrôle $a$ & Sortie $\ket{b\oplus a,a}$ & Effet sur la cible $B$ \\
\midrule
$\ket{00}$ & $0$ & $\ket{00}$ & aucun \\
$\ket{01}$ & $1$ & $\ket{11}$ & $0\to1$ \\
$\ket{10}$ & $0$ & $\ket{10}$ & aucun \\
$\ket{11}$ & $1$ & $\ket{01}$ & $1\to0$ \\
\bottomrule
\end{tabular}
\end{center}

\subsubsection{Matrice du CNOT}

La $j$-ième colonne d'une matrice $4\times4$ est l'image du $j$-ième vecteur de base
($\ket{00},\ket{01},\ket{10},\ket{11}$, dans cet ordre). La table donne :
\begin{align*}
  \mathrm{CNOT}\ket{00}&=\ket{00}=\begin{pmatrix}1\\0\\0\\0\end{pmatrix},&
  \mathrm{CNOT}\ket{01}&=\ket{11}=\begin{pmatrix}0\\0\\0\\1\end{pmatrix},\\[1.5ex]
  \mathrm{CNOT}\ket{10}&=\ket{10}=\begin{pmatrix}0\\0\\1\\0\end{pmatrix},&
  \mathrm{CNOT}\ket{11}&=\ket{01}=\begin{pmatrix}0\\1\\0\\0\end{pmatrix},
\end{align*}
d'où, en juxtaposant ces quatre colonnes,
\[
  \mathrm{CNOT}=\begin{pmatrix}
    1 & 0 & 0 & 0\\ 0 & 0 & 0 & 1\\ 0 & 0 & 1 & 0\\ 0 & 1 & 0 & 0
  \end{pmatrix}.
\]

On retrouve cette matrice avec la formule $I\otimes\ket{0}\bra{0}+X\otimes\ket{1}\bra{1}$ et
la règle $A\otimes B=(A_{ij}B)$ (section~\ref{sec:op2}), où
$\ket{0}\bra{0}=\begin{pmatrix}1&0\\0&0\end{pmatrix}$ et
$\ket{1}\bra{1}=\begin{pmatrix}0&0\\0&1\end{pmatrix}$ :
\[
  I\otimes\ket{0}\bra{0}=\begin{pmatrix}1&0&0&0\\0&0&0&0\\0&0&1&0\\0&0&0&0\end{pmatrix},
  \qquad
  X\otimes\ket{1}\bra{1}=\begin{pmatrix}0&0&0&0\\0&0&0&1\\0&0&0&0\\0&1&0&0\end{pmatrix},
\]
et la somme de ces deux matrices est la matrice du CNOT.

\subsubsection{Action sur un état quelconque}

Pour $\ket{\Psi}=c_{00}\ket{00}+c_{01}\ket{01}+c_{10}\ket{10}+c_{11}\ket{11}$ :
\[
  \mathrm{CNOT}\begin{pmatrix}c_{00}\\c_{01}\\c_{10}\\c_{11}\end{pmatrix}
  =\begin{pmatrix}
    1\cdot c_{00}+0\cdot c_{01}+0\cdot c_{10}+0\cdot c_{11}\\
    0\cdot c_{00}+0\cdot c_{01}+0\cdot c_{10}+1\cdot c_{11}\\
    0\cdot c_{00}+0\cdot c_{01}+1\cdot c_{10}+0\cdot c_{11}\\
    0\cdot c_{00}+1\cdot c_{01}+0\cdot c_{10}+0\cdot c_{11}
  \end{pmatrix}
  =\begin{pmatrix}c_{00}\\c_{11}\\c_{10}\\c_{01}\end{pmatrix}.
\]
Le CNOT échange les amplitudes de $\ket{01}$ et de $\ket{11}$, c'est-à-dire des deux états
dont le contrôle vaut $1$, et laisse les autres.

\paragraph{Exemple : contrôle en superposition.} Soit $\ket{0}_B\otimes\ket{+}_A
=\frac{1}{\sqrt2}\big(\ket{00}+\ket{01}\big)=\frac{1}{\sqrt2}(1,1,0,0)^{\mathsf T}$,
un état séparable. Alors
\[
  \mathrm{CNOT}\,\frac{1}{\sqrt2}\begin{pmatrix}1\\1\\0\\0\end{pmatrix}
  =\frac{1}{\sqrt2}\begin{pmatrix}1\\0\\0\\1\end{pmatrix}
  =\frac{1}{\sqrt2}\big(\ket{00}+\ket{11}\big)=\ket{\Phi^+}.
\]
Le résultat est intriqué ($D=c_{00}c_{11}-c_{01}c_{10}=\frac12\neq0$ ; section~\ref{sec:separable}) :
un CNOT dont le contrôle est en superposition crée de l'intrication.

\paragraph{Réversibilité.} Deux CNOT successifs redonnent l'état initial :
$\ket{b,a}\to\ket{b\oplus a,a}\to\ket{b\oplus a\oplus a,a}=\ket{b,a}$, donc
$\mathrm{CNOT}^2=\mathrm{Id}$. La matrice est réelle et symétrique, donc
$\mathrm{CNOT}^\dagger\mathrm{CNOT}=\mathrm{CNOT}^2=\mathrm{Id}$ : elle est bien unitaire.

\begin{remarque}[Contrôle en dernier ou en premier]
La matrice dépend de l'ordre d'écriture du ket. Notons $\mathrm{CNOT}_{21}$ le CNOT dont le
contrôle est le \emph{second} chiffre du ket (c'est la matrice ci-dessus) et
$\mathrm{CNOT}_{12}$ celui dont le contrôle est le \emph{premier} chiffre, utilisé à la
section précédente :
\[
  \mathrm{CNOT}_{12}=\ket{0}\bra{0}\otimes I+\ket{1}\bra{1}\otimes X
  =\begin{pmatrix}1&0&0&0\\0&1&0&0\\0&0&0&1\\0&0&1&0\end{pmatrix}.
\]
Un même circuit (contrôle $A$, cible $B$) a donc deux matrices, selon que l'on écrit
$A$ en dernier ($\mathrm{CNOT}_{21}$) ou en premier ($\mathrm{CNOT}_{12}$). Les deux sont
reliées par le SWAP (section suivante).
\end{remarque}

% ---------------------------------------------------------------------
\subsection{Porte SWAP}\label{sec:swap}
% ---------------------------------------------------------------------

\begin{definition}{Porte SWAP (échange)}
La porte SWAP échange les états des deux qubits :
\[
  \mathrm{SWAP}\big(\ket{\psi}\otimes\ket{\phi}\big)=\ket{\phi}\otimes\ket{\psi}.
\]
Sur les états de base, $\mathrm{SWAP}\ket{a,b}=\ket{b,a}$ ; on prolonge par linéarité.
\end{definition}

Son symbole est formé de deux croix reliées par un trait vertical
(figure~\ref{fig:swap}).

\begin{figure}[!ht]
\centering
\begin{tikzpicture}[>=Stealth, font=\small]
  \draw (0,1.2) -- (6,1.2);
  \draw (0,0) -- (6,0);
  \node[anchor=east] at (0,1.2) {$\ket{a}$};
  \node[anchor=east] at (0,0) {$\ket{b}$};
  \node[anchor=west] at (6,1.2) {$\ket{b}$};
  \node[anchor=west] at (6,0) {$\ket{a}$};
  \draw[insarouge, thick] (3,1.2) -- (3,0);
  \pic at (3,1.2) {croix};
  \pic at (3,0) {croix};
\end{tikzpicture}
\caption{Porte SWAP : les deux qubits échangent leurs états.}
\label{fig:swap}
\end{figure}

\subsubsection{Table et matrice}

\begin{center}
\renewcommand{\arraystretch}{1.3}
\begin{tabular}{@{}cc@{}}
\toprule
Entrée & Sortie \\
\midrule
$\ket{00}$ & $\ket{00}$ \\
$\ket{01}$ & $\ket{10}$ \\
$\ket{10}$ & $\ket{01}$ \\
$\ket{11}$ & $\ket{11}$ \\
\bottomrule
\end{tabular}
\end{center}

Les images des quatre vecteurs de base sont les colonnes de la matrice :
\[
  \mathrm{SWAP}\ket{00}=\begin{pmatrix}1\\0\\0\\0\end{pmatrix},\quad
  \mathrm{SWAP}\ket{01}=\begin{pmatrix}0\\0\\1\\0\end{pmatrix},\quad
  \mathrm{SWAP}\ket{10}=\begin{pmatrix}0\\1\\0\\0\end{pmatrix},\quad
  \mathrm{SWAP}\ket{11}=\begin{pmatrix}0\\0\\0\\1\end{pmatrix},
\]
\[
  \mathrm{SWAP}=\begin{pmatrix}
    1 & 0 & 0 & 0\\ 0 & 0 & 1 & 0\\ 0 & 1 & 0 & 0\\ 0 & 0 & 0 & 1
  \end{pmatrix}.
\]
Appliquée à $\ket{01}=(0,1,0,0)^{\mathsf T}$, cette matrice sélectionne sa deuxième colonne :
\[
  \begin{pmatrix}
    1 & 0 & 0 & 0\\ 0 & 0 & 1 & 0\\ 0 & 1 & 0 & 0\\ 0 & 0 & 0 & 1
  \end{pmatrix}
  \begin{pmatrix}0\\1\\0\\0\end{pmatrix}
  =\begin{pmatrix}0\\0\\1\\0\end{pmatrix}=\ket{10}.
\]
Sur un état quelconque, $\mathrm{SWAP}\,(c_{00},c_{01},c_{10},c_{11})^{\mathsf T}
=(c_{00},c_{10},c_{01},c_{11})^{\mathsf T}$ : les amplitudes de $\ket{01}$ et $\ket{10}$
sont échangées.

\subsubsection{Accord avec la définition}

Pour deux états $\ket{\psi}=(a_0,a_1)^{\mathsf T}$ et $\ket{\phi}=(b_0,b_1)^{\mathsf T}$ :
\[
  \mathrm{SWAP}\,(\ket{\psi}\otimes\ket{\phi})
  =\mathrm{SWAP}\begin{pmatrix}a_0b_0\\a_0b_1\\a_1b_0\\a_1b_1\end{pmatrix}
  =\begin{pmatrix}a_0b_0\\a_1b_0\\a_0b_1\\a_1b_1\end{pmatrix}
  =\begin{pmatrix}b_0\\b_1\end{pmatrix}\otimes\begin{pmatrix}a_0\\a_1\end{pmatrix}
  =\ket{\phi}\otimes\ket{\psi},
\]
ce qui est bien la définition. La porte s'écrit aussi avec des projecteurs :
\[
  \mathrm{SWAP}=\sum_{a,b\in\{0,1\}}\ket{a}\bra{b}\otimes\ket{b}\bra{a}.
\]
En effet, pour $\ket{c}\otimes\ket{d}$ on trouve
$\sum_{a,b}\ket{a}\braket{b}{c}\otimes\ket{b}\braket{a}{d}$, et seul le terme $b=c$, $a=d$
est non nul : le résultat est $\ket{d}\otimes\ket{c}$.

\paragraph{Exemples.}
\begin{itemize}[nosep]
  \item Deux états différents : $\mathrm{SWAP}\big(\ket{0}\otimes\ket{+}\big)
        =\mathrm{SWAP}\,\frac{1}{\sqrt2}(1,1,0,0)^{\mathsf T}=\frac{1}{\sqrt2}(1,0,1,0)^{\mathsf T}
        =\ket{+}\otimes\ket{0}$.
  \item L'état $\ket{\psi}=\frac{1}{\sqrt2}\ket{00}+\frac{1}{\sqrt2}\ket{11}$ ne change pas :
        $\mathrm{SWAP}\ket{\psi}=\frac{1}{\sqrt2}\ket{00}+\frac{1}{\sqrt2}\ket{11}=\ket{\psi}$,
        car $\ket{00}$ et $\ket{11}$ sont invariants. De même
        $\mathrm{SWAP}\ket{\Phi^-}=\ket{\Phi^-}$ et $\mathrm{SWAP}\ket{\Psi^+}=\ket{\Psi^+}$, alors que
        $\mathrm{SWAP}\ket{\Psi^-}=\frac{1}{\sqrt2}(\ket{10}-\ket{01})=-\ket{\Psi^-}$.
\end{itemize}

\paragraph{Propriétés.} Échanger deux fois redonne l'état de départ :
$\mathrm{SWAP}^2=\mathrm{Id}$ ; la matrice est réelle et symétrique, donc
$\mathrm{SWAP}^\dagger\mathrm{SWAP}=\mathrm{SWAP}^2=\mathrm{Id}$ : elle est unitaire.

\subsubsection{Lien avec le CNOT}

\begin{resultat}{SWAP échange les rôles de contrôle et de cible}
\[
  \mathrm{CNOT}_{21}=\mathrm{SWAP}\;\mathrm{CNOT}_{12}\;\mathrm{SWAP}.
\]
\end{resultat}

\emph{Sur les états de base.}
$\mathrm{SWAP}\,\mathrm{CNOT}_{12}\,\mathrm{SWAP}\ket{b,a}
=\mathrm{SWAP}\,\mathrm{CNOT}_{12}\ket{a,b}
=\mathrm{SWAP}\ket{a,\,a\oplus b}=\ket{a\oplus b,\,a}=\mathrm{CNOT}_{21}\ket{b,a}$.

\emph{Par les matrices.} Multiplier à gauche par $\mathrm{SWAP}$ échange les lignes $2$ et
$3$, multiplier à droite l'échange les colonnes $2$ et $3$ :
\[
  \mathrm{SWAP}\,\mathrm{CNOT}_{12}
  =\begin{pmatrix}1&0&0&0\\0&0&0&1\\0&1&0&0\\0&0&1&0\end{pmatrix},
  \qquad
  \mathrm{SWAP}\,\mathrm{CNOT}_{12}\,\mathrm{SWAP}
  =\begin{pmatrix}1&0&0&0\\0&0&0&1\\0&0&1&0\\0&1&0&0\end{pmatrix}=\mathrm{CNOT}_{21}.
\]

\begin{remarque}[SWAP avec trois CNOT]
$\mathrm{SWAP}=\mathrm{CNOT}_{12}\,\mathrm{CNOT}_{21}\,\mathrm{CNOT}_{12}$. En effet, sur
$\ket{a,b}$ (on applique les portes de droite à gauche) :
\[
  \ket{a,b}\ \xrightarrow{\ \mathrm{CNOT}_{12}\ }\ \ket{a,\,a\oplus b}
  \ \xrightarrow{\ \mathrm{CNOT}_{21}\ }\ \ket{a\oplus(a\oplus b),\,a\oplus b}=\ket{b,\,a\oplus b}
  \ \xrightarrow{\ \mathrm{CNOT}_{12}\ }\ \ket{b,\,b\oplus(a\oplus b)}=\ket{b,a}.
\]
Les deux membres coïncident sur les états de base, donc sur tous les états par linéarité.
\end{remarque}

% ---------------------------------------------------------------------
\subsection{Un circuit quantique : le générateur d'états de Bell}\label{sec:circuit-bell}
% ---------------------------------------------------------------------

On assemble une porte de Hadamard et un CNOT. Le système est $(A,B)$ avec le ket
$\ket{a,b}=\ket{a}_A\otimes\ket{b}_B$ ; d'après la convention de la section, $B$ est dessiné
en haut et $A$ en bas (figure~\ref{fig:circuit-bell}). Le circuit a deux étapes :
\begin{enumerate}[nosep]
  \item la porte $I\otimes H$ : identité sur $A$, Hadamard sur $B$ ;
  \item un CNOT de contrôle $B$ et de cible $A$ (la matrice $\mathrm{CNOT}_{21}$ de la
        section~\ref{sec:cnot}, puisque le contrôle est le second chiffre).
\end{enumerate}
Les traits pointillés délimitent les états $\ket{\Pi_0}$, $\ket{\Pi_1}$ et $\ket{\Pi_2}$ :
avant les portes, après $I\otimes H$, après le CNOT.

\begin{figure}[!ht]
\centering
\begin{tikzpicture}[>=Stealth, font=\small]
  % fils : B en haut, A en bas
  \draw (0,1.4) -- (7.6,1.4);
  \draw (0,0) -- (7.6,0);
  \node[anchor=east] at (0,1.4) {$\ket{0}_B$};
  \node[anchor=east] at (0,0) {$\ket{0}_A$};
  % première étape : H sur B, identité (pointillée) sur A
  \node[porte] at (2.2,1.4) {$H$};
  \node[draw=insagris, dashed, fill=white, minimum size=7mm] at (2.2,0) {$I$};
  % CNOT : contrôle B (haut), cible A (bas)
  \draw[insarouge, thick] (5,1.4) -- (5,0);
  \pic at (5,1.4) {ctl};
  \pic at (5,0) {cible};
  % accolade de sortie
  \draw[insagris, thick] (7.75,1.6) -- (7.9,1.6) -- (7.9,-0.2) -- (7.75,-0.2);
  \node[anchor=west] at (7.9,0.7) {$\ket{\beta^{00}}_{AB}$};
  % états intermédiaires
  \foreach \x in {0.7,3.6,6.4}
    \draw[dashed, insagris!70] (\x,1.9) -- (\x,-0.7);
  \node[anchor=north, font=\footnotesize] at (0.7,-0.7) {$\ket{\Pi_0}=\ket{00}$};
  \node[anchor=north, font=\footnotesize] at (3.6,-0.7) {$\ket{\Pi_1}=\ket{0}\otimes\ket{+}$};
  \node[anchor=north, font=\footnotesize] at (6.4,-0.7) {$\ket{\Pi_2}=\ket{\beta^{00}}$};
\end{tikzpicture}
\caption{Générateur d'états de Bell : $I\otimes H$, puis un CNOT de contrôle $B$ (fil du haut)
et de cible $A$ (fil du bas). L'entrée représentée est $\ket{00}$.}
\label{fig:circuit-bell}
\end{figure}

\subsubsection{Entrée \texorpdfstring{$\ket{00}$}{|00>} : calcul pas à pas}

\begin{align*}
  \ket{\Pi_0}&=\ket{0}_A\otimes\ket{0}_B=\ket{00},\\[1ex]
  \ket{\Pi_1}&=(I\otimes H)\ket{\Pi_0}=I\ket{0}_A\otimes H\ket{0}_B=\ket{0}\otimes\ket{+}
    =\frac{\ket{00}+\ket{01}}{\sqrt2},\\[1ex]
  \ket{\Pi_2}&=\mathrm{CNOT}\ket{\Pi_1}=\frac{\mathrm{CNOT}\ket{00}+\mathrm{CNOT}\ket{01}}{\sqrt2}
    =\frac{\ket{00}+\ket{11}}{\sqrt2}=\ket{\beta^{00}}=\ket{\Phi^+}.
\end{align*}
Dans $\ket{\Pi_1}$, le premier chiffre est celui de $A$ et le second celui de $B$. Le CNOT a
pour contrôle $B$ : $\ket{00}$ ($B=0$) ne change pas, $\ket{01}$ ($B=1$) devient $\ket{11}$
puisque $A$ est inversé. Le résultat est un \emph{état de Bell}.

Avec les matrices, la règle $A\otimes B=(A_{ij}B)$ donne
\[
  I\otimes H=\begin{pmatrix}1\cdot H&0\\0&1\cdot H\end{pmatrix}
  =\frac{1}{\sqrt2}\begin{pmatrix}1&1&0&0\\1&-1&0&0\\0&0&1&1\\0&0&1&-1\end{pmatrix}.
\]
Sa première colonne est $\ket{\Pi_1}=\frac{1}{\sqrt2}(1,1,0,0)^{\mathsf T}$ (image de $\ket{00}$).
Le CNOT échange les amplitudes de $\ket{01}$ et de $\ket{11}$ (section~\ref{sec:cnot}), donc
\[
  \ket{\Pi_2}=\mathrm{CNOT}\,\frac{1}{\sqrt2}\begin{pmatrix}1\\1\\0\\0\end{pmatrix}
  =\frac{1}{\sqrt2}\begin{pmatrix}1\\0\\0\\1\end{pmatrix}.
\]

\subsubsection{Les quatre entrées}

Prenons $B$ dans l'état $\ket{x}$ et $A$ dans l'état $\ket{y}$, avec $x,y\in\{0,1\}$. Alors
$\ket{\Pi_0}=\ket{y}_A\otimes\ket{x}_B=\ket{y,x}$. Comme
$H\ket{x}=\frac{1}{\sqrt2}\big(\ket{0}+(-1)^x\ket{1}\big)$ (c'est $\ket{+}$ pour $x=0$ et
$\ket{-}$ pour $x=1$),
\[
  \ket{\Pi_1}=\ket{y}\otimes H\ket{x}=\frac{1}{\sqrt2}\Big(\ket{y,0}+(-1)^x\ket{y,1}\Big).
\]
Le CNOT (contrôle $B$) ne change pas $\ket{y,0}$ et transforme $\ket{y,1}$ en
$\ket{\bar y,1}$, avec $\bar y=1-y$ :
\[
  \ket{\Pi_2}=\ket{\beta^{xy}}=\frac{1}{\sqrt2}\Big(\ket{y,0}+(-1)^x\ket{\bar y,1}\Big).
\]
Le premier indice de $\ket{\beta^{xy}}$ est l'entrée de $B$ (le qubit qui reçoit $H$), le
second celle de $A$ (la cible). Les quatre cas :

\begin{center}
\renewcommand{\arraystretch}{1.7}
\begin{tabular}{@{}ccccc@{}}
\toprule
$(B,A)=(x,y)$ & $\ket{\Pi_0}=\ket{y,x}$ & $\ket{\Pi_1}$ & $\ket{\Pi_2}=\ket{\beta^{xy}}$ & Nom \\
\midrule
$(0,0)$ & $\ket{00}$ & $\frac{1}{\sqrt2}(\ket{00}+\ket{01})$ & $\frac{1}{\sqrt2}(\ket{00}+\ket{11})$ & $\ket{\Phi^+}$ \\
$(1,0)$ & $\ket{01}$ & $\frac{1}{\sqrt2}(\ket{00}-\ket{01})$ & $\frac{1}{\sqrt2}(\ket{00}-\ket{11})$ & $\ket{\Phi^-}$ \\
$(0,1)$ & $\ket{10}$ & $\frac{1}{\sqrt2}(\ket{10}+\ket{11})$ & $\frac{1}{\sqrt2}(\ket{10}+\ket{01})$ & $\ket{\Psi^+}$ \\
$(1,1)$ & $\ket{11}$ & $\frac{1}{\sqrt2}(\ket{10}-\ket{11})$ & $\frac{1}{\sqrt2}(\ket{10}-\ket{01})$ & $-\ket{\Psi^-}$ \\
\bottomrule
\end{tabular}
\end{center}
Détail de la dernière ligne : $\ket{\Pi_1}=\ket{1}\otimes\ket{-}=\frac{1}{\sqrt2}(\ket{10}-\ket{11})$ ;
le CNOT laisse $\ket{10}$ et envoie $\ket{11}$ sur $\ket{01}$, d'où
$\frac{1}{\sqrt2}(\ket{10}-\ket{01})=-\frac{1}{\sqrt2}(\ket{01}-\ket{10})=-\ket{\Psi^-}$.

\paragraph{Forme matricielle.} $\mathrm{CNOT}_{21}$ échange les lignes $2$ et $4$ de
$I\otimes H$ (il échange les amplitudes de $\ket{01}$ et $\ket{11}$), donc
\[
  \mathrm{CNOT}_{21}\,(I\otimes H)=\frac{1}{\sqrt2}\begin{pmatrix}
    1 & 1 & 0 & 0\\ 0 & 0 & 1 & -1\\ 0 & 0 & 1 & 1\\ 1 & -1 & 0 & 0
  \end{pmatrix}.
\]
Ses colonnes sont les images de $\ket{00},\ket{01},\ket{10},\ket{11}$, c'est-à-dire de
$(B,A)=(0,0),(1,0),(0,1),(1,1)$ : ce sont $\ket{\beta^{00}}$, $\ket{\beta^{10}}$,
$\ket{\beta^{01}}$, $\ket{\beta^{11}}$, soit $\ket{\Phi^+}$, $\ket{\Phi^-}$, $\ket{\Psi^+}$
et $-\ket{\Psi^-}$.

\begin{remarque}[Comparaison avec la section précédente]
Le circuit de la figure~\ref{fig:bell} applique $H$ au \emph{premier} chiffre et prend celui-ci
pour contrôle ; ici $H$ et le contrôle portent sur le \emph{dernier} chiffre. Les deux
circuits se déduisent l'un de l'autre en échangeant les deux fils (conjugaison par le SWAP,
section~\ref{sec:swap}). Les quatre états de Bell obtenus sont les mêmes, avec les mêmes
indices, sauf le dernier : $-\ket{\Psi^-}$ ici, $+\ket{\Psi^-}$ là-bas. Ce signe est une phase
globale, sans effet sur les mesures (elle disparaît de $\rho$, section~\ref{sec:rho}).
\end{remarque}

\subsubsection{Matrice de densité de l'état de sortie}

Pour $\ket{\Psi}=\ket{\beta^{00}}=\frac{1}{\sqrt2}\ket{00}+\frac{1}{\sqrt2}\ket{11}$ :
\[
  \rho_{AB}=\ket{\Psi}\bra{\Psi}
  =\begin{pmatrix}\frac{1}{\sqrt2}\\0\\0\\\frac{1}{\sqrt2}\end{pmatrix}
   \begin{pmatrix}\frac{1}{\sqrt2}&0&0&\frac{1}{\sqrt2}\end{pmatrix}
  =\frac12\begin{pmatrix}1&0&0&1\\0&0&0&0\\0&0&0&0\\1&0&0&1\end{pmatrix}.
\]
Sous forme de ket-bra :
$\rho_{AB}=\frac12\big(\ket{00}\bra{00}+\ket{00}\bra{11}+\ket{11}\bra{00}+\ket{11}\bra{11}\big)$.
On vérifie $\operatorname{Tr}\rho_{AB}=\frac12(1+1)=1$ et
\[
  \rho_{AB}^2=\frac14\begin{pmatrix}2&0&0&2\\0&0&0&0\\0&0&0&0\\2&0&0&2\end{pmatrix}=\rho_{AB},
\]
donc $\operatorname{Tr}\rho_{AB}^2=1$ : le couple $(A,B)$ est dans un état \emph{pur}.

% ---------------------------------------------------------------------
\subsection{Information : variables aléatoires et états quantiques}\label{sec:info}
% ---------------------------------------------------------------------

\begin{definition}{Entropie, information mutuelle, entropie conditionnelle}
Soit $A$ une variable aléatoire de loi $p(a)$. Son \emph{entropie de Shannon}, en bits, est
\[
  H(A)=-\sum_a p(a)\log_2p(a)\qquad(\text{avec }0\log_20=0).
\]
Elle mesure l'incertitude moyenne sur $A$ : une pièce équilibrée a $H=1$ bit. Pour un couple
$(A,B)$ de loi $p(a,b)$, $H(A,B)=-\sum_{a,b}p(a,b)\log_2p(a,b)$, et
\begin{align*}
  I(A;B)&=H(A)+H(B)-H(A,B) &&\text{(information mutuelle)},\\
  H(A|B)&=H(A,B)-H(B)=H(A)-I(A;B) &&\text{(entropie conditionnelle)}.
\end{align*}
\end{definition}

Les deux écritures de $H(A|B)$ sont égales : la définition de $I(A;B)$ donne
$H(A,B)=H(A)+H(B)-I(A;B)$, donc $H(A,B)-H(B)=H(A)-I(A;B)$. L'information mutuelle $I(A;B)$
mesure ce que la connaissance de $B$ apprend sur $A$ ; $H(A|B)$ est l'incertitude qui reste sur
$A$ quand on connaît $B$.

\subsubsection{Deux variables aléatoires indépendantes}

Alice et Bob lancent chacun une pièce équilibrée, sans lien entre les deux lancers :
$p(a,b)=\frac14$ pour les quatre couples. Chaque pièce donne $0$ ou $1$ avec la probabilité $\frac12$, donc
\begin{align*}
  H(A)&=-\Big(\tfrac12\log_2\tfrac12+\tfrac12\log_2\tfrac12\Big)=1\text{ bit},\\
  H(B)&=1\text{ bit (même calcul)},\\
  H(A,B)&=-4\cdot\tfrac14\log_2\tfrac14=2\text{ bits},\\
  I(A;B)&=H(A)+H(B)-H(A,B)=1+1-2=0,\\
  H(A|B)&=H(A)-I(A;B)=1-0=H(A)=1\text{ bit}.
\end{align*}
Connaître $B$ n'apprend rien sur $A$.

\subsubsection{Deux variables solidaires}

Les deux pièces tombent toujours du même côté : $p(0,0)=p(1,1)=\frac12$ et
$p(0,1)=p(1,0)=0$. Les lois de chaque pièce sont inchangées ($p(a{=}0)=p(0,0)+p(0,1)=\frac12$,
etc.), donc $H(A)=H(B)=1$ bit, mais
\begin{align*}
  H(A,B)&=-2\cdot\tfrac12\log_2\tfrac12-2\cdot0\log_20=1\text{ bit},\\
  I(A;B)&=1+1-1=1\text{ bit},\qquad H(A|B)=1-1=0.
\end{align*}
Connaître $B$ détermine $A$ : plus aucune incertitude ne reste.

\subsubsection{Deux qubits dans un état de Bell}

Pour un système quantique, on remplace la loi de probabilité par la matrice de densité.

\begin{definition}{Trace partielle et entropie de von Neumann}
L'état du seul qubit $A$ d'un système $(A,B)$ de matrice de densité $\rho_{AB}$ est obtenu en
\emph{traçant sur $B$} : $\rho_A=\operatorname{Tr}_B(\rho_{AB})$. La trace partielle est
linéaire et remplace la partie $B$ de chaque terme par sa trace :
\[
  \operatorname{Tr}_B\big(\ket{a}\bra{a'}\otimes\ket{b}\bra{b'}\big)=\ket{a}\bra{a'}\;\braket{b'}{b}.
\]
L'\emph{entropie de von Neumann} d'un état $\rho$, de valeurs propres $\lambda_i$, est
$S(\rho)=-\operatorname{Tr}(\rho\log_2\rho)=-\sum_i\lambda_i\log_2\lambda_i$. En quantique,
$H(A)=S(\rho_A)$, $H(B)=S(\rho_B)$ et $H(A,B)=S(\rho_{AB})$ ; les définitions de
$I(A;B)$ et de $H(A|B)$ sont les mêmes.
\end{definition}

On part de l'état de Bell
$\rho_{AB}=\frac12\big(\ket{00}\bra{00}+\ket{00}\bra{11}+\ket{11}\bra{00}+\ket{11}\bra{11}\big)$
produit par le circuit de la section~\ref{sec:circuit-bell}. Chaque terme est un produit
d'opérateurs à un qubit :
$\ket{00}\bra{11}=\ket{0}\bra{1}\otimes\ket{0}\bra{1}$, etc. Alors
\begin{align*}
  \rho_A=\operatorname{Tr}_B(\rho_{AB})
  &=\tfrac12\Big(\ket{0}\bra{0}\,\braket{0}{0}+\ket{0}\bra{1}\,\braket{1}{0}
      +\ket{1}\bra{0}\,\braket{0}{1}+\ket{1}\bra{1}\,\braket{1}{1}\Big)\\
  &=\tfrac12\Big(\ket{0}\bra{0}\cdot1+\ket{0}\bra{1}\cdot0+\ket{1}\bra{0}\cdot0+\ket{1}\bra{1}\cdot1\Big)\\
  &=\tfrac12\big(\ket{0}\bra{0}+\ket{1}\bra{1}\big)=\tfrac12\,\mathrm{Id}
  =\begin{pmatrix}\frac12&0\\0&\frac12\end{pmatrix}.
\end{align*}
La formule des composantes de la section~\ref{sec:separable},
$(\rho_A)_{ij}=\rho_{(i0),(j0)}+\rho_{(i1),(j1)}$, donne le même résultat. En traçant sur $A$
au lieu de $B$, la règle est $\operatorname{Tr}_A(\ket{a}\bra{a'}\otimes\ket{b}\bra{b'})
=\braket{a'}{a}\,\ket{b}\bra{b'}$, et on trouve de même $\rho_B=\frac12\mathrm{Id}$.

Calcul des entropies :
\begin{itemize}[nosep]
  \item $\rho_{AB}$ est pur : ses valeurs propres sont $1,0,0,0$, donc
        $H(A,B)=S(\rho_{AB})=-1\cdot\log_21=0$ ;
  \item $\rho_A$ et $\rho_B$ ont pour valeurs propres $\frac12,\frac12$, donc
        $H(A)=H(B)=-\big(\frac12\log_2\frac12+\frac12\log_2\frac12\big)=1$ bit.
\end{itemize}
D'où
\begin{align*}
  I(A;B)&=H(A)+H(B)-H(A,B)=1+1-0=2\text{ bits},\\
  H(A|B)&=H(A,B)-H(B)=0-1=-1\text{ bit}\qquad(=H(A)-I(A;B)=1-2).
\end{align*}

\begin{remarque}[Une entropie conditionnelle négative]
Pour des variables classiques, $p(a,b)\leq p(b)$, donc $-\log_2p(a,b)\geq-\log_2p(b)$ et
$H(A,B)\geq H(B)$ : on a toujours $H(A|B)\geq0$, et $I(A;B)\leq H(A)$. Ici $H(A|B)=-1$ et
$I(A;B)=2>H(A)$ : ces valeurs n'ont pas d'équivalent classique, c'est une signature de
l'intrication. Le résultat est surprenant mais mathématiquement correct.
\end{remarque}

\begin{center}
\renewcommand{\arraystretch}{1.4}
\begin{tabular}{@{}lccc@{}}
\toprule
 & indépendantes & solidaires & état de Bell \\
\midrule
$H(A)$ & $1$ & $1$ & $1$ \\
$H(B)$ & $1$ & $1$ & $1$ \\
$H(A,B)$ & $2$ & $1$ & $0$ \\
$I(A;B)$ & $0$ & $1$ & $2$ \\
$H(A|B)$ & $1$ & $0$ & $-1$ \\
\bottomrule
\end{tabular}
\end{center}
(en bits ; les deux premières colonnes sont classiques, la troisième est quantique).

% ---------------------------------------------------------------------
\subsection{Système pur et système mélangé}\label{sec:pur-melange}
% ---------------------------------------------------------------------

Rappel (section~\ref{sec:pur-mixte}) : si $\lambda_i$ sont les valeurs propres de la matrice de
densité $\rho$, alors $\operatorname{Tr}\rho^2=\sum_i\lambda_i^2\leq1$. Le système est
\emph{pur} si $\operatorname{Tr}\rho^2=1$ (alors $\rho=\ket{\psi}\bra{\psi}$ et $\rho^2=\rho$),
et \emph{mélangé} si $\operatorname{Tr}\rho^2<1$. Pour un qubit,
$\operatorname{Tr}\rho^2=\frac{1+|\vec r|^2}{2}$ (section~\ref{sec:bloch}) est compris entre
$\frac12$ et $1$.

\subsubsection{Un système pur : exercice}

\emph{Énoncé.} Soit $\ket{\psi}=\frac{1}{\sqrt2}\ket{0}+\frac{1}{\sqrt2}\ket{1}$.
\begin{enumerate}[label=\alph*), nosep]
  \item Calculer la matrice de densité $\rho$.
  \item Calculer la pureté du système quantique.
  \item Calculer $\langle\sigma_x\rangle$ et $\langle\sigma_z\rangle$.
  \item Écrire $\ket{\psi}$ dans la base $X$.
\end{enumerate}

\paragraph{a) Matrice de densité.} $\rho=\ket{\psi}\bra{\psi}$ est le produit d'un vecteur
colonne par un vecteur ligne (les amplitudes sont réelles) :
\[
  \rho=\begin{pmatrix}\frac{1}{\sqrt2}\\[0.4ex]\frac{1}{\sqrt2}\end{pmatrix}
       \begin{pmatrix}\frac{1}{\sqrt2}&\frac{1}{\sqrt2}\end{pmatrix}
  =\begin{pmatrix}\frac12&\frac12\\[0.4ex]\frac12&\frac12\end{pmatrix}.
\]

\paragraph{b) Pureté.}
\[
  \rho^2=\begin{pmatrix}\frac12&\frac12\\[0.4ex]\frac12&\frac12\end{pmatrix}
         \begin{pmatrix}\frac12&\frac12\\[0.4ex]\frac12&\frac12\end{pmatrix}
  =\begin{pmatrix}\frac14+\frac14&\frac14+\frac14\\[0.4ex]\frac14+\frac14&\frac14+\frac14\end{pmatrix}
  =\begin{pmatrix}\frac12&\frac12\\[0.4ex]\frac12&\frac12\end{pmatrix}=\rho,
\]
donc $\operatorname{Tr}\rho^2=\operatorname{Tr}\rho=\frac12+\frac12=1$ : le système est \emph{pur}.

\paragraph{c) Valeurs moyennes.} Avec $\sigma_x=X=\begin{pmatrix}0&1\\1&0\end{pmatrix}$ et
$\sigma_z=Z=\begin{pmatrix}1&0\\0&-1\end{pmatrix}$ (section~\ref{sec:op-dirac}), la valeur
moyenne est $\langle\sigma\rangle=\operatorname{Tr}(\rho\,\sigma)$ (section~\ref{sec:rho}).
\begin{align*}
  \rho\,\sigma_x&=\begin{pmatrix}\frac12&\frac12\\[0.4ex]\frac12&\frac12\end{pmatrix}
     \begin{pmatrix}0&1\\1&0\end{pmatrix}
     =\begin{pmatrix}\frac12&\frac12\\[0.4ex]\frac12&\frac12\end{pmatrix},
     &\langle\sigma_x\rangle&=\operatorname{Tr}(\rho\,\sigma_x)=\tfrac12+\tfrac12=1,\\[1.5ex]
  \rho\,\sigma_z&=\begin{pmatrix}\frac12&\frac12\\[0.4ex]\frac12&\frac12\end{pmatrix}
     \begin{pmatrix}1&0\\0&-1\end{pmatrix}
     =\begin{pmatrix}\frac12&-\frac12\\[0.4ex]\frac12&-\frac12\end{pmatrix},
     &\langle\sigma_z\rangle&=\operatorname{Tr}(\rho\,\sigma_z)=\tfrac12-\tfrac12=0.
\end{align*}
On retrouve $\langle\sigma_z\rangle$ avec les probabilités de mesure : dans la base $Z$,
$\sigma_z$ vaut $+1$ pour $\ket{0}$ et $-1$ pour $\ket{1}$, et
$|\braket{0}{\psi}|^2=|\braket{1}{\psi}|^2=\frac12$, donc
$\langle\sigma_z\rangle=\frac12\times(+1)+\frac12\times(-1)=0$.

\paragraph{d) Écriture dans la base $X$.} La base $X$ est $\{\ket{+},\ket{-}\}$, orthonormée, donc
$\ket{+}\bra{+}+\ket{-}\bra{-}=\mathrm{Id}$ et
$\ket{\psi}=\ket{+}\braket{+}{\psi}+\ket{-}\braket{-}{\psi}$, avec
\[
  \braket{+}{\psi}=\tfrac{1}{\sqrt2}\Big(\tfrac{1}{\sqrt2}+\tfrac{1}{\sqrt2}\Big)=1,
  \qquad
  \braket{-}{\psi}=\tfrac{1}{\sqrt2}\Big(\tfrac{1}{\sqrt2}-\tfrac{1}{\sqrt2}\Big)=0 .
\]
Donc $\ket{\psi}=1\cdot\ket{+}+0\cdot\ket{-}=\ket{+}$ : dans la base $X$, l'état n'a qu'une
seule composante. Comme $\sigma_x\ket{+}=+\ket{+}$ et $\sigma_x\ket{-}=-\ket{-}$, une
mesure de $\sigma_x$ donne $+1$ avec certitude, et
$\langle\sigma_x\rangle=1\cdot(+1)+0\cdot(-1)=1$ : c'est le résultat de la question c).
Le vecteur de Bloch de $\ket{+}$ est $(\langle\sigma_x\rangle,\langle\sigma_y\rangle,\langle\sigma_z\rangle)=(1,0,0)$,
de norme $1$ : l'état est sur la sphère, en accord avec $\operatorname{Tr}\rho^2=\frac{1+1}{2}=1$.

\subsubsection{Un système mélangé : un qubit de la paire de Bell}

À la section~\ref{sec:info}, la trace sur $B$ de l'état de Bell a donné
$\rho_A=\frac12\mathrm{Id}=\begin{pmatrix}\frac12&0\\0&\frac12\end{pmatrix}$. Alors
\[
  \rho_A^2=\begin{pmatrix}\frac14&0\\0&\frac14\end{pmatrix},
  \qquad
  \operatorname{Tr}\big(\rho_A^2\big)=\frac14+\frac14=\frac12<1 :
\]
le qubit $A$ est dans un état \emph{mélangé}, et c'est le plus mélangé possible pour un qubit
(centre de la boule de Bloch, $\vec r=\vec0$). Pourtant le couple $(A,B)$ est pur
($\operatorname{Tr}\rho_{AB}^2=1$) : on connaît parfaitement l'état du couple, mais pas celui d'un
qubit pris seul. L'information est entièrement dans les corrélations entre $A$ et $B$, et
$S(\rho_A)=1$ bit.

\subsubsection{Un système mélangé : le mélange statistique de Bob}

Bob prépare un qubit dans l'état $\ket{0}$ avec la probabilité $p_1=\frac34$ (75\,\%) et dans
l'état $\ket{1}$ avec la probabilité $p_2=\frac14$ (25\,\%), sans que l'on sache lequel des deux a
été préparé. Ce n'est pas un état pur $\ket{\psi}$ : on décrit le système par la matrice de
densité $\rho=\sum_ip_i\ket{\psi_i}\bra{\psi_i}$ (section~\ref{sec:pur-mixte}) :
\[
  \rho=\underbrace{\tfrac34}_{p_1}\ket{0}\bra{0}+\underbrace{\tfrac14}_{p_2}\ket{1}\bra{1}
  =\begin{pmatrix}\frac34&0\\0&\frac14\end{pmatrix}.
\]
Alors
\[
  \rho^2=\begin{pmatrix}\frac{9}{16}&0\\0&\frac{1}{16}\end{pmatrix},
  \qquad
  \operatorname{Tr}\big(\rho^2\big)=\frac{9}{16}+\frac{1}{16}=\frac{10}{16}=\frac58<1 :
\]
le système est \emph{mélangé}.

Les valeurs moyennes sont les moyennes pondérées sur la préparation :
$\langle\sigma_z\rangle=\operatorname{Tr}(\rho\,\sigma_z)=\frac34-\frac14=\frac12
=p_1\cdot(+1)+p_2\cdot(-1)$, et $\langle\sigma_x\rangle=\operatorname{Tr}(\rho\,\sigma_x)=0$
(les termes hors diagonale de $\rho$ sont nuls). Le vecteur de Bloch est $(0,0,\frac12)$,
à l'intérieur de la boule, et $\operatorname{Tr}\rho^2=\frac{1+1/4}{2}=\frac58$. Son entropie
est $S(\rho)=-\frac34\log_2\frac34-\frac14\log_2\frac14=2-\frac34\log_23\approx0{,}811$ bit.

\begin{remarque}[Mélange et superposition]
Considérons l'état pur $\ket{\psi'}=\frac{\sqrt3}{2}\ket{0}+\frac12\ket{1}$, qui a les mêmes
probabilités en base $Z$ ($P(0)=\frac34$, $P(1)=\frac14$). Sa matrice de densité est
\[
  \rho'=\begin{pmatrix}\frac34&\frac{\sqrt3}{4}\\[0.4ex]\frac{\sqrt3}{4}&\frac14\end{pmatrix},
  \qquad
  \operatorname{Tr}\rho'^2=\frac{9}{16}+\frac1{16}+2\cdot\frac{3}{16}=1,
  \qquad
  \langle\sigma_x\rangle=2\cdot\frac{\sqrt3}{4}=\frac{\sqrt3}{2}.
\]
Le mélange de Bob a la même diagonale mais aucun terme hors diagonale ($\operatorname{Tr}\rho^2=\frac58$,
$\langle\sigma_x\rangle=0$) : une superposition n'est pas un tirage au hasard entre $\ket{0}$ et
$\ket{1}$. De plus, pour un état pur il existe une mesure au résultat certain (ici $\sigma_x$ pour
$\ket{+}$) ; pour le mélange de Bob, la meilleure probabilité d'obtenir un même résultat est
$75\,\%$ (mesure en base $Z$, résultat $\ket{0}$).
\end{remarque}

\subsubsection{Bilan}

\begin{center}
\renewcommand{\arraystretch}{1.5}
\begin{tabular}{@{}lcccl@{}}
\toprule
\textbf{Système} & $\rho$ & $\operatorname{Tr}\rho^2$ & $S(\rho)$ & \textbf{Nature} \\
\midrule
$\ket{\psi}=\ket{+}$ & ${\renewcommand{\arraystretch}{1}\begin{pmatrix}\frac12&\frac12\\\frac12&\frac12\end{pmatrix}}$ & $1$ & $0$ & pur \\[2ex]
Qubit $A$ d'une paire de Bell & ${\renewcommand{\arraystretch}{1}\begin{pmatrix}\frac12&0\\0&\frac12\end{pmatrix}}$ & $\frac12$ & $1$ bit & mélangé \\[2ex]
Source de Bob (75\,\%/25\,\%) & ${\renewcommand{\arraystretch}{1}\begin{pmatrix}\frac34&0\\0&\frac14\end{pmatrix}}$ & $\frac58$ & $\approx0{,}811$ bit & mélangé \\
\bottomrule
\end{tabular}
\end{center}

% ---------------------------------------------------------------------
\subsection{Algorithme quantique : la téléportation}\label{sec:teleportation}
% ---------------------------------------------------------------------

\emph{Problème.} Alice détient un qubit $Q$ dans un état qu'elle ne connaît pas,
$\ket{\varphi}=\alpha\ket{0}+\beta\ket{1}$ avec $|\alpha|^2+|\beta|^2=1$, et elle veut que Bob
obtienne cet état sur l'un de ses qubits. Elle ne peut pas mesurer $Q$ pour connaître
$\alpha$ et $\beta$ (une mesure ne donne qu'un bit et détruit l'état), elle ne peut pas le
copier (théorème de non-clonage) et elle n'envoie pas le qubit lui-même.

\begin{definition}{Protocole de téléportation quantique}
\textbf{Ressources.}
\begin{itemize}[nosep]
  \item Un \emph{e-bit} : Alice et Bob partagent une paire de qubits $(A,B)$ dans l'état de Bell
        $\ket{\Phi^+}_{BA}=\frac{1}{\sqrt2}\big(\ket{00}+\ket{11}\big)$ ; Alice détient $A$,
        Bob détient $B$.
  \item Le qubit $Q$ à téléporter, chez Alice.
  \item Un canal classique (deux bits).
\end{itemize}
Le système est écrit $(B,A,Q)$ : $\ket{b,a,q}=\ket{b}_B\otimes\ket{a}_A\otimes\ket{q}_Q$.

\textbf{Étapes.}
\begin{enumerate}[nosep]
  \item Alice applique un CNOT de contrôle $Q$ et de cible $A$.
  \item Alice applique une porte $H$ sur $Q$.
  \item Alice mesure $A$ et $Q$ dans la base de calcul, obtient deux bits $(a,q)$ et les envoie
        à Bob.
  \item Bob applique $X$ si $a=1$, puis $Z$ si $q=1$, à son qubit $B$.
\end{enumerate}
\end{definition}

\begin{figure}[!ht]
\centering
\begin{adjustbox}{max width=\linewidth}
\begin{tikzpicture}[>=Stealth, font=\small]
  % fils quantiques : Q en haut, A au milieu, B en bas
  \draw (0,2.4) -- (5.3,2.4);
  \draw (0,1.2) -- (5.3,1.2);
  \draw (0,0) -- (10.6,0);
  \node[anchor=east] at (0,2.4) {$\ket{\varphi}_Q$};
  \node[anchor=south west, text=insagris, font=\footnotesize] at (0.05,1.2) {$A$};
  \node[anchor=south west, text=insagris, font=\footnotesize] at (0.05,0) {$B$};
  % paire de Bell partagée
  \draw[insagris, thick] (-0.1,1.5) -- (-0.25,1.5) -- (-0.25,-0.3) -- (-0.1,-0.3);
  \node[anchor=east] at (-0.35,0.6) {$\ket{\Phi^+}_{BA}$};
  % CNOT : contrôle Q, cible A
  \draw[insarouge, thick] (2.6,2.4) -- (2.6,1.2);
  \pic at (2.6,2.4) {ctl};
  \pic at (2.6,1.2) {cible};
  % Hadamard sur Q
  \node[porte] at (4.2,2.4) {$H$};
  % mesures d'Alice
  \pic at (5.6,2.4) {mesure};
  \pic at (5.6,1.2) {mesure};
  % fils classiques (doubles)
  \draw[double, double distance=1.2pt, insagris] (5.9,1.2) -- (7.6,1.2) -- (7.6,0.35);
  \draw[double, double distance=1.2pt, insagris] (5.9,2.4) -- (9.6,2.4) -- (9.6,0.35);
  \node[above, font=\footnotesize] at (6.5,1.2) {$a$};
  \node[above, font=\footnotesize] at (6.5,2.4) {$q$};
  % corrections de Bob
  \node[porte] at (7.6,0) {$X$};
  \node[porte] at (9.6,0) {$Z$};
  \node[anchor=west] at (10.6,0) {$\ket{\varphi}_B$};
  % états intermédiaires
  \foreach \x in {1.4,3.4,4.9}
    \draw[dashed, insagris!70] (\x,2.8) -- (\x,-0.5);
  \node[anchor=north, font=\footnotesize] at (1.4,-0.5) {$\ket{\Pi_0}$};
  \node[anchor=north, font=\footnotesize] at (3.4,-0.5) {$\ket{\Pi_1}$};
  \node[anchor=north, font=\footnotesize] at (4.9,-0.5) {$\ket{\Pi_2}$};
\end{tikzpicture}
\end{adjustbox}
\caption{Circuit de téléportation. Alice détient $Q$ (dans $\ket{\varphi}=\alpha\ket{0}+\beta\ket{1}$)
et $A$, Bob détient $B$ ; $A$ et $B$ forment la paire $\ket{\Phi^+}_{BA}$. Les traits doubles
sont des fils \emph{classiques} qui portent les résultats $a$ (mesure de $A$) et $q$ (mesure de
$Q$) ; $X$ est appliqué si $a=1$, $Z$ si $q=1$. Le ket s'écrit $\ket{b,a,q}$ : le fil du haut
($Q$) est le dernier chiffre.}
\label{fig:teleportation}
\end{figure}

\subsubsection{Calcul pas à pas}

\paragraph{État initial $\ket{\Pi_0}$.} Alice et Bob partagent la paire, et $Q$ est indépendant
d'elle :
\begin{align*}
  \ket{\Pi_0}&=\ket{\Phi^+}_{BA}\otimes\ket{\varphi}_Q
    =\frac{1}{\sqrt2}\big(\ket{00}+\ket{11}\big)\otimes\big(\alpha\ket{0}+\beta\ket{1}\big)\\
  &=\frac{1}{\sqrt2}\Big(\alpha\ket{00}\ket{0}+\beta\ket{00}\ket{1}+\alpha\ket{11}\ket{0}+\beta\ket{11}\ket{1}\Big)\\
  &=\frac{\alpha\ket{000}+\beta\ket{001}+\alpha\ket{110}+\beta\ket{111}}{\sqrt2}.
\end{align*}

\paragraph{Après le CNOT $Q\to A$ : $\ket{\Pi_1}$.} Le CNOT (contrôle $Q$, dernier chiffre ;
cible $A$, chiffre du milieu) agit ainsi : $\ket{b,a,q}\to\ket{b,\,a\oplus q,\,q}$. Terme par terme :
\begin{center}
\renewcommand{\arraystretch}{1.3}
\begin{tabular}{@{}cccc@{}}
\toprule
Terme & $q$ & Effet sur $A$ & Résultat \\
\midrule
$\alpha\ket{000}$ & $0$ & aucun & $\alpha\ket{000}$ \\
$\beta\ket{001}$ & $1$ & $0\to1$ & $\beta\ket{011}$ \\
$\alpha\ket{110}$ & $0$ & aucun & $\alpha\ket{110}$ \\
$\beta\ket{111}$ & $1$ & $1\to0$ & $\beta\ket{101}$ \\
\bottomrule
\end{tabular}
\end{center}
donc
\[
  \ket{\Pi_1}=\frac{\alpha\ket{000}+\beta\ket{011}+\alpha\ket{110}+\beta\ket{101}}{\sqrt2}.
\]
On isole ensuite le dernier chiffre, celui de $Q$, qui va changer à l'étape suivante
(le premier ket est celui du couple $\ket{b,a}$) :
\[
  \ket{\Pi_1}=\frac{\alpha\ket{00}\otimes\ket{0}+\beta\ket{01}\otimes\ket{1}
    +\alpha\ket{11}\otimes\ket{0}+\beta\ket{10}\otimes\ket{1}}{\sqrt2}.
\]

\paragraph{Après la porte $H$ sur $Q$ : $\ket{\Pi_2}$.} On remplace $\ket{0}$ par
$H\ket{0}=\ket{+}$ et $\ket{1}$ par $H\ket{1}=\ket{-}$ dans le facteur $Q$ :
\[
  \ket{\Pi_2}=\frac{1}{\sqrt2}\Big(\alpha\ket{00}\otimes\ket{+}+\beta\ket{01}\otimes\ket{-}
    +\alpha\ket{11}\otimes\ket{+}+\beta\ket{10}\otimes\ket{-}\Big).
\]
Avec $\ket{\pm}=\frac{1}{\sqrt2}(\ket{0}\pm\ket{1})$, le préfacteur devient
$\frac{1}{\sqrt2}\cdot\frac{1}{\sqrt2}=\frac12$ :
\begin{align*}
  \ket{\Pi_2}&=\frac12\Big(\alpha\ket{00}\big(\ket{0}+\ket{1}\big)+\beta\ket{01}\big(\ket{0}-\ket{1}\big)
     +\alpha\ket{11}\big(\ket{0}+\ket{1}\big)+\beta\ket{10}\big(\ket{0}-\ket{1}\big)\Big)\\
  &=\frac12\Big(\alpha\ket{000}+\alpha\ket{001}+\beta\ket{010}-\beta\ket{011}
     +\alpha\ket{110}+\alpha\ket{111}+\beta\ket{100}-\beta\ket{101}\Big).
\end{align*}

\paragraph{Regroupement selon le résultat d'Alice.} Alice mesurera $A$ et $Q$, donc les deux
derniers chiffres. On regroupe les huit termes qui ont les mêmes chiffres $(a,q)$, en mettant
le ket de Bob en premier :
\begin{center}
\renewcommand{\arraystretch}{1.3}
\begin{tabular}{@{}cll@{}}
\toprule
$(a,q)$ & Termes de $\ket{\Pi_2}$ (facteur $\frac12$ omis) & Ket de Bob \\
\midrule
$(0,0)$ & $\alpha\ket{000}+\beta\ket{100}$ & $(\alpha\ket{0}+\beta\ket{1})\ket{00}$ \\
$(0,1)$ & $\alpha\ket{001}-\beta\ket{101}$ & $(\alpha\ket{0}-\beta\ket{1})\ket{01}$ \\
$(1,0)$ & $\beta\ket{010}+\alpha\ket{110}$ & $(\beta\ket{0}+\alpha\ket{1})\ket{10}$ \\
$(1,1)$ & $-\beta\ket{011}+\alpha\ket{111}$ & $(-\beta\ket{0}+\alpha\ket{1})\ket{11}$ \\
\bottomrule
\end{tabular}
\end{center}
d'où
\[
  \ket{\Pi_2}=\frac12\Big[\big(\alpha\ket{0}+\beta\ket{1}\big)\ket{00}
    +\big(\alpha\ket{0}-\beta\ket{1}\big)\ket{01}
    +\big(\beta\ket{0}+\alpha\ket{1}\big)\ket{10}
    +\big(\alpha\ket{1}-\beta\ket{0}\big)\ket{11}\Big].
\]
Dans chaque terme, le ket de gauche est l'état de $B$ et le ket de droite $\ket{a,q}$ celui du
couple $(A,Q)$ d'Alice.

\subsubsection{Mesure d'Alice et corrections de Bob}

Alice mesure $A$ et $Q$. Comme à la section~\ref{sec:mesure2}, la probabilité d'un résultat
$(a,q)$ est le carré de la norme du terme correspondant, et l'état de Bob devient le ket
associé (déjà normé). Ici
\[
  P(a,q)=\Big\|\tfrac12\ket{\phi_{aq}}\Big\|^2=\tfrac14\big(|\alpha|^2+|\beta|^2\big)=\tfrac14
  \quad\text{pour les quatre résultats,}
\]
qui sont donc équiprobables et \emph{indépendants} de $\alpha$ et $\beta$ : les deux bits
envoyés ne révèlent rien sur $\ket{\varphi}$. Les kets de Bob sont
$\ket{\phi_{00}}=\alpha\ket{0}+\beta\ket{1}$, $\ket{\phi_{01}}=\alpha\ket{0}-\beta\ket{1}$,
$\ket{\phi_{10}}=\beta\ket{0}+\alpha\ket{1}$ et $\ket{\phi_{11}}=-\beta\ket{0}+\alpha\ket{1}$.

\begin{center}
\renewcommand{\arraystretch}{1.5}
\begin{tabular}{@{}cccll@{}}
\toprule
Résultat $(a,q)$ & Probabilité & État de Bob & Lien avec $\ket{\varphi}$ & Correction de Bob \\
\midrule
$(0,0)$ & $\frac14$ & $\alpha\ket{0}+\beta\ket{1}$ & $\ket{\varphi}$ & $I$ (rien) \\
$(0,1)$ & $\frac14$ & $\alpha\ket{0}-\beta\ket{1}$ & $Z\ket{\varphi}$ & $Z$ \\
$(1,0)$ & $\frac14$ & $\beta\ket{0}+\alpha\ket{1}$ & $X\ket{\varphi}$ & $X$ \\
$(1,1)$ & $\frac14$ & $-\beta\ket{0}+\alpha\ket{1}$ & $XZ\ket{\varphi}$ & $X$ puis $Z$ (opérateur $ZX$) \\
\bottomrule
\end{tabular}
\end{center}

\paragraph{Vérification des quatre corrections.} Avec $X\ket{0}=\ket{1}$, $X\ket{1}=\ket{0}$,
$Z\ket{0}=\ket{0}$ et $Z\ket{1}=-\ket{1}$ :
\begin{align*}
  (0,0):\quad & I\big(\alpha\ket{0}+\beta\ket{1}\big)=\alpha\ket{0}+\beta\ket{1},\\
  (0,1):\quad & Z\big(\alpha\ket{0}-\beta\ket{1}\big)=\alpha\ket{0}-\beta Z\ket{1}=\alpha\ket{0}+\beta\ket{1},\\
  (1,0):\quad & X\big(\beta\ket{0}+\alpha\ket{1}\big)=\beta\ket{1}+\alpha\ket{0}=\alpha\ket{0}+\beta\ket{1},\\
  (1,1):\quad & X\big(-\beta\ket{0}+\alpha\ket{1}\big)=-\beta\ket{1}+\alpha\ket{0}=\alpha\ket{0}-\beta\ket{1},
     \ \text{puis}\ Z\big(\alpha\ket{0}-\beta\ket{1}\big)=\alpha\ket{0}+\beta\ket{1}.
\end{align*}
Dans le dernier cas, avec les matrices,
$ZX=\begin{pmatrix}1&0\\0&-1\end{pmatrix}\begin{pmatrix}0&1\\1&0\end{pmatrix}
=\begin{pmatrix}0&1\\-1&0\end{pmatrix}$ et
$ZX\begin{pmatrix}-\beta\\ \alpha\end{pmatrix}=\begin{pmatrix}\alpha\\ \beta\end{pmatrix}$.
L'ordre compte : l'état de Bob est $XZ\ket{\varphi}$ ($Z\ket{\varphi}=\alpha\ket{0}-\beta\ket{1}$,
puis $X$ donne $\alpha\ket{1}-\beta\ket{0}$), et on l'annule avec
$(XZ)^{-1}=Z^{-1}X^{-1}=ZX$ : on applique d'abord $X$, puis $Z$, comme dans le circuit.
Dans les quatre cas, Bob retrouve exactement $\ket{\varphi}=\alpha\ket{0}+\beta\ket{1}$.

\subsubsection{Ce que Bob sait avant de recevoir les bits}

Tant que Bob ignore $(a,q)$, son qubit est dans le mélange des quatre kets $\ket{\phi_{aq}}$, chacun
avec la probabilité $\frac14$ :
\begin{align*}
  \rho_B&=\frac14\sum_{a,q}\ket{\phi_{aq}}\bra{\phi_{aq}}\\
  &=\frac14\Bigg[\begin{pmatrix}|\alpha|^2&\alpha\beta^*\\ \alpha^*\beta&|\beta|^2\end{pmatrix}
    +\begin{pmatrix}|\alpha|^2&-\alpha\beta^*\\ -\alpha^*\beta&|\beta|^2\end{pmatrix}
    +\begin{pmatrix}|\beta|^2&\alpha^*\beta\\ \alpha\beta^*&|\alpha|^2\end{pmatrix}
    +\begin{pmatrix}|\beta|^2&-\alpha^*\beta\\ -\alpha\beta^*&|\alpha|^2\end{pmatrix}\Bigg]\\
  &=\frac14\begin{pmatrix}2(|\alpha|^2+|\beta|^2)&0\\0&2(|\alpha|^2+|\beta|^2)\end{pmatrix}
   =\frac12\begin{pmatrix}1&0\\0&1\end{pmatrix}=\frac12\,\mathrm{Id}.
\end{align*}
Sans les deux bits, Bob a donc un qubit complètement mélangé, qui ne dépend pas de $\alpha$ ni de
$\beta$ : il n'a reçu aucune information sur $\ket{\varphi}$. C'est l'envoi classique de $(a,q)$
qui permet de la retrouver, donc rien ne voyage plus vite que la lumière.

\begin{remarque}[Bilan de la téléportation]
\begin{itemize}[nosep]
  \item Ressources : $1$ e-bit (la paire de Bell) et $2$ bits classiques pour transmettre
        \emph{un qubit} d'état quelconque, sans qu'Alice connaisse $\alpha$ ni $\beta$.
  \item Pas de clonage : après la mesure, $A$ et $Q$ sont dans l'état de base $\ket{a,q}$.
        Alice n'a plus $\ket{\varphi}$ : l'état a été \emph{déplacé}, pas copié. L'intrication
        entre $A$ et $B$ a été consommée.
  \item Les quatre corrections $I,Z,X,ZX$ sont des portes contrôlées par des bits classiques
        (section~\ref{sec:controlees}).
\end{itemize}
\end{remarque}

% ---------------------------------------------------------------------
\subsection{Portes contrôlées : contrôlé-\texorpdfstring{$U$}{U}, Fredkin et Toffoli}\label{sec:controlees}
% ---------------------------------------------------------------------

Le CNOT est le cas le plus simple d'une idée générale : appliquer une opération à un ou
plusieurs qubits \emph{cibles} seulement si un ou plusieurs qubits de \emph{contrôle} valent $1$.
Dans cette section, comme à la section~\ref{sec:op2}, le contrôle est le \emph{premier} chiffre
du ket. Les qubits sont notés $x$, $y$, $z$ (le qubit $x$ est le contrôle) et
$\ket{x,y,z}=\ket{x}\otimes\ket{y}\otimes\ket{z}$ ; d'après la convention de la section, le
contrôle est dessiné en bas (figure~\ref{fig:controlees}). Pour éviter la confusion avec les
matrices de Pauli, on écrit ici $\sigma_x=X$ pour la porte NOT.

\begin{figure}[!ht]
\centering
\begin{adjustbox}{max width=\linewidth}
\begin{tikzpicture}[>=Stealth, font=\small]
  % (a) contrôlé-U : contrôle x en bas, cible y en haut
  \draw (0,1.2) -- (3.4,1.2);
  \draw (0,0) -- (3.4,0);
  \node[anchor=east] at (0,1.2) {$y$};
  \node[anchor=east] at (0,0) {$x$};
  \draw[insarouge, thick] (1.7,0) -- (1.7,0.85);
  \pic at (1.7,0) {ctl};
  \node[porte] at (1.7,1.2) {$U$};
  \node[anchor=north, font=\footnotesize] at (1.7,-0.5) {(a) contrôlé-$U$};
  % (b) contrôlé-SWAP (Fredkin) : x contrôle, y et z échangés
  \draw (5,2.4) -- (8.4,2.4);
  \draw (5,1.2) -- (8.4,1.2);
  \draw (5,0) -- (8.4,0);
  \node[anchor=east] at (5,2.4) {$z$};
  \node[anchor=east] at (5,1.2) {$y$};
  \node[anchor=east] at (5,0) {$x$};
  \draw[insarouge, thick] (6.7,0) -- (6.7,2.4);
  \pic at (6.7,0) {ctl};
  \pic at (6.7,1.2) {croix};
  \pic at (6.7,2.4) {croix};
  \node[anchor=north, font=\footnotesize] at (6.7,-0.5) {(b) Fredkin (C-SWAP)};
  % (c) Toffoli (CCNOT) : x et y contrôlent, z cible
  \draw (10,2.4) -- (13.4,2.4);
  \draw (10,1.2) -- (13.4,1.2);
  \draw (10,0) -- (13.4,0);
  \node[anchor=east] at (10,2.4) {$z$};
  \node[anchor=east] at (10,1.2) {$y$};
  \node[anchor=east] at (10,0) {$x$};
  \draw[insarouge, thick] (11.7,0) -- (11.7,2.4);
  \pic at (11.7,0) {ctl};
  \pic at (11.7,1.2) {ctl};
  \pic at (11.7,2.4) {cible};
  \node[anchor=north, font=\footnotesize] at (11.7,-0.5) {(c) Toffoli (CCNOT)};
\end{tikzpicture}
\end{adjustbox}
\caption{Portes contrôlées. Le contrôle $x$ est le premier chiffre du ket, donc le fil du bas.
(a) Le qubit $y$ subit $U$ si $x=1$ ; pour $U=X$ on retrouve le CNOT. (b) Si $x=1$, les qubits
$y$ et $z$ sont échangés. (c) Si $x=y=1$, le qubit $z$ est inversé.}
\label{fig:controlees}
\end{figure}

\subsubsection{Contrôlé-\texorpdfstring{$U$}{U}}

\begin{definition}{Porte contrôlée-$U$}
Soit $U$ une matrice unitaire $2\times2$. Sur le système $(x,y)$, de contrôle $x$ et de cible $y$,
\[
  \mathrm{C}U=\ket{0}\bra{0}_x\otimes I_y+\ket{1}\bra{1}_x\otimes U_y .
\]
Le premier terme agit quand le contrôle est $\ket{0}$ : la cible ne change pas ($I$). Le second
agit quand le contrôle est $\ket{1}$ : la cible subit $U$.
\end{definition}

Comme $\ket{0}\bra{0}\ket{x}=\delta_{x0}\ket{0}$ et $\ket{1}\bra{1}\ket{x}=\delta_{x1}\ket{1}$,
\[
  \mathrm{C}U\big(\ket{x}\otimes\ket{y}\big)=\ket{x}\otimes U^x\ket{y},\qquad U^0=I,\ U^1=U .
\]
Par la règle $A\otimes B=(A_{ij}B)$, $\ket{0}\bra{0}\otimes I=\begin{pmatrix}I&0\\0&0\end{pmatrix}$
et $\ket{1}\bra{1}\otimes U=\begin{pmatrix}0&0\\0&U\end{pmatrix}$ (blocs $2\times2$), donc
\[
  \mathrm{C}U=\begin{pmatrix}I&0\\0&U\end{pmatrix}
  =\begin{pmatrix}1&0&0&0\\0&1&0&0\\0&0&u_{00}&u_{01}\\0&0&u_{10}&u_{11}\end{pmatrix}
  \qquad\text{pour }U=\begin{pmatrix}u_{00}&u_{01}\\u_{10}&u_{11}\end{pmatrix}.
\]
Cette matrice est unitaire car $I^\dagger I=I$ et $U^\dagger U=I$.

\paragraph{Le CNOT est un cas particulier.} Pour $U=\sigma_x$,
$\mathrm{C}\sigma_x=\ket{0}\bra{0}_x\otimes I_y+\ket{1}\bra{1}_x\otimes\sigma_x$ est le CNOT de
contrôle $x$ (premier chiffre), c'est-à-dire $\mathrm{CNOT}_{12}$ de la section~\ref{sec:cnot}.
On calcule l'action sur les quatre états de base, avec $\braket{0}{0}=\braket{1}{1}=1$,
$\braket{0}{1}=\braket{1}{0}=0$, $\sigma_x\ket{0}=\ket{1}$ et $\sigma_x\ket{1}=\ket{0}$ :
\begin{align*}
  \mathrm{C}\sigma_x\ket{00}&=\big[\ket{0}\bra{0}_x\otimes I_y\big]\ket{0}_x\ket{0}_y
      +\big[\ket{1}\bra{1}_x\otimes\sigma_x\big]\ket{0}_x\ket{0}_y\\
    &=\ket{0}\braket{0}{0}\otimes I\ket{0}+\ket{1}\braket{1}{0}\otimes\sigma_x\ket{0}
     =\ket{0}\otimes\ket{0}+0=\ket{00},\\[1.2ex]
  \mathrm{C}\sigma_x\ket{01}&=\ket{0}\braket{0}{0}\otimes I\ket{1}+\ket{1}\braket{1}{0}\otimes\sigma_x\ket{1}
     =\ket{0}\otimes\ket{1}+0=\ket{01},\\[1.2ex]
  \mathrm{C}\sigma_x\ket{10}&=\ket{0}\braket{0}{1}\otimes I\ket{0}+\ket{1}\braket{1}{1}\otimes\sigma_x\ket{0}
     =0+\ket{1}\otimes\ket{1}=\ket{11},\\[1.2ex]
  \mathrm{C}\sigma_x\ket{11}&=\ket{0}\braket{0}{1}\otimes I\ket{1}+\ket{1}\braket{1}{1}\otimes\sigma_x\ket{1}
     =0+\ket{1}\otimes\ket{0}=\ket{10}.
\end{align*}
On retrouve la table du CNOT : la cible n'est inversée que si le contrôle vaut $1$.

\paragraph{Un autre exemple : $U=Z$.} Ici
\[
  \mathrm{C}Z=\begin{pmatrix}I&0\\0&Z\end{pmatrix}=\mathrm{diag}(1,1,1,-1) :
\]
la porte multiplie $\ket{11}$ par $-1$ et ne touche pas les autres états de base. Elle est
symétrique entre les deux qubits. Elle se déduit du CNOT par
$\mathrm{C}Z=(I\otimes H)\,\mathrm{C}X\,(I\otimes H)$, où $\mathrm{C}X$ est le contrôlé-$X$.
En effet, $I\otimes H=\mathrm{diag}(H,H)$ (deux blocs $2\times2$ égaux à $H$), et le produit de
matrices diagonales par blocs se fait bloc par bloc :
\[
  (I\otimes H)\begin{pmatrix}I&0\\0&X\end{pmatrix}(I\otimes H)
  =\begin{pmatrix}H\,I\,H&0\\0&H\,X\,H\end{pmatrix}
  =\begin{pmatrix}I&0\\0&HXH\end{pmatrix},
\]
car $H^2=I$, et
\[
  HXH=\frac12\begin{pmatrix}1&1\\1&-1\end{pmatrix}\begin{pmatrix}0&1\\1&0\end{pmatrix}\begin{pmatrix}1&1\\1&-1\end{pmatrix}
  =\frac12\begin{pmatrix}1&1\\1&-1\end{pmatrix}\begin{pmatrix}1&-1\\1&1\end{pmatrix}
  =\frac12\begin{pmatrix}2&0\\0&-2\end{pmatrix}=Z
\]
(section~\ref{sec:hadamard} : $H$ échange les axes $x$ et $z$).

\subsubsection{Contrôlé-SWAP : la porte de Fredkin}

\begin{definition}{Porte de Fredkin (C-SWAP)}
Sur trois qubits $(x,y,z)$ (espace de dimension $2^3=8$), avec $x$ pour contrôle,
\[
  \mathrm{C\text{-}SWAP}=\ket{0}\bra{0}_x\otimes I_y\otimes I_z+\ket{1}\bra{1}_x\otimes\mathrm{SWAP}_{yz}.
\]
Si $x=0$, rien ne change ; si $x=1$, les qubits $y$ et $z$ sont échangés.
\end{definition}

Calcul sur quatre états de base (avec $I\ket{ab}=\ket{ab}$ et $\mathrm{SWAP}\ket{ab}=\ket{ba}$) :
\begin{align*}
  \mathrm{C\text{-}SWAP}\ket{000}&=\ket{0}\braket{0}{0}\otimes I\ket{00}+\ket{1}\braket{1}{0}\otimes\mathrm{SWAP}\ket{00}
     =\ket{0}\otimes\ket{00}+0=\ket{000},\\[1.2ex]
  \mathrm{C\text{-}SWAP}\ket{101}&=\ket{0}\braket{0}{1}\otimes I\ket{01}+\ket{1}\braket{1}{1}\otimes\mathrm{SWAP}\ket{01}
     =0+\ket{1}\otimes\ket{10}=\ket{110},\\[1.2ex]
  \mathrm{C\text{-}SWAP}\ket{110}&=\ket{0}\braket{0}{1}\otimes I\ket{10}+\ket{1}\braket{1}{1}\otimes\mathrm{SWAP}\ket{10}
     =0+\ket{1}\otimes\ket{01}=\ket{101},\\[1.2ex]
  \mathrm{C\text{-}SWAP}\ket{010}&=\ket{0}\braket{0}{0}\otimes I\ket{10}+\ket{1}\braket{1}{0}\otimes\mathrm{SWAP}\ket{10}
     =\ket{0}\otimes\ket{10}+0=\ket{010}.
\end{align*}
Les huit états de base donnent :
\begin{center}
\renewcommand{\arraystretch}{1.3}
\begin{tabular}{@{}ccccccccc@{}}
\toprule
Entrée & $\ket{000}$ & $\ket{001}$ & $\ket{010}$ & $\ket{011}$ & $\ket{100}$ & $\ket{101}$ & $\ket{110}$ & $\ket{111}$ \\
\midrule
Sortie & $\ket{000}$ & $\ket{001}$ & $\ket{010}$ & $\ket{011}$ & $\ket{100}$ & $\ket{110}$ & $\ket{101}$ & $\ket{111}$ \\
\bottomrule
\end{tabular}
\end{center}
Seuls $\ket{101}$ et $\ket{110}$ changent (contrôle à $1$ et $y\neq z$) ; pour $\ket{100}$ et
$\ket{111}$ l'échange de deux qubits égaux ne change rien. Dans la base
$\ket{000},\ket{001},\dots,\ket{111}$ (ordre binaire), la matrice est
$\mathrm{diag}(I_4,\mathrm{SWAP})$ (deux blocs $4\times4$) :
\[
  \mathrm{C\text{-}SWAP}=\begingroup\setlength{\arraycolsep}{3pt}\begin{pmatrix}
    1&0&0&0&0&0&0&0\\
    0&1&0&0&0&0&0&0\\
    0&0&1&0&0&0&0&0\\
    0&0&0&1&0&0&0&0\\
    0&0&0&0&1&0&0&0\\
    0&0&0&0&0&0&1&0\\
    0&0&0&0&0&1&0&0\\
    0&0&0&0&0&0&0&1
  \end{pmatrix}\endgroup .
\]
C'est l'identité dont les lignes $6$ et $7$ sont échangées.

\paragraph{Contrôle en superposition.} Avec $\ket{+}_x\otimes\ket{0}_y\otimes\ket{1}_z
=\frac{1}{\sqrt2}\big(\ket{001}+\ket{101}\big)$ :
\[
  \mathrm{C\text{-}SWAP}\,\frac{1}{\sqrt2}\big(\ket{001}+\ket{101}\big)
  =\frac{1}{\sqrt2}\big(\ket{001}+\ket{110}\big).
\]
Le résultat n'est pas un produit : les coefficients de $\ket{x}\otimes\ket{yz}$ forment la
matrice (lignes $x=0,1$, colonnes $yz=00,01,10,11$)
\[
  \frac{1}{\sqrt2}\begin{pmatrix}0&1&0&0\\0&0&1&0\end{pmatrix},
\]
dont les deux lignes ne sont pas proportionnelles. Le qubit $x$ est intriqué avec la paire
$(y,z)$.

\subsubsection{Contrôlé-contrôlé-NOT : la porte de Toffoli}

\begin{definition}{Porte de Toffoli (CCNOT)}
Sur trois qubits $(x,y,z)$, avec $x$ et $y$ pour contrôles et $z$ pour cible,
\[
  \mathrm{CCNOT}=\ket{0}\bra{0}_x\otimes I_y\otimes I_z
   +\ket{1}\bra{1}_x\otimes\Big[\ket{0}\bra{0}_y\otimes I_z+\ket{1}\bra{1}_y\otimes\sigma_x\Big].
\]
Le crochet est un CNOT (contrôle $y$, cible $z$), appliqué seulement si $x=1$ : si $x=0$, rien ;
si $x=1$ et $y=0$, rien ; si $x=1$ et $y=1$, le qubit $z$ est inversé.
\end{definition}

Calcul de $\mathrm{CCNOT}\ket{110}$ :
\begin{align*}
  \mathrm{CCNOT}\ket{110}&=\ket{0}\braket{0}{1}\otimes I\ket{1}\otimes I\ket{0}
     +\ket{1}\braket{1}{1}\otimes\Big[\ket{0}\braket{0}{1}\otimes I\ket{0}
       +\ket{1}\braket{1}{1}\otimes\sigma_x\ket{0}\Big]\\
  &=0+\ket{1}\otimes\Big[0+\ket{1}\otimes\ket{1}\Big]=\ket{111}.
\end{align*}
Pour les huit états de base, $\mathrm{CCNOT}\ket{x,y,z}=\ket{x,\,y,\,z\oplus xy}$ :
\begin{center}
\renewcommand{\arraystretch}{1.3}
\begin{tabular}{@{}ccccccccc@{}}
\toprule
Entrée & $\ket{000}$ & $\ket{001}$ & $\ket{010}$ & $\ket{011}$ & $\ket{100}$ & $\ket{101}$ & $\ket{110}$ & $\ket{111}$ \\
\midrule
Sortie & $\ket{000}$ & $\ket{001}$ & $\ket{010}$ & $\ket{011}$ & $\ket{100}$ & $\ket{101}$ & $\ket{111}$ & $\ket{110}$ \\
\bottomrule
\end{tabular}
\end{center}
Seuls $\ket{110}$ et $\ket{111}$ sont échangés : la matrice est $\mathrm{diag}(I_6,X)$
(un bloc $6\times6$ égal à l'identité, puis $X$),
\[
  \mathrm{CCNOT}=\begingroup\setlength{\arraycolsep}{3pt}\begin{pmatrix}
    1&0&0&0&0&0&0&0\\
    0&1&0&0&0&0&0&0\\
    0&0&1&0&0&0&0&0\\
    0&0&0&1&0&0&0&0\\
    0&0&0&0&1&0&0&0\\
    0&0&0&0&0&1&0&0\\
    0&0&0&0&0&0&0&1\\
    0&0&0&0&0&0&1&0
  \end{pmatrix}\endgroup .
\]

\begin{remarque}[Portes réversibles]
Les portes de Fredkin et de Toffoli sont des permutations des états de base qui échangent
seulement deux d'entre eux : elles sont leurs propres inverses ($\mathrm{C\text{-}SWAP}^2=\mathrm{CCNOT}^2=\mathrm{Id}$)
et unitaires. Avec $z=0$, la Toffoli calcule le ET logique $x\wedge y$ sans rien effacer :
$\ket{x,y,0}\to\ket{x,y,xy}$ ; avec $z=1$, elle calcule le NON-ET. Elle sert à écrire tout calcul
classique de façon réversible.
\end{remarque}

% ---------------------------------------------------------------------
\subsection{Récapitulatif du chapitre}
% ---------------------------------------------------------------------

\begin{center}
\renewcommand{\arraystretch}{1.5}
\begin{tabular}{@{}lp{12cm}@{}}
\toprule
\textbf{Notion} & \textbf{Formule ou résultat} \\
\midrule
CNOT & $\ket{b,a}\to\ket{b\oplus a,a}$ (contrôle $a$ écrit en dernier), $\mathrm{CNOT}^2=\mathrm{Id}$ \\
SWAP & $\mathrm{SWAP}\ket{a,b}=\ket{b,a}$, \quad $\mathrm{CNOT}_{21}=\mathrm{SWAP}\,\mathrm{CNOT}_{12}\,\mathrm{SWAP}$ \\
Circuit de Bell & $(I\otimes H)$ puis CNOT : $\ket{\beta^{xy}}=\frac{1}{\sqrt2}\big(\ket{y,0}+(-1)^x\ket{\bar y,1}\big)$, $\ket{\beta^{00}}=\ket{\Phi^+}$ \\
Information & $I(A;B)=H(A)+H(B)-H(A,B)$, \quad $H(A|B)=H(A,B)-H(B)$ \\
Trace partielle & $\operatorname{Tr}_B\big(\ket{a}\bra{a'}\otimes\ket{b}\bra{b'}\big)=\ket{a}\bra{a'}\braket{b'}{b}$ ; état de Bell : $\rho_A=\frac12\mathrm{Id}$ \\
État de Bell (bits) & $H(A,B)=0$, \quad $H(A)=H(B)=1$, \quad $I(A;B)=2$, \quad $H(A|B)=-1$ \\
Pur / mélangé & $\operatorname{Tr}\rho^2=1$ / $<1$ ; $\ket{+}$ : $1$ ; qubit d'une paire de Bell : $\frac12$ ; source de Bob : $\frac58$ \\
Téléportation & $1$ e-bit $+\ 2$ bits classiques ; correction : $X$ si $a=1$, puis $Z$ si $q=1$ ; $P(a,q)=\frac14$ \\
Contrôlé-$U$ & $\ket{0}\bra{0}\otimes I+\ket{1}\bra{1}\otimes U=\mathrm{diag}(I,U)$ ; $\ket{x,y}\to\ket{x}\otimes U^x\ket{y}$ \\
Fredkin & $\ket{0}\bra{0}\otimes I\otimes I+\ket{1}\bra{1}\otimes\mathrm{SWAP}$ ; $\ket{101}\leftrightarrow\ket{110}$ \\
Toffoli & $\ket{x,y,z}\to\ket{x,y,z\oplus xy}$ ; $\ket{110}\leftrightarrow\ket{111}$ \\
\bottomrule
\end{tabular}
\end{center}
.
%  Comme chapitre4.tex, chapitre5.tex et chapitre6.tex, il ouvre une NOUVELLE
%  SECTION (numéro 5). Pour le rattacher à la section précédente, remplacer
%  \section par \subsection et chaque \subsection par \subsubsection.
%  Pas de préambule : environnements (definition, resultat, remarque),
%  macros (\ii, \ee, \dd, \ket, \bra, \braket) et couleurs (insarouge,
%  insagris) sont définis dans main.tex. Les symboles de circuits (contrôle,
%  cible, croix, mesure) sont définis juste après avec \tikzset.
% =====================================================================

% Macros de Dirac (déjà définies dans main.tex récent ; ce bloc évite l'erreur
% « Undefined control sequence \ket » si main.tex est une ancienne version).
\providecommand{\ket}[1]{|#1\rangle}
\providecommand{\bra}[1]{\langle #1|}
\providecommand{\braket}[2]{\langle #1|#2\rangle}

% Symboles de circuits pour les figures de ce chapitre.
\tikzset{
  ctl/.pic   = {\fill[insarouge] (0,0) circle (2.6pt);},
  cible/.pic = {\draw[insarouge, thick, fill=white] (0,0) circle (0.28);
                \draw[insarouge, thick] (-0.28,0) -- (0.28,0);
                \draw[insarouge, thick] (0,-0.28) -- (0,0.28);},
  croix/.pic = {\draw[insarouge, thick] (-0.17,-0.17) -- (0.17,0.17);
                \draw[insarouge, thick] (-0.17,0.17) -- (0.17,-0.17);},
  mesure/.pic = {\draw[thick, fill=white] (-0.3,-0.27) rectangle (0.3,0.27);
                 \draw[thick] (-0.19,-0.13) arc[start angle=155, end angle=25, radius=0.21];
                 \draw[thick, -{Stealth[length=1.3mm]}] (0,-0.17) -- (0.14,0.12);},
  porte/.style = {draw=insarouge, fill=insarouge!8, minimum size=7mm, thick},
}

% =====================================================================
\section{Portes multi-qubits, circuits et téléportation}\label{sec:circuits}
% =====================================================================

La section précédente a décrit les états de deux qubits. On étudie ici les portes qui
agissent sur plusieurs qubits : le CNOT et le SWAP, puis les portes contrôlées à deux et à
trois qubits (contrôlé-$U$, Fredkin, Toffoli). On assemble ensuite des portes dans un
circuit qui fabrique un état de Bell, on mesure l'information contenue dans un tel état,
on distingue les états purs des états mélangés, et on démontre le protocole de
téléportation. Tous les calculs sont détaillés.

\begin{remarque}[Convention pour les circuits]
Dans un circuit, le temps va de gauche à droite et chaque qubit a un fil. Les fils sont
dessinés dans l'ordre \emph{inverse} de l'écriture des kets : le fil du \emph{haut} porte le
\emph{dernier} chiffre du ket, le fil du \emph{bas} le premier. Par exemple, si le qubit $A$
est dessiné au-dessus du qubit $B$, le ket s'écrit $\ket{b,a}=\ket{b}_B\otimes\ket{a}_A$.
La figure~\ref{fig:bell} de la section précédente suit la même convention.
\end{remarque}

% ---------------------------------------------------------------------
\subsection{Porte CNOT}\label{sec:cnot}
% ---------------------------------------------------------------------

\begin{definition}{Porte CNOT (NOT contrôlé)}
La porte CNOT agit sur deux qubits : le qubit $A$ est le \emph{contrôle}, le qubit $B$ est la
\emph{cible}. La cible est inversée (porte $X$) si le contrôle vaut $1$ et reste inchangée
s'il vaut $0$. Avec le système écrit $(B,A)$, soit $\ket{b,a}=\ket{b}_B\otimes\ket{a}_A$ :
\[
  \mathrm{CNOT}\,\ket{b,a}=\ket{b\oplus a,\,a},
  \qquad
  \mathrm{CNOT}=I\otimes\ket{0}\bra{0}+X\otimes\ket{1}\bra{1},
\]
où $\oplus$ est l'addition modulo $2$ ($0\oplus0=1\oplus1=0$ et $0\oplus1=1\oplus0=1$).
\end{definition}

Sur le circuit de la figure~\ref{fig:cnot}, le point noir posé sur le fil de $A$ marque le
contrôle et le symbole $\oplus$ posé sur le fil de $B$ marque la cible.

\begin{figure}[!ht]
\centering
\begin{tikzpicture}[>=Stealth, font=\small]
  \draw (0,1.2) -- (6,1.2);
  \draw (0,0) -- (6,0);
  \node[anchor=east] at (0,1.2) {$\ket{a}_A$};
  \node[anchor=east] at (0,0) {$\ket{b}_B$};
  \node[anchor=west] at (6,1.2) {$\ket{a}_A$};
  \node[anchor=west] at (6,0) {$\ket{b\oplus a}_B$};
  \draw[insarouge, thick] (3,1.2) -- (3,0);
  \pic at (3,1.2) {ctl};
  \pic at (3,0) {cible};
  \node[anchor=south, text=insagris] at (3,1.5) {qubit de contrôle};
  \node[anchor=north, text=insagris] at (3,-0.5) {qubit cible};
\end{tikzpicture}
\caption{Porte CNOT : le fil du haut ($A$) est le contrôle, le fil du bas ($B$) est la
cible. Le ket s'écrit $\ket{b,a}$ (le fil du haut donne le dernier chiffre).}
\label{fig:cnot}
\end{figure}

\subsubsection{Table de vérité}

\begin{center}
\renewcommand{\arraystretch}{1.3}
\begin{tabular}{@{}cccc@{}}
\toprule
Entrée $\ket{b,a}$ & Contrôle $a$ & Sortie $\ket{b\oplus a,a}$ & Effet sur la cible $B$ \\
\midrule
$\ket{00}$ & $0$ & $\ket{00}$ & aucun \\
$\ket{01}$ & $1$ & $\ket{11}$ & $0\to1$ \\
$\ket{10}$ & $0$ & $\ket{10}$ & aucun \\
$\ket{11}$ & $1$ & $\ket{01}$ & $1\to0$ \\
\bottomrule
\end{tabular}
\end{center}

\subsubsection{Matrice du CNOT}

La $j$-ième colonne d'une matrice $4\times4$ est l'image du $j$-ième vecteur de base
($\ket{00},\ket{01},\ket{10},\ket{11}$, dans cet ordre). La table donne :
\begin{align*}
  \mathrm{CNOT}\ket{00}&=\ket{00}=\begin{pmatrix}1\\0\\0\\0\end{pmatrix},&
  \mathrm{CNOT}\ket{01}&=\ket{11}=\begin{pmatrix}0\\0\\0\\1\end{pmatrix},\\[1.5ex]
  \mathrm{CNOT}\ket{10}&=\ket{10}=\begin{pmatrix}0\\0\\1\\0\end{pmatrix},&
  \mathrm{CNOT}\ket{11}&=\ket{01}=\begin{pmatrix}0\\1\\0\\0\end{pmatrix},
\end{align*}
d'où, en juxtaposant ces quatre colonnes,
\[
  \mathrm{CNOT}=\begin{pmatrix}
    1 & 0 & 0 & 0\\ 0 & 0 & 0 & 1\\ 0 & 0 & 1 & 0\\ 0 & 1 & 0 & 0
  \end{pmatrix}.
\]

On retrouve cette matrice avec la formule $I\otimes\ket{0}\bra{0}+X\otimes\ket{1}\bra{1}$ et
la règle $A\otimes B=(A_{ij}B)$ (section~\ref{sec:op2}), où
$\ket{0}\bra{0}=\begin{pmatrix}1&0\\0&0\end{pmatrix}$ et
$\ket{1}\bra{1}=\begin{pmatrix}0&0\\0&1\end{pmatrix}$ :
\[
  I\otimes\ket{0}\bra{0}=\begin{pmatrix}1&0&0&0\\0&0&0&0\\0&0&1&0\\0&0&0&0\end{pmatrix},
  \qquad
  X\otimes\ket{1}\bra{1}=\begin{pmatrix}0&0&0&0\\0&0&0&1\\0&0&0&0\\0&1&0&0\end{pmatrix},
\]
et la somme de ces deux matrices est la matrice du CNOT.

\subsubsection{Action sur un état quelconque}

Pour $\ket{\Psi}=c_{00}\ket{00}+c_{01}\ket{01}+c_{10}\ket{10}+c_{11}\ket{11}$ :
\[
  \mathrm{CNOT}\begin{pmatrix}c_{00}\\c_{01}\\c_{10}\\c_{11}\end{pmatrix}
  =\begin{pmatrix}
    1\cdot c_{00}+0\cdot c_{01}+0\cdot c_{10}+0\cdot c_{11}\\
    0\cdot c_{00}+0\cdot c_{01}+0\cdot c_{10}+1\cdot c_{11}\\
    0\cdot c_{00}+0\cdot c_{01}+1\cdot c_{10}+0\cdot c_{11}\\
    0\cdot c_{00}+1\cdot c_{01}+0\cdot c_{10}+0\cdot c_{11}
  \end{pmatrix}
  =\begin{pmatrix}c_{00}\\c_{11}\\c_{10}\\c_{01}\end{pmatrix}.
\]
Le CNOT échange les amplitudes de $\ket{01}$ et de $\ket{11}$, c'est-à-dire des deux états
dont le contrôle vaut $1$, et laisse les autres.

\paragraph{Exemple : contrôle en superposition.} Soit $\ket{0}_B\otimes\ket{+}_A
=\frac{1}{\sqrt2}\big(\ket{00}+\ket{01}\big)=\frac{1}{\sqrt2}(1,1,0,0)^{\mathsf T}$,
un état séparable. Alors
\[
  \mathrm{CNOT}\,\frac{1}{\sqrt2}\begin{pmatrix}1\\1\\0\\0\end{pmatrix}
  =\frac{1}{\sqrt2}\begin{pmatrix}1\\0\\0\\1\end{pmatrix}
  =\frac{1}{\sqrt2}\big(\ket{00}+\ket{11}\big)=\ket{\Phi^+}.
\]
Le résultat est intriqué ($D=c_{00}c_{11}-c_{01}c_{10}=\frac12\neq0$ ; section~\ref{sec:separable}) :
un CNOT dont le contrôle est en superposition crée de l'intrication.

\paragraph{Réversibilité.} Deux CNOT successifs redonnent l'état initial :
$\ket{b,a}\to\ket{b\oplus a,a}\to\ket{b\oplus a\oplus a,a}=\ket{b,a}$, donc
$\mathrm{CNOT}^2=\mathrm{Id}$. La matrice est réelle et symétrique, donc
$\mathrm{CNOT}^\dagger\mathrm{CNOT}=\mathrm{CNOT}^2=\mathrm{Id}$ : elle est bien unitaire.

\begin{remarque}[Contrôle en dernier ou en premier]
La matrice dépend de l'ordre d'écriture du ket. Notons $\mathrm{CNOT}_{21}$ le CNOT dont le
contrôle est le \emph{second} chiffre du ket (c'est la matrice ci-dessus) et
$\mathrm{CNOT}_{12}$ celui dont le contrôle est le \emph{premier} chiffre, utilisé à la
section précédente :
\[
  \mathrm{CNOT}_{12}=\ket{0}\bra{0}\otimes I+\ket{1}\bra{1}\otimes X
  =\begin{pmatrix}1&0&0&0\\0&1&0&0\\0&0&0&1\\0&0&1&0\end{pmatrix}.
\]
Un même circuit (contrôle $A$, cible $B$) a donc deux matrices, selon que l'on écrit
$A$ en dernier ($\mathrm{CNOT}_{21}$) ou en premier ($\mathrm{CNOT}_{12}$). Les deux sont
reliées par le SWAP (section suivante).
\end{remarque}

% ---------------------------------------------------------------------
\subsection{Porte SWAP}\label{sec:swap}
% ---------------------------------------------------------------------

\begin{definition}{Porte SWAP (échange)}
La porte SWAP échange les états des deux qubits :
\[
  \mathrm{SWAP}\big(\ket{\psi}\otimes\ket{\phi}\big)=\ket{\phi}\otimes\ket{\psi}.
\]
Sur les états de base, $\mathrm{SWAP}\ket{a,b}=\ket{b,a}$ ; on prolonge par linéarité.
\end{definition}

Son symbole est formé de deux croix reliées par un trait vertical
(figure~\ref{fig:swap}).

\begin{figure}[!ht]
\centering
\begin{tikzpicture}[>=Stealth, font=\small]
  \draw (0,1.2) -- (6,1.2);
  \draw (0,0) -- (6,0);
  \node[anchor=east] at (0,1.2) {$\ket{a}$};
  \node[anchor=east] at (0,0) {$\ket{b}$};
  \node[anchor=west] at (6,1.2) {$\ket{b}$};
  \node[anchor=west] at (6,0) {$\ket{a}$};
  \draw[insarouge, thick] (3,1.2) -- (3,0);
  \pic at (3,1.2) {croix};
  \pic at (3,0) {croix};
\end{tikzpicture}
\caption{Porte SWAP : les deux qubits échangent leurs états.}
\label{fig:swap}
\end{figure}

\subsubsection{Table et matrice}

\begin{center}
\renewcommand{\arraystretch}{1.3}
\begin{tabular}{@{}cc@{}}
\toprule
Entrée & Sortie \\
\midrule
$\ket{00}$ & $\ket{00}$ \\
$\ket{01}$ & $\ket{10}$ \\
$\ket{10}$ & $\ket{01}$ \\
$\ket{11}$ & $\ket{11}$ \\
\bottomrule
\end{tabular}
\end{center}

Les images des quatre vecteurs de base sont les colonnes de la matrice :
\[
  \mathrm{SWAP}\ket{00}=\begin{pmatrix}1\\0\\0\\0\end{pmatrix},\quad
  \mathrm{SWAP}\ket{01}=\begin{pmatrix}0\\0\\1\\0\end{pmatrix},\quad
  \mathrm{SWAP}\ket{10}=\begin{pmatrix}0\\1\\0\\0\end{pmatrix},\quad
  \mathrm{SWAP}\ket{11}=\begin{pmatrix}0\\0\\0\\1\end{pmatrix},
\]
\[
  \mathrm{SWAP}=\begin{pmatrix}
    1 & 0 & 0 & 0\\ 0 & 0 & 1 & 0\\ 0 & 1 & 0 & 0\\ 0 & 0 & 0 & 1
  \end{pmatrix}.
\]
Appliquée à $\ket{01}=(0,1,0,0)^{\mathsf T}$, cette matrice sélectionne sa deuxième colonne :
\[
  \begin{pmatrix}
    1 & 0 & 0 & 0\\ 0 & 0 & 1 & 0\\ 0 & 1 & 0 & 0\\ 0 & 0 & 0 & 1
  \end{pmatrix}
  \begin{pmatrix}0\\1\\0\\0\end{pmatrix}
  =\begin{pmatrix}0\\0\\1\\0\end{pmatrix}=\ket{10}.
\]
Sur un état quelconque, $\mathrm{SWAP}\,(c_{00},c_{01},c_{10},c_{11})^{\mathsf T}
=(c_{00},c_{10},c_{01},c_{11})^{\mathsf T}$ : les amplitudes de $\ket{01}$ et $\ket{10}$
sont échangées.

\subsubsection{Accord avec la définition}

Pour deux états $\ket{\psi}=(a_0,a_1)^{\mathsf T}$ et $\ket{\phi}=(b_0,b_1)^{\mathsf T}$ :
\[
  \mathrm{SWAP}\,(\ket{\psi}\otimes\ket{\phi})
  =\mathrm{SWAP}\begin{pmatrix}a_0b_0\\a_0b_1\\a_1b_0\\a_1b_1\end{pmatrix}
  =\begin{pmatrix}a_0b_0\\a_1b_0\\a_0b_1\\a_1b_1\end{pmatrix}
  =\begin{pmatrix}b_0\\b_1\end{pmatrix}\otimes\begin{pmatrix}a_0\\a_1\end{pmatrix}
  =\ket{\phi}\otimes\ket{\psi},
\]
ce qui est bien la définition. La porte s'écrit aussi avec des projecteurs :
\[
  \mathrm{SWAP}=\sum_{a,b\in\{0,1\}}\ket{a}\bra{b}\otimes\ket{b}\bra{a}.
\]
En effet, pour $\ket{c}\otimes\ket{d}$ on trouve
$\sum_{a,b}\ket{a}\braket{b}{c}\otimes\ket{b}\braket{a}{d}$, et seul le terme $b=c$, $a=d$
est non nul : le résultat est $\ket{d}\otimes\ket{c}$.

\paragraph{Exemples.}
\begin{itemize}[nosep]
  \item Deux états différents : $\mathrm{SWAP}\big(\ket{0}\otimes\ket{+}\big)
        =\mathrm{SWAP}\,\frac{1}{\sqrt2}(1,1,0,0)^{\mathsf T}=\frac{1}{\sqrt2}(1,0,1,0)^{\mathsf T}
        =\ket{+}\otimes\ket{0}$.
  \item L'état $\ket{\psi}=\frac{1}{\sqrt2}\ket{00}+\frac{1}{\sqrt2}\ket{11}$ ne change pas :
        $\mathrm{SWAP}\ket{\psi}=\frac{1}{\sqrt2}\ket{00}+\frac{1}{\sqrt2}\ket{11}=\ket{\psi}$,
        car $\ket{00}$ et $\ket{11}$ sont invariants. De même
        $\mathrm{SWAP}\ket{\Phi^-}=\ket{\Phi^-}$ et $\mathrm{SWAP}\ket{\Psi^+}=\ket{\Psi^+}$, alors que
        $\mathrm{SWAP}\ket{\Psi^-}=\frac{1}{\sqrt2}(\ket{10}-\ket{01})=-\ket{\Psi^-}$.
\end{itemize}

\paragraph{Propriétés.} Échanger deux fois redonne l'état de départ :
$\mathrm{SWAP}^2=\mathrm{Id}$ ; la matrice est réelle et symétrique, donc
$\mathrm{SWAP}^\dagger\mathrm{SWAP}=\mathrm{SWAP}^2=\mathrm{Id}$ : elle est unitaire.

\subsubsection{Lien avec le CNOT}

\begin{resultat}{SWAP échange les rôles de contrôle et de cible}
\[
  \mathrm{CNOT}_{21}=\mathrm{SWAP}\;\mathrm{CNOT}_{12}\;\mathrm{SWAP}.
\]
\end{resultat}

\emph{Sur les états de base.}
$\mathrm{SWAP}\,\mathrm{CNOT}_{12}\,\mathrm{SWAP}\ket{b,a}
=\mathrm{SWAP}\,\mathrm{CNOT}_{12}\ket{a,b}
=\mathrm{SWAP}\ket{a,\,a\oplus b}=\ket{a\oplus b,\,a}=\mathrm{CNOT}_{21}\ket{b,a}$.

\emph{Par les matrices.} Multiplier à gauche par $\mathrm{SWAP}$ échange les lignes $2$ et
$3$, multiplier à droite l'échange les colonnes $2$ et $3$ :
\[
  \mathrm{SWAP}\,\mathrm{CNOT}_{12}
  =\begin{pmatrix}1&0&0&0\\0&0&0&1\\0&1&0&0\\0&0&1&0\end{pmatrix},
  \qquad
  \mathrm{SWAP}\,\mathrm{CNOT}_{12}\,\mathrm{SWAP}
  =\begin{pmatrix}1&0&0&0\\0&0&0&1\\0&0&1&0\\0&1&0&0\end{pmatrix}=\mathrm{CNOT}_{21}.
\]

\begin{remarque}[SWAP avec trois CNOT]
$\mathrm{SWAP}=\mathrm{CNOT}_{12}\,\mathrm{CNOT}_{21}\,\mathrm{CNOT}_{12}$. En effet, sur
$\ket{a,b}$ (on applique les portes de droite à gauche) :
\[
  \ket{a,b}\ \xrightarrow{\ \mathrm{CNOT}_{12}\ }\ \ket{a,\,a\oplus b}
  \ \xrightarrow{\ \mathrm{CNOT}_{21}\ }\ \ket{a\oplus(a\oplus b),\,a\oplus b}=\ket{b,\,a\oplus b}
  \ \xrightarrow{\ \mathrm{CNOT}_{12}\ }\ \ket{b,\,b\oplus(a\oplus b)}=\ket{b,a}.
\]
Les deux membres coïncident sur les états de base, donc sur tous les états par linéarité.
\end{remarque}

% ---------------------------------------------------------------------
\subsection{Un circuit quantique : le générateur d'états de Bell}\label{sec:circuit-bell}
% ---------------------------------------------------------------------

On assemble une porte de Hadamard et un CNOT. Le système est $(A,B)$ avec le ket
$\ket{a,b}=\ket{a}_A\otimes\ket{b}_B$ ; d'après la convention de la section, $B$ est dessiné
en haut et $A$ en bas (figure~\ref{fig:circuit-bell}). Le circuit a deux étapes :
\begin{enumerate}[nosep]
  \item la porte $I\otimes H$ : identité sur $A$, Hadamard sur $B$ ;
  \item un CNOT de contrôle $B$ et de cible $A$ (la matrice $\mathrm{CNOT}_{21}$ de la
        section~\ref{sec:cnot}, puisque le contrôle est le second chiffre).
\end{enumerate}
Les traits pointillés délimitent les états $\ket{\Pi_0}$, $\ket{\Pi_1}$ et $\ket{\Pi_2}$ :
avant les portes, après $I\otimes H$, après le CNOT.

\begin{figure}[!ht]
\centering
\begin{tikzpicture}[>=Stealth, font=\small]
  % fils : B en haut, A en bas
  \draw (0,1.4) -- (7.6,1.4);
  \draw (0,0) -- (7.6,0);
  \node[anchor=east] at (0,1.4) {$\ket{0}_B$};
  \node[anchor=east] at (0,0) {$\ket{0}_A$};
  % première étape : H sur B, identité (pointillée) sur A
  \node[porte] at (2.2,1.4) {$H$};
  \node[draw=insagris, dashed, fill=white, minimum size=7mm] at (2.2,0) {$I$};
  % CNOT : contrôle B (haut), cible A (bas)
  \draw[insarouge, thick] (5,1.4) -- (5,0);
  \pic at (5,1.4) {ctl};
  \pic at (5,0) {cible};
  % accolade de sortie
  \draw[insagris, thick] (7.75,1.6) -- (7.9,1.6) -- (7.9,-0.2) -- (7.75,-0.2);
  \node[anchor=west] at (7.9,0.7) {$\ket{\beta^{00}}_{AB}$};
  % états intermédiaires
  \foreach \x in {0.7,3.6,6.4}
    \draw[dashed, insagris!70] (\x,1.9) -- (\x,-0.7);
  \node[anchor=north, font=\footnotesize] at (0.7,-0.7) {$\ket{\Pi_0}=\ket{00}$};
  \node[anchor=north, font=\footnotesize] at (3.6,-0.7) {$\ket{\Pi_1}=\ket{0}\otimes\ket{+}$};
  \node[anchor=north, font=\footnotesize] at (6.4,-0.7) {$\ket{\Pi_2}=\ket{\beta^{00}}$};
\end{tikzpicture}
\caption{Générateur d'états de Bell : $I\otimes H$, puis un CNOT de contrôle $B$ (fil du haut)
et de cible $A$ (fil du bas). L'entrée représentée est $\ket{00}$.}
\label{fig:circuit-bell}
\end{figure}

\subsubsection{Entrée \texorpdfstring{$\ket{00}$}{|00>} : calcul pas à pas}

\begin{align*}
  \ket{\Pi_0}&=\ket{0}_A\otimes\ket{0}_B=\ket{00},\\[1ex]
  \ket{\Pi_1}&=(I\otimes H)\ket{\Pi_0}=I\ket{0}_A\otimes H\ket{0}_B=\ket{0}\otimes\ket{+}
    =\frac{\ket{00}+\ket{01}}{\sqrt2},\\[1ex]
  \ket{\Pi_2}&=\mathrm{CNOT}\ket{\Pi_1}=\frac{\mathrm{CNOT}\ket{00}+\mathrm{CNOT}\ket{01}}{\sqrt2}
    =\frac{\ket{00}+\ket{11}}{\sqrt2}=\ket{\beta^{00}}=\ket{\Phi^+}.
\end{align*}
Dans $\ket{\Pi_1}$, le premier chiffre est celui de $A$ et le second celui de $B$. Le CNOT a
pour contrôle $B$ : $\ket{00}$ ($B=0$) ne change pas, $\ket{01}$ ($B=1$) devient $\ket{11}$
puisque $A$ est inversé. Le résultat est un \emph{état de Bell}.

Avec les matrices, la règle $A\otimes B=(A_{ij}B)$ donne
\[
  I\otimes H=\begin{pmatrix}1\cdot H&0\\0&1\cdot H\end{pmatrix}
  =\frac{1}{\sqrt2}\begin{pmatrix}1&1&0&0\\1&-1&0&0\\0&0&1&1\\0&0&1&-1\end{pmatrix}.
\]
Sa première colonne est $\ket{\Pi_1}=\frac{1}{\sqrt2}(1,1,0,0)^{\mathsf T}$ (image de $\ket{00}$).
Le CNOT échange les amplitudes de $\ket{01}$ et de $\ket{11}$ (section~\ref{sec:cnot}), donc
\[
  \ket{\Pi_2}=\mathrm{CNOT}\,\frac{1}{\sqrt2}\begin{pmatrix}1\\1\\0\\0\end{pmatrix}
  =\frac{1}{\sqrt2}\begin{pmatrix}1\\0\\0\\1\end{pmatrix}.
\]

\subsubsection{Les quatre entrées}

Prenons $B$ dans l'état $\ket{x}$ et $A$ dans l'état $\ket{y}$, avec $x,y\in\{0,1\}$. Alors
$\ket{\Pi_0}=\ket{y}_A\otimes\ket{x}_B=\ket{y,x}$. Comme
$H\ket{x}=\frac{1}{\sqrt2}\big(\ket{0}+(-1)^x\ket{1}\big)$ (c'est $\ket{+}$ pour $x=0$ et
$\ket{-}$ pour $x=1$),
\[
  \ket{\Pi_1}=\ket{y}\otimes H\ket{x}=\frac{1}{\sqrt2}\Big(\ket{y,0}+(-1)^x\ket{y,1}\Big).
\]
Le CNOT (contrôle $B$) ne change pas $\ket{y,0}$ et transforme $\ket{y,1}$ en
$\ket{\bar y,1}$, avec $\bar y=1-y$ :
\[
  \ket{\Pi_2}=\ket{\beta^{xy}}=\frac{1}{\sqrt2}\Big(\ket{y,0}+(-1)^x\ket{\bar y,1}\Big).
\]
Le premier indice de $\ket{\beta^{xy}}$ est l'entrée de $B$ (le qubit qui reçoit $H$), le
second celle de $A$ (la cible). Les quatre cas :

\begin{center}
\renewcommand{\arraystretch}{1.7}
\begin{tabular}{@{}ccccc@{}}
\toprule
$(B,A)=(x,y)$ & $\ket{\Pi_0}=\ket{y,x}$ & $\ket{\Pi_1}$ & $\ket{\Pi_2}=\ket{\beta^{xy}}$ & Nom \\
\midrule
$(0,0)$ & $\ket{00}$ & $\frac{1}{\sqrt2}(\ket{00}+\ket{01})$ & $\frac{1}{\sqrt2}(\ket{00}+\ket{11})$ & $\ket{\Phi^+}$ \\
$(1,0)$ & $\ket{01}$ & $\frac{1}{\sqrt2}(\ket{00}-\ket{01})$ & $\frac{1}{\sqrt2}(\ket{00}-\ket{11})$ & $\ket{\Phi^-}$ \\
$(0,1)$ & $\ket{10}$ & $\frac{1}{\sqrt2}(\ket{10}+\ket{11})$ & $\frac{1}{\sqrt2}(\ket{10}+\ket{01})$ & $\ket{\Psi^+}$ \\
$(1,1)$ & $\ket{11}$ & $\frac{1}{\sqrt2}(\ket{10}-\ket{11})$ & $\frac{1}{\sqrt2}(\ket{10}-\ket{01})$ & $-\ket{\Psi^-}$ \\
\bottomrule
\end{tabular}
\end{center}
Détail de la dernière ligne : $\ket{\Pi_1}=\ket{1}\otimes\ket{-}=\frac{1}{\sqrt2}(\ket{10}-\ket{11})$ ;
le CNOT laisse $\ket{10}$ et envoie $\ket{11}$ sur $\ket{01}$, d'où
$\frac{1}{\sqrt2}(\ket{10}-\ket{01})=-\frac{1}{\sqrt2}(\ket{01}-\ket{10})=-\ket{\Psi^-}$.

\paragraph{Forme matricielle.} $\mathrm{CNOT}_{21}$ échange les lignes $2$ et $4$ de
$I\otimes H$ (il échange les amplitudes de $\ket{01}$ et $\ket{11}$), donc
\[
  \mathrm{CNOT}_{21}\,(I\otimes H)=\frac{1}{\sqrt2}\begin{pmatrix}
    1 & 1 & 0 & 0\\ 0 & 0 & 1 & -1\\ 0 & 0 & 1 & 1\\ 1 & -1 & 0 & 0
  \end{pmatrix}.
\]
Ses colonnes sont les images de $\ket{00},\ket{01},\ket{10},\ket{11}$, c'est-à-dire de
$(B,A)=(0,0),(1,0),(0,1),(1,1)$ : ce sont $\ket{\beta^{00}}$, $\ket{\beta^{10}}$,
$\ket{\beta^{01}}$, $\ket{\beta^{11}}$, soit $\ket{\Phi^+}$, $\ket{\Phi^-}$, $\ket{\Psi^+}$
et $-\ket{\Psi^-}$.

\begin{remarque}[Comparaison avec la section précédente]
Le circuit de la figure~\ref{fig:bell} applique $H$ au \emph{premier} chiffre et prend celui-ci
pour contrôle ; ici $H$ et le contrôle portent sur le \emph{dernier} chiffre. Les deux
circuits se déduisent l'un de l'autre en échangeant les deux fils (conjugaison par le SWAP,
section~\ref{sec:swap}). Les quatre états de Bell obtenus sont les mêmes, avec les mêmes
indices, sauf le dernier : $-\ket{\Psi^-}$ ici, $+\ket{\Psi^-}$ là-bas. Ce signe est une phase
globale, sans effet sur les mesures (elle disparaît de $\rho$, section~\ref{sec:rho}).
\end{remarque}

\subsubsection{Matrice de densité de l'état de sortie}

Pour $\ket{\Psi}=\ket{\beta^{00}}=\frac{1}{\sqrt2}\ket{00}+\frac{1}{\sqrt2}\ket{11}$ :
\[
  \rho_{AB}=\ket{\Psi}\bra{\Psi}
  =\begin{pmatrix}\frac{1}{\sqrt2}\\0\\0\\\frac{1}{\sqrt2}\end{pmatrix}
   \begin{pmatrix}\frac{1}{\sqrt2}&0&0&\frac{1}{\sqrt2}\end{pmatrix}
  =\frac12\begin{pmatrix}1&0&0&1\\0&0&0&0\\0&0&0&0\\1&0&0&1\end{pmatrix}.
\]
Sous forme de ket-bra :
$\rho_{AB}=\frac12\big(\ket{00}\bra{00}+\ket{00}\bra{11}+\ket{11}\bra{00}+\ket{11}\bra{11}\big)$.
On vérifie $\operatorname{Tr}\rho_{AB}=\frac12(1+1)=1$ et
\[
  \rho_{AB}^2=\frac14\begin{pmatrix}2&0&0&2\\0&0&0&0\\0&0&0&0\\2&0&0&2\end{pmatrix}=\rho_{AB},
\]
donc $\operatorname{Tr}\rho_{AB}^2=1$ : le couple $(A,B)$ est dans un état \emph{pur}.

% ---------------------------------------------------------------------
\subsection{Information : variables aléatoires et états quantiques}\label{sec:info}
% ---------------------------------------------------------------------

\begin{definition}{Entropie, information mutuelle, entropie conditionnelle}
Soit $A$ une variable aléatoire de loi $p(a)$. Son \emph{entropie de Shannon}, en bits, est
\[
  H(A)=-\sum_a p(a)\log_2p(a)\qquad(\text{avec }0\log_20=0).
\]
Elle mesure l'incertitude moyenne sur $A$ : une pièce équilibrée a $H=1$ bit. Pour un couple
$(A,B)$ de loi $p(a,b)$, $H(A,B)=-\sum_{a,b}p(a,b)\log_2p(a,b)$, et
\begin{align*}
  I(A;B)&=H(A)+H(B)-H(A,B) &&\text{(information mutuelle)},\\
  H(A|B)&=H(A,B)-H(B)=H(A)-I(A;B) &&\text{(entropie conditionnelle)}.
\end{align*}
\end{definition}

Les deux écritures de $H(A|B)$ sont égales : la définition de $I(A;B)$ donne
$H(A,B)=H(A)+H(B)-I(A;B)$, donc $H(A,B)-H(B)=H(A)-I(A;B)$. L'information mutuelle $I(A;B)$
mesure ce que la connaissance de $B$ apprend sur $A$ ; $H(A|B)$ est l'incertitude qui reste sur
$A$ quand on connaît $B$.

\subsubsection{Deux variables aléatoires indépendantes}

Alice et Bob lancent chacun une pièce équilibrée, sans lien entre les deux lancers :
$p(a,b)=\frac14$ pour les quatre couples. Chaque pièce donne $0$ ou $1$ avec la probabilité $\frac12$, donc
\begin{align*}
  H(A)&=-\Big(\tfrac12\log_2\tfrac12+\tfrac12\log_2\tfrac12\Big)=1\text{ bit},\\
  H(B)&=1\text{ bit (même calcul)},\\
  H(A,B)&=-4\cdot\tfrac14\log_2\tfrac14=2\text{ bits},\\
  I(A;B)&=H(A)+H(B)-H(A,B)=1+1-2=0,\\
  H(A|B)&=H(A)-I(A;B)=1-0=H(A)=1\text{ bit}.
\end{align*}
Connaître $B$ n'apprend rien sur $A$.

\subsubsection{Deux variables solidaires}

Les deux pièces tombent toujours du même côté : $p(0,0)=p(1,1)=\frac12$ et
$p(0,1)=p(1,0)=0$. Les lois de chaque pièce sont inchangées ($p(a{=}0)=p(0,0)+p(0,1)=\frac12$,
etc.), donc $H(A)=H(B)=1$ bit, mais
\begin{align*}
  H(A,B)&=-2\cdot\tfrac12\log_2\tfrac12-2\cdot0\log_20=1\text{ bit},\\
  I(A;B)&=1+1-1=1\text{ bit},\qquad H(A|B)=1-1=0.
\end{align*}
Connaître $B$ détermine $A$ : plus aucune incertitude ne reste.

\subsubsection{Deux qubits dans un état de Bell}

Pour un système quantique, on remplace la loi de probabilité par la matrice de densité.

\begin{definition}{Trace partielle et entropie de von Neumann}
L'état du seul qubit $A$ d'un système $(A,B)$ de matrice de densité $\rho_{AB}$ est obtenu en
\emph{traçant sur $B$} : $\rho_A=\operatorname{Tr}_B(\rho_{AB})$. La trace partielle est
linéaire et remplace la partie $B$ de chaque terme par sa trace :
\[
  \operatorname{Tr}_B\big(\ket{a}\bra{a'}\otimes\ket{b}\bra{b'}\big)=\ket{a}\bra{a'}\;\braket{b'}{b}.
\]
L'\emph{entropie de von Neumann} d'un état $\rho$, de valeurs propres $\lambda_i$, est
$S(\rho)=-\operatorname{Tr}(\rho\log_2\rho)=-\sum_i\lambda_i\log_2\lambda_i$. En quantique,
$H(A)=S(\rho_A)$, $H(B)=S(\rho_B)$ et $H(A,B)=S(\rho_{AB})$ ; les définitions de
$I(A;B)$ et de $H(A|B)$ sont les mêmes.
\end{definition}

On part de l'état de Bell
$\rho_{AB}=\frac12\big(\ket{00}\bra{00}+\ket{00}\bra{11}+\ket{11}\bra{00}+\ket{11}\bra{11}\big)$
produit par le circuit de la section~\ref{sec:circuit-bell}. Chaque terme est un produit
d'opérateurs à un qubit :
$\ket{00}\bra{11}=\ket{0}\bra{1}\otimes\ket{0}\bra{1}$, etc. Alors
\begin{align*}
  \rho_A=\operatorname{Tr}_B(\rho_{AB})
  &=\tfrac12\Big(\ket{0}\bra{0}\,\braket{0}{0}+\ket{0}\bra{1}\,\braket{1}{0}
      +\ket{1}\bra{0}\,\braket{0}{1}+\ket{1}\bra{1}\,\braket{1}{1}\Big)\\
  &=\tfrac12\Big(\ket{0}\bra{0}\cdot1+\ket{0}\bra{1}\cdot0+\ket{1}\bra{0}\cdot0+\ket{1}\bra{1}\cdot1\Big)\\
  &=\tfrac12\big(\ket{0}\bra{0}+\ket{1}\bra{1}\big)=\tfrac12\,\mathrm{Id}
  =\begin{pmatrix}\frac12&0\\0&\frac12\end{pmatrix}.
\end{align*}
La formule des composantes de la section~\ref{sec:separable},
$(\rho_A)_{ij}=\rho_{(i0),(j0)}+\rho_{(i1),(j1)}$, donne le même résultat. En traçant sur $A$
au lieu de $B$, la règle est $\operatorname{Tr}_A(\ket{a}\bra{a'}\otimes\ket{b}\bra{b'})
=\braket{a'}{a}\,\ket{b}\bra{b'}$, et on trouve de même $\rho_B=\frac12\mathrm{Id}$.

Calcul des entropies :
\begin{itemize}[nosep]
  \item $\rho_{AB}$ est pur : ses valeurs propres sont $1,0,0,0$, donc
        $H(A,B)=S(\rho_{AB})=-1\cdot\log_21=0$ ;
  \item $\rho_A$ et $\rho_B$ ont pour valeurs propres $\frac12,\frac12$, donc
        $H(A)=H(B)=-\big(\frac12\log_2\frac12+\frac12\log_2\frac12\big)=1$ bit.
\end{itemize}
D'où
\begin{align*}
  I(A;B)&=H(A)+H(B)-H(A,B)=1+1-0=2\text{ bits},\\
  H(A|B)&=H(A,B)-H(B)=0-1=-1\text{ bit}\qquad(=H(A)-I(A;B)=1-2).
\end{align*}

\begin{remarque}[Une entropie conditionnelle négative]
Pour des variables classiques, $p(a,b)\leq p(b)$, donc $-\log_2p(a,b)\geq-\log_2p(b)$ et
$H(A,B)\geq H(B)$ : on a toujours $H(A|B)\geq0$, et $I(A;B)\leq H(A)$. Ici $H(A|B)=-1$ et
$I(A;B)=2>H(A)$ : ces valeurs n'ont pas d'équivalent classique, c'est une signature de
l'intrication. Le résultat est surprenant mais mathématiquement correct.
\end{remarque}

\begin{center}
\renewcommand{\arraystretch}{1.4}
\begin{tabular}{@{}lccc@{}}
\toprule
 & indépendantes & solidaires & état de Bell \\
\midrule
$H(A)$ & $1$ & $1$ & $1$ \\
$H(B)$ & $1$ & $1$ & $1$ \\
$H(A,B)$ & $2$ & $1$ & $0$ \\
$I(A;B)$ & $0$ & $1$ & $2$ \\
$H(A|B)$ & $1$ & $0$ & $-1$ \\
\bottomrule
\end{tabular}
\end{center}
(en bits ; les deux premières colonnes sont classiques, la troisième est quantique).

% ---------------------------------------------------------------------
\subsection{Système pur et système mélangé}\label{sec:pur-melange}
% ---------------------------------------------------------------------

Rappel (section~\ref{sec:pur-mixte}) : si $\lambda_i$ sont les valeurs propres de la matrice de
densité $\rho$, alors $\operatorname{Tr}\rho^2=\sum_i\lambda_i^2\leq1$. Le système est
\emph{pur} si $\operatorname{Tr}\rho^2=1$ (alors $\rho=\ket{\psi}\bra{\psi}$ et $\rho^2=\rho$),
et \emph{mélangé} si $\operatorname{Tr}\rho^2<1$. Pour un qubit,
$\operatorname{Tr}\rho^2=\frac{1+|\vec r|^2}{2}$ (section~\ref{sec:bloch}) est compris entre
$\frac12$ et $1$.

\subsubsection{Un système pur : exercice}

\emph{Énoncé.} Soit $\ket{\psi}=\frac{1}{\sqrt2}\ket{0}+\frac{1}{\sqrt2}\ket{1}$.
\begin{enumerate}[label=\alph*), nosep]
  \item Calculer la matrice de densité $\rho$.
  \item Calculer la pureté du système quantique.
  \item Calculer $\langle\sigma_x\rangle$ et $\langle\sigma_z\rangle$.
  \item Écrire $\ket{\psi}$ dans la base $X$.
\end{enumerate}

\paragraph{a) Matrice de densité.} $\rho=\ket{\psi}\bra{\psi}$ est le produit d'un vecteur
colonne par un vecteur ligne (les amplitudes sont réelles) :
\[
  \rho=\begin{pmatrix}\frac{1}{\sqrt2}\\[0.4ex]\frac{1}{\sqrt2}\end{pmatrix}
       \begin{pmatrix}\frac{1}{\sqrt2}&\frac{1}{\sqrt2}\end{pmatrix}
  =\begin{pmatrix}\frac12&\frac12\\[0.4ex]\frac12&\frac12\end{pmatrix}.
\]

\paragraph{b) Pureté.}
\[
  \rho^2=\begin{pmatrix}\frac12&\frac12\\[0.4ex]\frac12&\frac12\end{pmatrix}
         \begin{pmatrix}\frac12&\frac12\\[0.4ex]\frac12&\frac12\end{pmatrix}
  =\begin{pmatrix}\frac14+\frac14&\frac14+\frac14\\[0.4ex]\frac14+\frac14&\frac14+\frac14\end{pmatrix}
  =\begin{pmatrix}\frac12&\frac12\\[0.4ex]\frac12&\frac12\end{pmatrix}=\rho,
\]
donc $\operatorname{Tr}\rho^2=\operatorname{Tr}\rho=\frac12+\frac12=1$ : le système est \emph{pur}.

\paragraph{c) Valeurs moyennes.} Avec $\sigma_x=X=\begin{pmatrix}0&1\\1&0\end{pmatrix}$ et
$\sigma_z=Z=\begin{pmatrix}1&0\\0&-1\end{pmatrix}$ (section~\ref{sec:op-dirac}), la valeur
moyenne est $\langle\sigma\rangle=\operatorname{Tr}(\rho\,\sigma)$ (section~\ref{sec:rho}).
\begin{align*}
  \rho\,\sigma_x&=\begin{pmatrix}\frac12&\frac12\\[0.4ex]\frac12&\frac12\end{pmatrix}
     \begin{pmatrix}0&1\\1&0\end{pmatrix}
     =\begin{pmatrix}\frac12&\frac12\\[0.4ex]\frac12&\frac12\end{pmatrix},
     &\langle\sigma_x\rangle&=\operatorname{Tr}(\rho\,\sigma_x)=\tfrac12+\tfrac12=1,\\[1.5ex]
  \rho\,\sigma_z&=\begin{pmatrix}\frac12&\frac12\\[0.4ex]\frac12&\frac12\end{pmatrix}
     \begin{pmatrix}1&0\\0&-1\end{pmatrix}
     =\begin{pmatrix}\frac12&-\frac12\\[0.4ex]\frac12&-\frac12\end{pmatrix},
     &\langle\sigma_z\rangle&=\operatorname{Tr}(\rho\,\sigma_z)=\tfrac12-\tfrac12=0.
\end{align*}
On retrouve $\langle\sigma_z\rangle$ avec les probabilités de mesure : dans la base $Z$,
$\sigma_z$ vaut $+1$ pour $\ket{0}$ et $-1$ pour $\ket{1}$, et
$|\braket{0}{\psi}|^2=|\braket{1}{\psi}|^2=\frac12$, donc
$\langle\sigma_z\rangle=\frac12\times(+1)+\frac12\times(-1)=0$.

\paragraph{d) Écriture dans la base $X$.} La base $X$ est $\{\ket{+},\ket{-}\}$, orthonormée, donc
$\ket{+}\bra{+}+\ket{-}\bra{-}=\mathrm{Id}$ et
$\ket{\psi}=\ket{+}\braket{+}{\psi}+\ket{-}\braket{-}{\psi}$, avec
\[
  \braket{+}{\psi}=\tfrac{1}{\sqrt2}\Big(\tfrac{1}{\sqrt2}+\tfrac{1}{\sqrt2}\Big)=1,
  \qquad
  \braket{-}{\psi}=\tfrac{1}{\sqrt2}\Big(\tfrac{1}{\sqrt2}-\tfrac{1}{\sqrt2}\Big)=0 .
\]
Donc $\ket{\psi}=1\cdot\ket{+}+0\cdot\ket{-}=\ket{+}$ : dans la base $X$, l'état n'a qu'une
seule composante. Comme $\sigma_x\ket{+}=+\ket{+}$ et $\sigma_x\ket{-}=-\ket{-}$, une
mesure de $\sigma_x$ donne $+1$ avec certitude, et
$\langle\sigma_x\rangle=1\cdot(+1)+0\cdot(-1)=1$ : c'est le résultat de la question c).
Le vecteur de Bloch de $\ket{+}$ est $(\langle\sigma_x\rangle,\langle\sigma_y\rangle,\langle\sigma_z\rangle)=(1,0,0)$,
de norme $1$ : l'état est sur la sphère, en accord avec $\operatorname{Tr}\rho^2=\frac{1+1}{2}=1$.

\subsubsection{Un système mélangé : un qubit de la paire de Bell}

À la section~\ref{sec:info}, la trace sur $B$ de l'état de Bell a donné
$\rho_A=\frac12\mathrm{Id}=\begin{pmatrix}\frac12&0\\0&\frac12\end{pmatrix}$. Alors
\[
  \rho_A^2=\begin{pmatrix}\frac14&0\\0&\frac14\end{pmatrix},
  \qquad
  \operatorname{Tr}\big(\rho_A^2\big)=\frac14+\frac14=\frac12<1 :
\]
le qubit $A$ est dans un état \emph{mélangé}, et c'est le plus mélangé possible pour un qubit
(centre de la boule de Bloch, $\vec r=\vec0$). Pourtant le couple $(A,B)$ est pur
($\operatorname{Tr}\rho_{AB}^2=1$) : on connaît parfaitement l'état du couple, mais pas celui d'un
qubit pris seul. L'information est entièrement dans les corrélations entre $A$ et $B$, et
$S(\rho_A)=1$ bit.

\subsubsection{Un système mélangé : le mélange statistique de Bob}

Bob prépare un qubit dans l'état $\ket{0}$ avec la probabilité $p_1=\frac34$ (75\,\%) et dans
l'état $\ket{1}$ avec la probabilité $p_2=\frac14$ (25\,\%), sans que l'on sache lequel des deux a
été préparé. Ce n'est pas un état pur $\ket{\psi}$ : on décrit le système par la matrice de
densité $\rho=\sum_ip_i\ket{\psi_i}\bra{\psi_i}$ (section~\ref{sec:pur-mixte}) :
\[
  \rho=\underbrace{\tfrac34}_{p_1}\ket{0}\bra{0}+\underbrace{\tfrac14}_{p_2}\ket{1}\bra{1}
  =\begin{pmatrix}\frac34&0\\0&\frac14\end{pmatrix}.
\]
Alors
\[
  \rho^2=\begin{pmatrix}\frac{9}{16}&0\\0&\frac{1}{16}\end{pmatrix},
  \qquad
  \operatorname{Tr}\big(\rho^2\big)=\frac{9}{16}+\frac{1}{16}=\frac{10}{16}=\frac58<1 :
\]
le système est \emph{mélangé}.

Les valeurs moyennes sont les moyennes pondérées sur la préparation :
$\langle\sigma_z\rangle=\operatorname{Tr}(\rho\,\sigma_z)=\frac34-\frac14=\frac12
=p_1\cdot(+1)+p_2\cdot(-1)$, et $\langle\sigma_x\rangle=\operatorname{Tr}(\rho\,\sigma_x)=0$
(les termes hors diagonale de $\rho$ sont nuls). Le vecteur de Bloch est $(0,0,\frac12)$,
à l'intérieur de la boule, et $\operatorname{Tr}\rho^2=\frac{1+1/4}{2}=\frac58$. Son entropie
est $S(\rho)=-\frac34\log_2\frac34-\frac14\log_2\frac14=2-\frac34\log_23\approx0{,}811$ bit.

\begin{remarque}[Mélange et superposition]
Considérons l'état pur $\ket{\psi'}=\frac{\sqrt3}{2}\ket{0}+\frac12\ket{1}$, qui a les mêmes
probabilités en base $Z$ ($P(0)=\frac34$, $P(1)=\frac14$). Sa matrice de densité est
\[
  \rho'=\begin{pmatrix}\frac34&\frac{\sqrt3}{4}\\[0.4ex]\frac{\sqrt3}{4}&\frac14\end{pmatrix},
  \qquad
  \operatorname{Tr}\rho'^2=\frac{9}{16}+\frac1{16}+2\cdot\frac{3}{16}=1,
  \qquad
  \langle\sigma_x\rangle=2\cdot\frac{\sqrt3}{4}=\frac{\sqrt3}{2}.
\]
Le mélange de Bob a la même diagonale mais aucun terme hors diagonale ($\operatorname{Tr}\rho^2=\frac58$,
$\langle\sigma_x\rangle=0$) : une superposition n'est pas un tirage au hasard entre $\ket{0}$ et
$\ket{1}$. De plus, pour un état pur il existe une mesure au résultat certain (ici $\sigma_x$ pour
$\ket{+}$) ; pour le mélange de Bob, la meilleure probabilité d'obtenir un même résultat est
$75\,\%$ (mesure en base $Z$, résultat $\ket{0}$).
\end{remarque}

\subsubsection{Bilan}

\begin{center}
\renewcommand{\arraystretch}{1.5}
\begin{tabular}{@{}lcccl@{}}
\toprule
\textbf{Système} & $\rho$ & $\operatorname{Tr}\rho^2$ & $S(\rho)$ & \textbf{Nature} \\
\midrule
$\ket{\psi}=\ket{+}$ & ${\renewcommand{\arraystretch}{1}\begin{pmatrix}\frac12&\frac12\\\frac12&\frac12\end{pmatrix}}$ & $1$ & $0$ & pur \\[2ex]
Qubit $A$ d'une paire de Bell & ${\renewcommand{\arraystretch}{1}\begin{pmatrix}\frac12&0\\0&\frac12\end{pmatrix}}$ & $\frac12$ & $1$ bit & mélangé \\[2ex]
Source de Bob (75\,\%/25\,\%) & ${\renewcommand{\arraystretch}{1}\begin{pmatrix}\frac34&0\\0&\frac14\end{pmatrix}}$ & $\frac58$ & $\approx0{,}811$ bit & mélangé \\
\bottomrule
\end{tabular}
\end{center}

% ---------------------------------------------------------------------
\subsection{Algorithme quantique : la téléportation}\label{sec:teleportation}
% ---------------------------------------------------------------------

\emph{Problème.} Alice détient un qubit $Q$ dans un état qu'elle ne connaît pas,
$\ket{\varphi}=\alpha\ket{0}+\beta\ket{1}$ avec $|\alpha|^2+|\beta|^2=1$, et elle veut que Bob
obtienne cet état sur l'un de ses qubits. Elle ne peut pas mesurer $Q$ pour connaître
$\alpha$ et $\beta$ (une mesure ne donne qu'un bit et détruit l'état), elle ne peut pas le
copier (théorème de non-clonage) et elle n'envoie pas le qubit lui-même.

\begin{definition}{Protocole de téléportation quantique}
\textbf{Ressources.}
\begin{itemize}[nosep]
  \item Un \emph{e-bit} : Alice et Bob partagent une paire de qubits $(A,B)$ dans l'état de Bell
        $\ket{\Phi^+}_{BA}=\frac{1}{\sqrt2}\big(\ket{00}+\ket{11}\big)$ ; Alice détient $A$,
        Bob détient $B$.
  \item Le qubit $Q$ à téléporter, chez Alice.
  \item Un canal classique (deux bits).
\end{itemize}
Le système est écrit $(B,A,Q)$ : $\ket{b,a,q}=\ket{b}_B\otimes\ket{a}_A\otimes\ket{q}_Q$.

\textbf{Étapes.}
\begin{enumerate}[nosep]
  \item Alice applique un CNOT de contrôle $Q$ et de cible $A$.
  \item Alice applique une porte $H$ sur $Q$.
  \item Alice mesure $A$ et $Q$ dans la base de calcul, obtient deux bits $(a,q)$ et les envoie
        à Bob.
  \item Bob applique $X$ si $a=1$, puis $Z$ si $q=1$, à son qubit $B$.
\end{enumerate}
\end{definition}

\begin{figure}[!ht]
\centering
\begin{adjustbox}{max width=\linewidth}
\begin{tikzpicture}[>=Stealth, font=\small]
  % fils quantiques : Q en haut, A au milieu, B en bas
  \draw (0,2.4) -- (5.3,2.4);
  \draw (0,1.2) -- (5.3,1.2);
  \draw (0,0) -- (10.6,0);
  \node[anchor=east] at (0,2.4) {$\ket{\varphi}_Q$};
  \node[anchor=south west, text=insagris, font=\footnotesize] at (0.05,1.2) {$A$};
  \node[anchor=south west, text=insagris, font=\footnotesize] at (0.05,0) {$B$};
  % paire de Bell partagée
  \draw[insagris, thick] (-0.1,1.5) -- (-0.25,1.5) -- (-0.25,-0.3) -- (-0.1,-0.3);
  \node[anchor=east] at (-0.35,0.6) {$\ket{\Phi^+}_{BA}$};
  % CNOT : contrôle Q, cible A
  \draw[insarouge, thick] (2.6,2.4) -- (2.6,1.2);
  \pic at (2.6,2.4) {ctl};
  \pic at (2.6,1.2) {cible};
  % Hadamard sur Q
  \node[porte] at (4.2,2.4) {$H$};
  % mesures d'Alice
  \pic at (5.6,2.4) {mesure};
  \pic at (5.6,1.2) {mesure};
  % fils classiques (doubles)
  \draw[double, double distance=1.2pt, insagris] (5.9,1.2) -- (7.6,1.2) -- (7.6,0.35);
  \draw[double, double distance=1.2pt, insagris] (5.9,2.4) -- (9.6,2.4) -- (9.6,0.35);
  \node[above, font=\footnotesize] at (6.5,1.2) {$a$};
  \node[above, font=\footnotesize] at (6.5,2.4) {$q$};
  % corrections de Bob
  \node[porte] at (7.6,0) {$X$};
  \node[porte] at (9.6,0) {$Z$};
  \node[anchor=west] at (10.6,0) {$\ket{\varphi}_B$};
  % états intermédiaires
  \foreach \x in {1.4,3.4,4.9}
    \draw[dashed, insagris!70] (\x,2.8) -- (\x,-0.5);
  \node[anchor=north, font=\footnotesize] at (1.4,-0.5) {$\ket{\Pi_0}$};
  \node[anchor=north, font=\footnotesize] at (3.4,-0.5) {$\ket{\Pi_1}$};
  \node[anchor=north, font=\footnotesize] at (4.9,-0.5) {$\ket{\Pi_2}$};
\end{tikzpicture}
\end{adjustbox}
\caption{Circuit de téléportation. Alice détient $Q$ (dans $\ket{\varphi}=\alpha\ket{0}+\beta\ket{1}$)
et $A$, Bob détient $B$ ; $A$ et $B$ forment la paire $\ket{\Phi^+}_{BA}$. Les traits doubles
sont des fils \emph{classiques} qui portent les résultats $a$ (mesure de $A$) et $q$ (mesure de
$Q$) ; $X$ est appliqué si $a=1$, $Z$ si $q=1$. Le ket s'écrit $\ket{b,a,q}$ : le fil du haut
($Q$) est le dernier chiffre.}
\label{fig:teleportation}
\end{figure}

\subsubsection{Calcul pas à pas}

\paragraph{État initial $\ket{\Pi_0}$.} Alice et Bob partagent la paire, et $Q$ est indépendant
d'elle :
\begin{align*}
  \ket{\Pi_0}&=\ket{\Phi^+}_{BA}\otimes\ket{\varphi}_Q
    =\frac{1}{\sqrt2}\big(\ket{00}+\ket{11}\big)\otimes\big(\alpha\ket{0}+\beta\ket{1}\big)\\
  &=\frac{1}{\sqrt2}\Big(\alpha\ket{00}\ket{0}+\beta\ket{00}\ket{1}+\alpha\ket{11}\ket{0}+\beta\ket{11}\ket{1}\Big)\\
  &=\frac{\alpha\ket{000}+\beta\ket{001}+\alpha\ket{110}+\beta\ket{111}}{\sqrt2}.
\end{align*}

\paragraph{Après le CNOT $Q\to A$ : $\ket{\Pi_1}$.} Le CNOT (contrôle $Q$, dernier chiffre ;
cible $A$, chiffre du milieu) agit ainsi : $\ket{b,a,q}\to\ket{b,\,a\oplus q,\,q}$. Terme par terme :
\begin{center}
\renewcommand{\arraystretch}{1.3}
\begin{tabular}{@{}cccc@{}}
\toprule
Terme & $q$ & Effet sur $A$ & Résultat \\
\midrule
$\alpha\ket{000}$ & $0$ & aucun & $\alpha\ket{000}$ \\
$\beta\ket{001}$ & $1$ & $0\to1$ & $\beta\ket{011}$ \\
$\alpha\ket{110}$ & $0$ & aucun & $\alpha\ket{110}$ \\
$\beta\ket{111}$ & $1$ & $1\to0$ & $\beta\ket{101}$ \\
\bottomrule
\end{tabular}
\end{center}
donc
\[
  \ket{\Pi_1}=\frac{\alpha\ket{000}+\beta\ket{011}+\alpha\ket{110}+\beta\ket{101}}{\sqrt2}.
\]
On isole ensuite le dernier chiffre, celui de $Q$, qui va changer à l'étape suivante
(le premier ket est celui du couple $\ket{b,a}$) :
\[
  \ket{\Pi_1}=\frac{\alpha\ket{00}\otimes\ket{0}+\beta\ket{01}\otimes\ket{1}
    +\alpha\ket{11}\otimes\ket{0}+\beta\ket{10}\otimes\ket{1}}{\sqrt2}.
\]

\paragraph{Après la porte $H$ sur $Q$ : $\ket{\Pi_2}$.} On remplace $\ket{0}$ par
$H\ket{0}=\ket{+}$ et $\ket{1}$ par $H\ket{1}=\ket{-}$ dans le facteur $Q$ :
\[
  \ket{\Pi_2}=\frac{1}{\sqrt2}\Big(\alpha\ket{00}\otimes\ket{+}+\beta\ket{01}\otimes\ket{-}
    +\alpha\ket{11}\otimes\ket{+}+\beta\ket{10}\otimes\ket{-}\Big).
\]
Avec $\ket{\pm}=\frac{1}{\sqrt2}(\ket{0}\pm\ket{1})$, le préfacteur devient
$\frac{1}{\sqrt2}\cdot\frac{1}{\sqrt2}=\frac12$ :
\begin{align*}
  \ket{\Pi_2}&=\frac12\Big(\alpha\ket{00}\big(\ket{0}+\ket{1}\big)+\beta\ket{01}\big(\ket{0}-\ket{1}\big)
     +\alpha\ket{11}\big(\ket{0}+\ket{1}\big)+\beta\ket{10}\big(\ket{0}-\ket{1}\big)\Big)\\
  &=\frac12\Big(\alpha\ket{000}+\alpha\ket{001}+\beta\ket{010}-\beta\ket{011}
     +\alpha\ket{110}+\alpha\ket{111}+\beta\ket{100}-\beta\ket{101}\Big).
\end{align*}

\paragraph{Regroupement selon le résultat d'Alice.} Alice mesurera $A$ et $Q$, donc les deux
derniers chiffres. On regroupe les huit termes qui ont les mêmes chiffres $(a,q)$, en mettant
le ket de Bob en premier :
\begin{center}
\renewcommand{\arraystretch}{1.3}
\begin{tabular}{@{}cll@{}}
\toprule
$(a,q)$ & Termes de $\ket{\Pi_2}$ (facteur $\frac12$ omis) & Ket de Bob \\
\midrule
$(0,0)$ & $\alpha\ket{000}+\beta\ket{100}$ & $(\alpha\ket{0}+\beta\ket{1})\ket{00}$ \\
$(0,1)$ & $\alpha\ket{001}-\beta\ket{101}$ & $(\alpha\ket{0}-\beta\ket{1})\ket{01}$ \\
$(1,0)$ & $\beta\ket{010}+\alpha\ket{110}$ & $(\beta\ket{0}+\alpha\ket{1})\ket{10}$ \\
$(1,1)$ & $-\beta\ket{011}+\alpha\ket{111}$ & $(-\beta\ket{0}+\alpha\ket{1})\ket{11}$ \\
\bottomrule
\end{tabular}
\end{center}
d'où
\[
  \ket{\Pi_2}=\frac12\Big[\big(\alpha\ket{0}+\beta\ket{1}\big)\ket{00}
    +\big(\alpha\ket{0}-\beta\ket{1}\big)\ket{01}
    +\big(\beta\ket{0}+\alpha\ket{1}\big)\ket{10}
    +\big(\alpha\ket{1}-\beta\ket{0}\big)\ket{11}\Big].
\]
Dans chaque terme, le ket de gauche est l'état de $B$ et le ket de droite $\ket{a,q}$ celui du
couple $(A,Q)$ d'Alice.

\subsubsection{Mesure d'Alice et corrections de Bob}

Alice mesure $A$ et $Q$. Comme à la section~\ref{sec:mesure2}, la probabilité d'un résultat
$(a,q)$ est le carré de la norme du terme correspondant, et l'état de Bob devient le ket
associé (déjà normé). Ici
\[
  P(a,q)=\Big\|\tfrac12\ket{\phi_{aq}}\Big\|^2=\tfrac14\big(|\alpha|^2+|\beta|^2\big)=\tfrac14
  \quad\text{pour les quatre résultats,}
\]
qui sont donc équiprobables et \emph{indépendants} de $\alpha$ et $\beta$ : les deux bits
envoyés ne révèlent rien sur $\ket{\varphi}$. Les kets de Bob sont
$\ket{\phi_{00}}=\alpha\ket{0}+\beta\ket{1}$, $\ket{\phi_{01}}=\alpha\ket{0}-\beta\ket{1}$,
$\ket{\phi_{10}}=\beta\ket{0}+\alpha\ket{1}$ et $\ket{\phi_{11}}=-\beta\ket{0}+\alpha\ket{1}$.

\begin{center}
\renewcommand{\arraystretch}{1.5}
\begin{tabular}{@{}cccll@{}}
\toprule
Résultat $(a,q)$ & Probabilité & État de Bob & Lien avec $\ket{\varphi}$ & Correction de Bob \\
\midrule
$(0,0)$ & $\frac14$ & $\alpha\ket{0}+\beta\ket{1}$ & $\ket{\varphi}$ & $I$ (rien) \\
$(0,1)$ & $\frac14$ & $\alpha\ket{0}-\beta\ket{1}$ & $Z\ket{\varphi}$ & $Z$ \\
$(1,0)$ & $\frac14$ & $\beta\ket{0}+\alpha\ket{1}$ & $X\ket{\varphi}$ & $X$ \\
$(1,1)$ & $\frac14$ & $-\beta\ket{0}+\alpha\ket{1}$ & $XZ\ket{\varphi}$ & $X$ puis $Z$ (opérateur $ZX$) \\
\bottomrule
\end{tabular}
\end{center}

\paragraph{Vérification des quatre corrections.} Avec $X\ket{0}=\ket{1}$, $X\ket{1}=\ket{0}$,
$Z\ket{0}=\ket{0}$ et $Z\ket{1}=-\ket{1}$ :
\begin{align*}
  (0,0):\quad & I\big(\alpha\ket{0}+\beta\ket{1}\big)=\alpha\ket{0}+\beta\ket{1},\\
  (0,1):\quad & Z\big(\alpha\ket{0}-\beta\ket{1}\big)=\alpha\ket{0}-\beta Z\ket{1}=\alpha\ket{0}+\beta\ket{1},\\
  (1,0):\quad & X\big(\beta\ket{0}+\alpha\ket{1}\big)=\beta\ket{1}+\alpha\ket{0}=\alpha\ket{0}+\beta\ket{1},\\
  (1,1):\quad & X\big(-\beta\ket{0}+\alpha\ket{1}\big)=-\beta\ket{1}+\alpha\ket{0}=\alpha\ket{0}-\beta\ket{1},
     \ \text{puis}\ Z\big(\alpha\ket{0}-\beta\ket{1}\big)=\alpha\ket{0}+\beta\ket{1}.
\end{align*}
Dans le dernier cas, avec les matrices,
$ZX=\begin{pmatrix}1&0\\0&-1\end{pmatrix}\begin{pmatrix}0&1\\1&0\end{pmatrix}
=\begin{pmatrix}0&1\\-1&0\end{pmatrix}$ et
$ZX\begin{pmatrix}-\beta\\ \alpha\end{pmatrix}=\begin{pmatrix}\alpha\\ \beta\end{pmatrix}$.
L'ordre compte : l'état de Bob est $XZ\ket{\varphi}$ ($Z\ket{\varphi}=\alpha\ket{0}-\beta\ket{1}$,
puis $X$ donne $\alpha\ket{1}-\beta\ket{0}$), et on l'annule avec
$(XZ)^{-1}=Z^{-1}X^{-1}=ZX$ : on applique d'abord $X$, puis $Z$, comme dans le circuit.
Dans les quatre cas, Bob retrouve exactement $\ket{\varphi}=\alpha\ket{0}+\beta\ket{1}$.

\subsubsection{Ce que Bob sait avant de recevoir les bits}

Tant que Bob ignore $(a,q)$, son qubit est dans le mélange des quatre kets $\ket{\phi_{aq}}$, chacun
avec la probabilité $\frac14$ :
\begin{align*}
  \rho_B&=\frac14\sum_{a,q}\ket{\phi_{aq}}\bra{\phi_{aq}}\\
  &=\frac14\Bigg[\begin{pmatrix}|\alpha|^2&\alpha\beta^*\\ \alpha^*\beta&|\beta|^2\end{pmatrix}
    +\begin{pmatrix}|\alpha|^2&-\alpha\beta^*\\ -\alpha^*\beta&|\beta|^2\end{pmatrix}
    +\begin{pmatrix}|\beta|^2&\alpha^*\beta\\ \alpha\beta^*&|\alpha|^2\end{pmatrix}
    +\begin{pmatrix}|\beta|^2&-\alpha^*\beta\\ -\alpha\beta^*&|\alpha|^2\end{pmatrix}\Bigg]\\
  &=\frac14\begin{pmatrix}2(|\alpha|^2+|\beta|^2)&0\\0&2(|\alpha|^2+|\beta|^2)\end{pmatrix}
   =\frac12\begin{pmatrix}1&0\\0&1\end{pmatrix}=\frac12\,\mathrm{Id}.
\end{align*}
Sans les deux bits, Bob a donc un qubit complètement mélangé, qui ne dépend pas de $\alpha$ ni de
$\beta$ : il n'a reçu aucune information sur $\ket{\varphi}$. C'est l'envoi classique de $(a,q)$
qui permet de la retrouver, donc rien ne voyage plus vite que la lumière.

\begin{remarque}[Bilan de la téléportation]
\begin{itemize}[nosep]
  \item Ressources : $1$ e-bit (la paire de Bell) et $2$ bits classiques pour transmettre
        \emph{un qubit} d'état quelconque, sans qu'Alice connaisse $\alpha$ ni $\beta$.
  \item Pas de clonage : après la mesure, $A$ et $Q$ sont dans l'état de base $\ket{a,q}$.
        Alice n'a plus $\ket{\varphi}$ : l'état a été \emph{déplacé}, pas copié. L'intrication
        entre $A$ et $B$ a été consommée.
  \item Les quatre corrections $I,Z,X,ZX$ sont des portes contrôlées par des bits classiques
        (section~\ref{sec:controlees}).
\end{itemize}
\end{remarque}

% ---------------------------------------------------------------------
\subsection{Portes contrôlées : contrôlé-\texorpdfstring{$U$}{U}, Fredkin et Toffoli}\label{sec:controlees}
% ---------------------------------------------------------------------

Le CNOT est le cas le plus simple d'une idée générale : appliquer une opération à un ou
plusieurs qubits \emph{cibles} seulement si un ou plusieurs qubits de \emph{contrôle} valent $1$.
Dans cette section, comme à la section~\ref{sec:op2}, le contrôle est le \emph{premier} chiffre
du ket. Les qubits sont notés $x$, $y$, $z$ (le qubit $x$ est le contrôle) et
$\ket{x,y,z}=\ket{x}\otimes\ket{y}\otimes\ket{z}$ ; d'après la convention de la section, le
contrôle est dessiné en bas (figure~\ref{fig:controlees}). Pour éviter la confusion avec les
matrices de Pauli, on écrit ici $\sigma_x=X$ pour la porte NOT.

\begin{figure}[!ht]
\centering
\begin{adjustbox}{max width=\linewidth}
\begin{tikzpicture}[>=Stealth, font=\small]
  % (a) contrôlé-U : contrôle x en bas, cible y en haut
  \draw (0,1.2) -- (3.4,1.2);
  \draw (0,0) -- (3.4,0);
  \node[anchor=east] at (0,1.2) {$y$};
  \node[anchor=east] at (0,0) {$x$};
  \draw[insarouge, thick] (1.7,0) -- (1.7,0.85);
  \pic at (1.7,0) {ctl};
  \node[porte] at (1.7,1.2) {$U$};
  \node[anchor=north, font=\footnotesize] at (1.7,-0.5) {(a) contrôlé-$U$};
  % (b) contrôlé-SWAP (Fredkin) : x contrôle, y et z échangés
  \draw (5,2.4) -- (8.4,2.4);
  \draw (5,1.2) -- (8.4,1.2);
  \draw (5,0) -- (8.4,0);
  \node[anchor=east] at (5,2.4) {$z$};
  \node[anchor=east] at (5,1.2) {$y$};
  \node[anchor=east] at (5,0) {$x$};
  \draw[insarouge, thick] (6.7,0) -- (6.7,2.4);
  \pic at (6.7,0) {ctl};
  \pic at (6.7,1.2) {croix};
  \pic at (6.7,2.4) {croix};
  \node[anchor=north, font=\footnotesize] at (6.7,-0.5) {(b) Fredkin (C-SWAP)};
  % (c) Toffoli (CCNOT) : x et y contrôlent, z cible
  \draw (10,2.4) -- (13.4,2.4);
  \draw (10,1.2) -- (13.4,1.2);
  \draw (10,0) -- (13.4,0);
  \node[anchor=east] at (10,2.4) {$z$};
  \node[anchor=east] at (10,1.2) {$y$};
  \node[anchor=east] at (10,0) {$x$};
  \draw[insarouge, thick] (11.7,0) -- (11.7,2.4);
  \pic at (11.7,0) {ctl};
  \pic at (11.7,1.2) {ctl};
  \pic at (11.7,2.4) {cible};
  \node[anchor=north, font=\footnotesize] at (11.7,-0.5) {(c) Toffoli (CCNOT)};
\end{tikzpicture}
\end{adjustbox}
\caption{Portes contrôlées. Le contrôle $x$ est le premier chiffre du ket, donc le fil du bas.
(a) Le qubit $y$ subit $U$ si $x=1$ ; pour $U=X$ on retrouve le CNOT. (b) Si $x=1$, les qubits
$y$ et $z$ sont échangés. (c) Si $x=y=1$, le qubit $z$ est inversé.}
\label{fig:controlees}
\end{figure}

\subsubsection{Contrôlé-\texorpdfstring{$U$}{U}}

\begin{definition}{Porte contrôlée-$U$}
Soit $U$ une matrice unitaire $2\times2$. Sur le système $(x,y)$, de contrôle $x$ et de cible $y$,
\[
  \mathrm{C}U=\ket{0}\bra{0}_x\otimes I_y+\ket{1}\bra{1}_x\otimes U_y .
\]
Le premier terme agit quand le contrôle est $\ket{0}$ : la cible ne change pas ($I$). Le second
agit quand le contrôle est $\ket{1}$ : la cible subit $U$.
\end{definition}

Comme $\ket{0}\bra{0}\ket{x}=\delta_{x0}\ket{0}$ et $\ket{1}\bra{1}\ket{x}=\delta_{x1}\ket{1}$,
\[
  \mathrm{C}U\big(\ket{x}\otimes\ket{y}\big)=\ket{x}\otimes U^x\ket{y},\qquad U^0=I,\ U^1=U .
\]
Par la règle $A\otimes B=(A_{ij}B)$, $\ket{0}\bra{0}\otimes I=\begin{pmatrix}I&0\\0&0\end{pmatrix}$
et $\ket{1}\bra{1}\otimes U=\begin{pmatrix}0&0\\0&U\end{pmatrix}$ (blocs $2\times2$), donc
\[
  \mathrm{C}U=\begin{pmatrix}I&0\\0&U\end{pmatrix}
  =\begin{pmatrix}1&0&0&0\\0&1&0&0\\0&0&u_{00}&u_{01}\\0&0&u_{10}&u_{11}\end{pmatrix}
  \qquad\text{pour }U=\begin{pmatrix}u_{00}&u_{01}\\u_{10}&u_{11}\end{pmatrix}.
\]
Cette matrice est unitaire car $I^\dagger I=I$ et $U^\dagger U=I$.

\paragraph{Le CNOT est un cas particulier.} Pour $U=\sigma_x$,
$\mathrm{C}\sigma_x=\ket{0}\bra{0}_x\otimes I_y+\ket{1}\bra{1}_x\otimes\sigma_x$ est le CNOT de
contrôle $x$ (premier chiffre), c'est-à-dire $\mathrm{CNOT}_{12}$ de la section~\ref{sec:cnot}.
On calcule l'action sur les quatre états de base, avec $\braket{0}{0}=\braket{1}{1}=1$,
$\braket{0}{1}=\braket{1}{0}=0$, $\sigma_x\ket{0}=\ket{1}$ et $\sigma_x\ket{1}=\ket{0}$ :
\begin{align*}
  \mathrm{C}\sigma_x\ket{00}&=\big[\ket{0}\bra{0}_x\otimes I_y\big]\ket{0}_x\ket{0}_y
      +\big[\ket{1}\bra{1}_x\otimes\sigma_x\big]\ket{0}_x\ket{0}_y\\
    &=\ket{0}\braket{0}{0}\otimes I\ket{0}+\ket{1}\braket{1}{0}\otimes\sigma_x\ket{0}
     =\ket{0}\otimes\ket{0}+0=\ket{00},\\[1.2ex]
  \mathrm{C}\sigma_x\ket{01}&=\ket{0}\braket{0}{0}\otimes I\ket{1}+\ket{1}\braket{1}{0}\otimes\sigma_x\ket{1}
     =\ket{0}\otimes\ket{1}+0=\ket{01},\\[1.2ex]
  \mathrm{C}\sigma_x\ket{10}&=\ket{0}\braket{0}{1}\otimes I\ket{0}+\ket{1}\braket{1}{1}\otimes\sigma_x\ket{0}
     =0+\ket{1}\otimes\ket{1}=\ket{11},\\[1.2ex]
  \mathrm{C}\sigma_x\ket{11}&=\ket{0}\braket{0}{1}\otimes I\ket{1}+\ket{1}\braket{1}{1}\otimes\sigma_x\ket{1}
     =0+\ket{1}\otimes\ket{0}=\ket{10}.
\end{align*}
On retrouve la table du CNOT : la cible n'est inversée que si le contrôle vaut $1$.

\paragraph{Un autre exemple : $U=Z$.} Ici
\[
  \mathrm{C}Z=\begin{pmatrix}I&0\\0&Z\end{pmatrix}=\mathrm{diag}(1,1,1,-1) :
\]
la porte multiplie $\ket{11}$ par $-1$ et ne touche pas les autres états de base. Elle est
symétrique entre les deux qubits. Elle se déduit du CNOT par
$\mathrm{C}Z=(I\otimes H)\,\mathrm{C}X\,(I\otimes H)$, où $\mathrm{C}X$ est le contrôlé-$X$.
En effet, $I\otimes H=\mathrm{diag}(H,H)$ (deux blocs $2\times2$ égaux à $H$), et le produit de
matrices diagonales par blocs se fait bloc par bloc :
\[
  (I\otimes H)\begin{pmatrix}I&0\\0&X\end{pmatrix}(I\otimes H)
  =\begin{pmatrix}H\,I\,H&0\\0&H\,X\,H\end{pmatrix}
  =\begin{pmatrix}I&0\\0&HXH\end{pmatrix},
\]
car $H^2=I$, et
\[
  HXH=\frac12\begin{pmatrix}1&1\\1&-1\end{pmatrix}\begin{pmatrix}0&1\\1&0\end{pmatrix}\begin{pmatrix}1&1\\1&-1\end{pmatrix}
  =\frac12\begin{pmatrix}1&1\\1&-1\end{pmatrix}\begin{pmatrix}1&-1\\1&1\end{pmatrix}
  =\frac12\begin{pmatrix}2&0\\0&-2\end{pmatrix}=Z
\]
(section~\ref{sec:hadamard} : $H$ échange les axes $x$ et $z$).

\subsubsection{Contrôlé-SWAP : la porte de Fredkin}

\begin{definition}{Porte de Fredkin (C-SWAP)}
Sur trois qubits $(x,y,z)$ (espace de dimension $2^3=8$), avec $x$ pour contrôle,
\[
  \mathrm{C\text{-}SWAP}=\ket{0}\bra{0}_x\otimes I_y\otimes I_z+\ket{1}\bra{1}_x\otimes\mathrm{SWAP}_{yz}.
\]
Si $x=0$, rien ne change ; si $x=1$, les qubits $y$ et $z$ sont échangés.
\end{definition}

Calcul sur quatre états de base (avec $I\ket{ab}=\ket{ab}$ et $\mathrm{SWAP}\ket{ab}=\ket{ba}$) :
\begin{align*}
  \mathrm{C\text{-}SWAP}\ket{000}&=\ket{0}\braket{0}{0}\otimes I\ket{00}+\ket{1}\braket{1}{0}\otimes\mathrm{SWAP}\ket{00}
     =\ket{0}\otimes\ket{00}+0=\ket{000},\\[1.2ex]
  \mathrm{C\text{-}SWAP}\ket{101}&=\ket{0}\braket{0}{1}\otimes I\ket{01}+\ket{1}\braket{1}{1}\otimes\mathrm{SWAP}\ket{01}
     =0+\ket{1}\otimes\ket{10}=\ket{110},\\[1.2ex]
  \mathrm{C\text{-}SWAP}\ket{110}&=\ket{0}\braket{0}{1}\otimes I\ket{10}+\ket{1}\braket{1}{1}\otimes\mathrm{SWAP}\ket{10}
     =0+\ket{1}\otimes\ket{01}=\ket{101},\\[1.2ex]
  \mathrm{C\text{-}SWAP}\ket{010}&=\ket{0}\braket{0}{0}\otimes I\ket{10}+\ket{1}\braket{1}{0}\otimes\mathrm{SWAP}\ket{10}
     =\ket{0}\otimes\ket{10}+0=\ket{010}.
\end{align*}
Les huit états de base donnent :
\begin{center}
\renewcommand{\arraystretch}{1.3}
\begin{tabular}{@{}ccccccccc@{}}
\toprule
Entrée & $\ket{000}$ & $\ket{001}$ & $\ket{010}$ & $\ket{011}$ & $\ket{100}$ & $\ket{101}$ & $\ket{110}$ & $\ket{111}$ \\
\midrule
Sortie & $\ket{000}$ & $\ket{001}$ & $\ket{010}$ & $\ket{011}$ & $\ket{100}$ & $\ket{110}$ & $\ket{101}$ & $\ket{111}$ \\
\bottomrule
\end{tabular}
\end{center}
Seuls $\ket{101}$ et $\ket{110}$ changent (contrôle à $1$ et $y\neq z$) ; pour $\ket{100}$ et
$\ket{111}$ l'échange de deux qubits égaux ne change rien. Dans la base
$\ket{000},\ket{001},\dots,\ket{111}$ (ordre binaire), la matrice est
$\mathrm{diag}(I_4,\mathrm{SWAP})$ (deux blocs $4\times4$) :
\[
  \mathrm{C\text{-}SWAP}=\begingroup\setlength{\arraycolsep}{3pt}\begin{pmatrix}
    1&0&0&0&0&0&0&0\\
    0&1&0&0&0&0&0&0\\
    0&0&1&0&0&0&0&0\\
    0&0&0&1&0&0&0&0\\
    0&0&0&0&1&0&0&0\\
    0&0&0&0&0&0&1&0\\
    0&0&0&0&0&1&0&0\\
    0&0&0&0&0&0&0&1
  \end{pmatrix}\endgroup .
\]
C'est l'identité dont les lignes $6$ et $7$ sont échangées.

\paragraph{Contrôle en superposition.} Avec $\ket{+}_x\otimes\ket{0}_y\otimes\ket{1}_z
=\frac{1}{\sqrt2}\big(\ket{001}+\ket{101}\big)$ :
\[
  \mathrm{C\text{-}SWAP}\,\frac{1}{\sqrt2}\big(\ket{001}+\ket{101}\big)
  =\frac{1}{\sqrt2}\big(\ket{001}+\ket{110}\big).
\]
Le résultat n'est pas un produit : les coefficients de $\ket{x}\otimes\ket{yz}$ forment la
matrice (lignes $x=0,1$, colonnes $yz=00,01,10,11$)
\[
  \frac{1}{\sqrt2}\begin{pmatrix}0&1&0&0\\0&0&1&0\end{pmatrix},
\]
dont les deux lignes ne sont pas proportionnelles. Le qubit $x$ est intriqué avec la paire
$(y,z)$.

\subsubsection{Contrôlé-contrôlé-NOT : la porte de Toffoli}

\begin{definition}{Porte de Toffoli (CCNOT)}
Sur trois qubits $(x,y,z)$, avec $x$ et $y$ pour contrôles et $z$ pour cible,
\[
  \mathrm{CCNOT}=\ket{0}\bra{0}_x\otimes I_y\otimes I_z
   +\ket{1}\bra{1}_x\otimes\Big[\ket{0}\bra{0}_y\otimes I_z+\ket{1}\bra{1}_y\otimes\sigma_x\Big].
\]
Le crochet est un CNOT (contrôle $y$, cible $z$), appliqué seulement si $x=1$ : si $x=0$, rien ;
si $x=1$ et $y=0$, rien ; si $x=1$ et $y=1$, le qubit $z$ est inversé.
\end{definition}

Calcul de $\mathrm{CCNOT}\ket{110}$ :
\begin{align*}
  \mathrm{CCNOT}\ket{110}&=\ket{0}\braket{0}{1}\otimes I\ket{1}\otimes I\ket{0}
     +\ket{1}\braket{1}{1}\otimes\Big[\ket{0}\braket{0}{1}\otimes I\ket{0}
       +\ket{1}\braket{1}{1}\otimes\sigma_x\ket{0}\Big]\\
  &=0+\ket{1}\otimes\Big[0+\ket{1}\otimes\ket{1}\Big]=\ket{111}.
\end{align*}
Pour les huit états de base, $\mathrm{CCNOT}\ket{x,y,z}=\ket{x,\,y,\,z\oplus xy}$ :
\begin{center}
\renewcommand{\arraystretch}{1.3}
\begin{tabular}{@{}ccccccccc@{}}
\toprule
Entrée & $\ket{000}$ & $\ket{001}$ & $\ket{010}$ & $\ket{011}$ & $\ket{100}$ & $\ket{101}$ & $\ket{110}$ & $\ket{111}$ \\
\midrule
Sortie & $\ket{000}$ & $\ket{001}$ & $\ket{010}$ & $\ket{011}$ & $\ket{100}$ & $\ket{101}$ & $\ket{111}$ & $\ket{110}$ \\
\bottomrule
\end{tabular}
\end{center}
Seuls $\ket{110}$ et $\ket{111}$ sont échangés : la matrice est $\mathrm{diag}(I_6,X)$
(un bloc $6\times6$ égal à l'identité, puis $X$),
\[
  \mathrm{CCNOT}=\begingroup\setlength{\arraycolsep}{3pt}\begin{pmatrix}
    1&0&0&0&0&0&0&0\\
    0&1&0&0&0&0&0&0\\
    0&0&1&0&0&0&0&0\\
    0&0&0&1&0&0&0&0\\
    0&0&0&0&1&0&0&0\\
    0&0&0&0&0&1&0&0\\
    0&0&0&0&0&0&0&1\\
    0&0&0&0&0&0&1&0
  \end{pmatrix}\endgroup .
\]

\begin{remarque}[Portes réversibles]
Les portes de Fredkin et de Toffoli sont des permutations des états de base qui échangent
seulement deux d'entre eux : elles sont leurs propres inverses ($\mathrm{C\text{-}SWAP}^2=\mathrm{CCNOT}^2=\mathrm{Id}$)
et unitaires. Avec $z=0$, la Toffoli calcule le ET logique $x\wedge y$ sans rien effacer :
$\ket{x,y,0}\to\ket{x,y,xy}$ ; avec $z=1$, elle calcule le NON-ET. Elle sert à écrire tout calcul
classique de façon réversible.
\end{remarque}

% ---------------------------------------------------------------------
\subsection{Récapitulatif du chapitre}
% ---------------------------------------------------------------------

\begin{center}
\renewcommand{\arraystretch}{1.5}
\begin{tabular}{@{}lp{12cm}@{}}
\toprule
\textbf{Notion} & \textbf{Formule ou résultat} \\
\midrule
CNOT & $\ket{b,a}\to\ket{b\oplus a,a}$ (contrôle $a$ écrit en dernier), $\mathrm{CNOT}^2=\mathrm{Id}$ \\
SWAP & $\mathrm{SWAP}\ket{a,b}=\ket{b,a}$, \quad $\mathrm{CNOT}_{21}=\mathrm{SWAP}\,\mathrm{CNOT}_{12}\,\mathrm{SWAP}$ \\
Circuit de Bell & $(I\otimes H)$ puis CNOT : $\ket{\beta^{xy}}=\frac{1}{\sqrt2}\big(\ket{y,0}+(-1)^x\ket{\bar y,1}\big)$, $\ket{\beta^{00}}=\ket{\Phi^+}$ \\
Information & $I(A;B)=H(A)+H(B)-H(A,B)$, \quad $H(A|B)=H(A,B)-H(B)$ \\
Trace partielle & $\operatorname{Tr}_B\big(\ket{a}\bra{a'}\otimes\ket{b}\bra{b'}\big)=\ket{a}\bra{a'}\braket{b'}{b}$ ; état de Bell : $\rho_A=\frac12\mathrm{Id}$ \\
État de Bell (bits) & $H(A,B)=0$, \quad $H(A)=H(B)=1$, \quad $I(A;B)=2$, \quad $H(A|B)=-1$ \\
Pur / mélangé & $\operatorname{Tr}\rho^2=1$ / $<1$ ; $\ket{+}$ : $1$ ; qubit d'une paire de Bell : $\frac12$ ; source de Bob : $\frac58$ \\
Téléportation & $1$ e-bit $+\ 2$ bits classiques ; correction : $X$ si $a=1$, puis $Z$ si $q=1$ ; $P(a,q)=\frac14$ \\
Contrôlé-$U$ & $\ket{0}\bra{0}\otimes I+\ket{1}\bra{1}\otimes U=\mathrm{diag}(I,U)$ ; $\ket{x,y}\to\ket{x}\otimes U^x\ket{y}$ \\
Fredkin & $\ket{0}\bra{0}\otimes I\otimes I+\ket{1}\bra{1}\otimes\mathrm{SWAP}$ ; $\ket{101}\leftrightarrow\ket{110}$ \\
Toffoli & $\ket{x,y,z}\to\ket{x,y,z\oplus xy}$ ; $\ket{110}\leftrightarrow\ket{111}$ \\
\bottomrule
\end{tabular}
\end{center}
